\documentclass[11pt,a4paper]{article}
\usepackage[utf8]{inputenc}\usepackage[T1]{fontenc}
\usepackage[spanish,es-noquoting]{babel}
\usepackage[margin=2.3cm]{geometry}
\usepackage{longtable,booktabs,array,xcolor,fancyhdr,titlesec,enumitem,hyperref}
\usepackage[most]{tcolorbox}
\definecolor{resalte}{RGB}{255,243,166}
\definecolor{gris}{RGB}{95,95,95}
\newtcolorbox{pasaje}{colback=resalte,colframe=resalte!62!black,boxrule=0.3pt,left=5pt,right=5pt,top=4pt,bottom=4pt,before skip=4pt,after skip=7pt,breakable}
\hypersetup{colorlinks=true,linkcolor=black,urlcolor=black,pdftitle={Respaldo documental de las citas}}
\pagestyle{fancy}\fancyhf{}\fancyhead[L]{\small Respaldo documental de las citas}\fancyhead[R]{\small\thepage}\renewcommand{\headrulewidth}{0.4pt}
\titleformat{\section}{\large\bfseries}{\thesection.}{0.6em}{}
\titleformat{\subsection}{\normalsize\bfseries}{}{0em}{}
\setlist{nosep,leftmargin=1.4em}
\newcommand{\campo}[1]{\textcolor{gris}{\footnotesize #1}}
\begin{document}
\begin{center}
{\LARGE\bfseries Respaldo documental de las citas}\\[3pt]
{\large EduFEM --- tesis de grado}\\[8pt]
\begin{minipage}{0.86\textwidth}\small
Este documento acompaña a la tesis y no forma parte de ella. Registra, para cada fuente citada, el ejemplar exactamente consultado, la página donde se encuentra lo que la tesis le atribuye y el pasaje textual que lo sostiene. Las citas del cuerpo de la tesis llevan localizador de página cuando sostienen una ecuación, un valor numérico, un umbral, una definición o una atribución concreta, y van sin localizador cuando son de encuadre (regla declarada en la Introducción de la tesis); este documento registra, para todas, la página impresa y el pasaje.\\[4pt]
Todas las páginas indicadas son \textbf{páginas impresas del ejemplar}, leídas del encabezado o del pie de la propia página. No se calcularon por desplazamiento: en varios ejemplares el desfase entre la página del visor y la impresa no es constante.
\end{minipage}\end{center}
\vspace{4pt}\hrule\vspace{10pt}
\section{Estado de las fuentes}
\begin{longtable}{@{}p{4.6cm}p{2.1cm}p{1.5cm}p{1.3cm}p{5.4cm}@{}}
\toprule \textbf{Fuente} & \textbf{Cotejo} & \textbf{Riesgo} & \textbf{Citas} & \textbf{Observación} \\ \midrule \endhead
Teaching the finite element method fundamentals to... & COINCIDE & baja & 10 & Coincide con lo declarado. \\ \addlinespace[2pt]
Metodologia de la Investigacion Cientifica & COINCIDE & ninguna & 5 & Coincide con lo declarado. \\ \addlinespace[2pt]
Learning from Examples: Instructional Principles from the... & COINCIDE & ninguna & 2 & Coincide con lo declarado. \\ \addlinespace[2pt]
Finite Element Procedures & COINCIDE & ninguna & 10 & Coincide con lo declarado. \\ \addlinespace[2pt]
The ICAP Framework: Linking Cognitive Engagement to... & COINCIDE & ninguna & 2 & Coincide con lo declarado. \\ \addlinespace[2pt]
Concepts and Applications of Finite Element Analysis & COINCIDE & ninguna & 15 & Coincide con lo declarado. \\ \addlinespace[2pt]
CSI Analysis Reference Manual --- For SAP2000, ETABS, SAFE... & COINCIDE & ninguna & 4 & Coincide con lo declarado. \\ \addlinespace[2pt]
La investigacion tecnologica: Investigar, idear e innovar... & COINCIDE & ninguna & 3 & Coincide con lo declarado. \\ \addlinespace[2pt]
Array programming with NumPy & COINCIDE & ninguna & 4 & Coincide con lo declarado. \\ \addlinespace[2pt]
The Finite Element Method: Linear Static and Dynamic... & COINCIDE & ninguna & 10 & Coincide con lo declarado. \\ \addlinespace[2pt]
Interactive Simulation of Eigenmodes for Instruction of... & COINCIDE & ninguna & 9 & Coincide con lo declarado. \\ \addlinespace[2pt]
Interactive Simulation of Finite Element Equation... & COINCIDE & ninguna & 13 & Coincide con lo declarado. \\ \addlinespace[2pt]
Verification and Validation in Scientific Computing & COINCIDE & ninguna & 9 & Coincide con lo declarado. \\ \addlinespace[2pt]
Structural Analysis with the Finite Element Method.... & COINCIDE & ninguna & 8 & Coincide con lo declarado. \\ \addlinespace[2pt]
FEM education in undergraduate studies: Industry-informed... & COINCIDE & ninguna & 9 & Coincide con lo declarado. \\ \addlinespace[2pt]
An Introduction to the Finite Element Method & COINCIDE & ninguna & 6 & Coincide con lo declarado. \\ \addlinespace[2pt]
The learning effects of computer simulations in science... & COINCIDE & ninguna & 2 & Coincide con lo declarado. \\ \addlinespace[2pt]
Code Verification by the Method of Manufactured Solutions & COINCIDE & ninguna & 6 & Coincide con lo declarado. \\ \addlinespace[2pt]
Metodología de la investigación: las rutas cuantitativa,... & COINCIDE & ninguna & 3 & Coincide con lo declarado. \\ \addlinespace[2pt]
The Verdict Geometric Quality Library & COINCIDE & ninguna & 7 & Coincide con lo declarado. \\ \addlinespace[2pt]
An Analysis of the Finite Element Method & COINCIDE & ninguna & 3 & Coincide con lo declarado. \\ \addlinespace[2pt]
ED-Elas2D: An Educational Program for Computer-aided... & COINCIDE & ninguna & 11 & Coincide con lo declarado. \\ \addlinespace[2pt]
Theory of Elasticity & COINCIDE & ninguna & 5 & Coincide con lo declarado. \\ \addlinespace[2pt]
SciPy 1.0: fundamental algorithms for scientific... & COINCIDE & ninguna & 6 & Coincide con lo declarado. \\ \addlinespace[2pt]
The Finite Element Method: Its Basis and Fundamentals & COINCIDE & ninguna & 11 & Coincide con lo declarado. \\ \addlinespace[2pt]
\bottomrule\end{longtable}
\section{Afirmaciones sin respaldo o con respaldo parcial}
\noindent\small Lo que sigue es el resultado más accionable de la revisión: afirmaciones de la tesis cuya fuente citada no las sostiene, o las sostiene solo en parte.\par\vspace{6pt}
\subsection{\texttt{bishay2020teaching}}
\begin{itemize}\small
\item \textbf{[PARCIAL]} NOTA PREVIA SOBRE PAGINACION (afecta a todos los renglones): el .bib declara pages = \{1007--1027\} (referencias.bib:148), pero el ejemplar PDF disponible imprime la paginacion 1-21. Encabezados: 'BISHAY \textbar{}3', '4\textbar{}', ... '\textbar{} 21'; pie de la p.1: 'Comput Appl Eng Educ. 2020;1-21.  wileyonlinelibrary.com/journal/cae'. Todas las paginas...
\\ \campo{El ejemplar es la version online-first (21 paginas, DOI 10.1002/cae.22281). El volumen 28, numero 4 y el rango 1007-1027 no aparecen impresos en ninguna pagina; probablemente sean correctos (paginacion final del fasciculo) pero no verificables con este PDF. Convendria o bien conseguir el PDF del fasciculo, o bien dejar constancia de que el rango proviene del registro editorial y no del ejemplar consultado.}
\item \textbf{[PARCIAL]} 01\_introduccion.tex:21 - 'La literatura sobre ensenanza del metodo coincide en que la simulacion interactiva y la manipulacion directa de los objetos matematicos del MEF favorecen un aprendizaje mas profundo que la mera exposicion teorica o el uso de herramientas opacas'
\\ \campo{Bishay respalda la mitad de la afirmacion: que construir y manipular el codigo paso a paso es mas efectivo que la exposicion teorica, y que el software comercial opaco es un problema (p. 2). NO respalda la parte de 'simulacion interactiva': la palabra 'interactive' no aparece en el articulo y las herramientas son scripts de MATLAB sin interfaz interactiva. La afirmacion sobre interactividad se apoya en lee2015interactive, no en esta fuente.}
\item \textbf{[PARCIAL]} 03\_diseno\_implementacion.tex:9 - 'El software educativo interactivo demuestra eficacia para consolidar conceptos abstractos del analisis estructural'
\\ \campo{ATENCION: hay un dato incomodo que la tesis deberia respetar. Bishay si mide eficacia (mejora estadisticamente significativa, p = .048, dos cohortes de 27 alumnos), pero en la p. 17 atribuye explicitamente esa mejora 'principalmente' a la actividad de generacion de tareas entre pares, no a las herramientas computacionales, que solo la 'facilitaron'. Citar a Bishay como evidencia de que 'el software educativo interactivo demuestra eficacia' sobre-atribuye...}
\item \textbf{[PARCIAL]} 03\_diseno\_implementacion.tex:155 - 'El diseno de estos modulos parte de la evidencia sobre la interactividad como motor del aprendizaje en ingenieria \textbackslash{}autocite\{lee2015interactive,bishay2020teaching\}'
\\ \campo{Bishay documenta 'engagement' (implicacion, motivacion intrinseca) como motor del aprendizaje, pero ese engagement proviene de la AUTORIA del alumno sobre el problema y la tarea, no de la interactividad de una interfaz. Es un respaldo oblicuo. Si la tesis quiere sostener 'interactividad' con esta fuente, conviene reformular a 'la implicacion activa del alumno' o dejar la carga en lee2015interactive.}
\item \textbf{[PARCIAL]} 05\_conclusiones.tex:51 - recomendacion de 'un estudio con estudiantes [...] mediante una comparacion antes-despues e instrumentos de percepcion'; 'los antecedentes revisados \textbackslash{}autocite\{perezsantiago2023fem,bishay2020teaching,lee2015interactive\} indican, en cambio, que conceptos conviene poner a prueba'
\\ \campo{Bishay si aporta un modelo metodologico replicable (examen comun + encuestas de autopercepcion), lo que sostiene la parte de 'instrumentos de percepcion'. Pero su diseno NO es 'antes-despues' (pretest/postest) sino una comparacion entre dos cohortes independientes con indice de desempeno normalizado por la nota del prerrequisito. Si la tesis quiere apoyarse en este antecedente, deberia decir 'comparacion entre grupos' o citarlo como diseno alternativo.}
\end{itemize}
{\small\campo{Sin respaldo localizable:} TRES afirmaciones citadas a esta clave NO tienen respaldo localizable en el articulo:

(1) 02\_marco\_teorico.tex:29, tabla comparativa, columna \textquotedbl{}Licencia: Libre\textquotedbl{} para el constructor y analizador de reticulados. NO hay ninguna declaracion de licencia, disponibilidad publica, repositorio ni caracter libre/abierto en las 21 paginas. Busque explicitamente \textquotedbl{}available\textquotedbl{}, \textquotedbl{}download\textquotedbl{}, \textquotedbl{}open source\textquotedbl{}, \textquotedbl{}free of charge\textquotedbl{}, \textquotedbl{}supporting information\textquotedbl{} y \textquotedbl{}license\textquotedbl{}: el unico verbo de entrega es \textquotedbl{}another main code called 'Main\_Truss\_Builder' was also provided to them\textquotedbl{} (p. 6), es decir, entregado a los alumnos del curso. Ademas, las herramientas corren sobre MATLAB, plataforma comercial (p. 4: \textquotedbl{}writing a simple computer code on MATLAB\textquotedbl{}), lo que hace dudoso etiquetarlas como \textquotedbl{}Libre\textquotedbl{}. RECOMENDACION: cambiar esa celda a \textquotedbl{}No declarada\textquotedbl{} o \textquotedbl{}Academico (requiere MATLAB)\textquotedbl{}. Es el hallazgo mas grave de esta verificacion.

(2) 02\_marco\_teorico.tex:29, tabla comparativa, columna \textquotedbl{}2D/3D\textquotedbl{}: la tesis asigna \textquotedbl{}1D/2D\textquotedbl{} a Bishay. El articulo cubre explicita y extensamente reticulados 2D Y 3D (p. 5: \textquotedbl{}The last code for the full analysis of 2D trusses is extended to analyze 3D trusses where now each node has three translational degrees of freedom\textquotedbl{}; p. 11-12: figuras de torres de transmision, montana rusa, torre Eiffel, chasis de Formula, todas 3D; el titulo mismo del articulo habla de \textquotedbl{}2D and 3D trusses\textquotedbl{}). Bajo el encabezado \textquotedbl{}2D/3D\textquotedbl{} de esa columna, el valor correcto seria \textquotedbl{}2D/3D (barras)\textquotedbl{}. Si lo que se quiso indicar es la...\par}
\subsection{\texttt{alvarez\_metodologia}}
{\small\campo{Sin respaldo localizable:} Ninguna de las cuatro citas queda sin respaldo: las cuatro afirmaciones concretas atribuidas a esta fuente (contradiccion/insuficiencia como origen del problema; objeto de estudio; campo de accion como parte del objeto; modelo como representacion ideal en la que el investigador abstrae y sistematiza lo esencial) estan localizadas en paginas impresas concretas (7-8, 8, 13 y 21 respectivamente).

Dos observaciones que si conviene registrar:

(1) ROTULO NO RESPALDADO POR LA FUENTE. La expresion \textquotedbl{}modelo de simulacion numerica\textquotedbl{} (02b\_diseno\_metodologico.tex:14) no figura en el libro: la busqueda de la raiz \textquotedbl{}simulaci\textquotedbl{} en las 80 paginas del PDF devuelve cero ocurrencias, y Alvarez de Zayas no propone ninguna tipologia de modelos que la incluya. La cita respalda unicamente la clausula entre rayas (\textquotedbl{}representacion ideal del objeto...\textquotedbl{}), lo cual la redaccion actual permite leer correctamente, pero el rotulo es aporte del autor de la tesis, no de la fuente. Si algun miembro del tribunal lo lee como categoria tomada de Alvarez, seria una atribucion incorrecta.

(2) NINGUNA DE LAS CUATRO CITAS LLEVA LOCALIZADOR. Las cuatro son \textbackslash{}autocite\{alvarez\_metodologia\} a secas, sobre una monografia de 80 paginas. Para las categorias del marco metodologico (objeto, campo de accion, contradiccion, modelo) el tribunal UATF suele exigir pagina. Localizadores propuestos: p. 7-8 (contradiccion/insuficiencia), p. 8 (objeto de estudio), p. 13 (campo de accion), p. 21 (modelo), y opcionalmente p. 22 (variable...\par}
\subsection{\texttt{bathe2014fem}}
\begin{itemize}\small
\item \textbf{[PARCIAL]} «La condicion det J \textgreater{} 0 en todo el elemento garantiza que el mapeo sea biyectivo y el elemento valido [...] Un determinante nulo o negativo indica un elemento degenerado o invertido» --- tesis/capitulos/02\_marco\_teorico.tex:176 (\textbackslash{}autocite\{bathe2014fem,cook2002concepts\}, sin localizador)
\\ \campo{Bathe respalda el vinculo biyeccion \textless{}-\textgreater{} Jacobiano NO SINGULAR y el criterio geometrico (angulos interiores \textless{} 180 grados), y define det J. Lo que NO dice en ninguna de estas paginas es la condicion de SIGNO estricto det J \textgreater{} 0 ni el termino 'invertido': eso es formulacion de Cook. Si se conserva la doble cita esta bien; si se pusiera localizador, seria \textbackslash{}autocite[p.\textasciitilde{}346]\{bathe2014fem\}.}
\item \textbf{[PARCIAL]} «Existen tecnicas correctivas [...] entre ellas la integracion reducida selectiva (SRI), que subintegra solo los terminos responsables del bloqueo, y la formulacion B-bar, que proyecta la parte volumetrica o de corte de la deformacion» --- tesis/capitulos/02\_marco\_teorico.tex:257 y, con la misma redaccion,...
\\ \campo{La mitad SRI esta respaldada de forma casi literal (incluido «subintegra solo los terminos responsables del bloqueo» = «only the shear terms are not integrated exactly», p. 477). La mitad B-bar NO: 'B-bar' / 'B-barra' no aparece en el indice ni en el texto de esta fuente; el tratamiento equivalente de Bathe es la formulacion mixta desplazamiento/presion (u/p, u/p-c, secs. 4.4.3 y 4.5, pp. 276--300). Es decir, para B-bar el respaldo debe recaer solo en...}
\item \textbf{[PARCIAL]} «El elemento Q4 exhibe el fenomeno de bloqueo por cortante con claridad [...] comportamiento caracteristico de un elemento que resulta excesivamente rigido cuando debe representar flexion mediante una interpolacion bilineal incapaz de reproducir el modo de curvatura pura» --- tesis/capitulos/04\_resultados.tex:172...
\\ \campo{Bathe respalda el principio general (la formulacion en desplazamientos SOBREESTIMA la rigidez, p. 476) y el mecanismo de la contribucion espuria de corte, pero lo desarrolla para elementos ESTRUCTURALES (viga, placa, cascara): el indice remite 'Shear locking' a las pp. 275, 332, 404, 408, 424, 444, todas de vigas/placas/cascaras. NO encontre en esta fuente el enunciado del bloqueo por cortante del cuadrilatero bilineal Q4 de continuo plano ni la frase...}
\item \textbf{[PARCIAL]} «la cadena cinematica N -\textgreater{} J -\textgreater{} B precede a la constitutiva, y la deformacion eps = B u antecede a la tension sigma = D eps, en consonancia con la presentacion clasica de la literatura» --- tesis/capitulos/04\_resultados.tex:209 (\textbackslash{}autocite\{bathe2014fem,cook2002concepts\})
\\ \campo{Cita de encuadre, no de contenido: el orden expositivo interpolaciones h\_i (5.19) -\textgreater{} operador Jacobiano J (5.24) -\textgreater{} matriz B (5.26) -\textgreater{} ley constitutiva C en (5.27) es efectivamente el de Bathe, de modo que la afirmacion 'en consonancia con la presentacion clasica' es verificable en pp. 345--347. No hay una frase unica que la enuncie; se comprueba por la secuencia del capitulo 5.}
\end{itemize}
{\small\campo{Sin respaldo localizable:} (1) DISCREPANCIA DE EDICION que afecta a TODOS los localizadores: el .bib declara «edition = \{2\}, publisher = \{K. J. Bathe\}, location = \{Watertown, MA\}, year = \{2014\}, isbn = \{978-0-9790049-5-7\}», pero el ejemplar consultado es la edicion de Prentice Hall de 1996 (ISBN 0-13-301458-4 impreso en la contraportada, PDF pag. 1052; indice de contenidos en pp. v-xii, 12 capitulos, References en p. 1013, Index en p. 1029). Todas las paginas que doy arriba son de ESE ejemplar. Si se conserva la entrada .bib como ed. 2014, ninguna de estas paginas es citable tal cual: hay que corregir la entrada al ejemplar de 1996 o localizar de nuevo en la ed. 2014.

(2) Termino «hourglass» (02\_marco\_teorico.tex:244 y :255): NO figura en esta fuente. No aparece en el indice (seccion H, p. 1032) ni en el texto de las secs. 5.5.5-5.5.6. Bathe habla de «spurious zero energy modes» y «phantom modes» (p. 472). La equivalencia terminologica que la tesis presenta entre parentesis debe atribuirse a Cook, no a Bathe.

(3) Formulacion «B-bar» (02\_marco\_teorico.tex:257 y 05\_conclusiones.tex:45): NO encontre respaldo en esta fuente. No hay entrada 'B-bar' ni 'B-barra' en el indice (seccion B, p. 1029) y las secs. 4.4.3, 4.5 y 5.5.6 tratan el mismo problema con formulaciones mixtas desplazamiento/presion (u/p, u/p-c, elemento 9/3, pp. 276-300). Solo la parte 'SRI' de esa frase esta respaldada por Bathe (p. 476-477); la parte 'B-bar' depende exclusivamente de Hughes.

(4) Bloqueo por cortante del Q4 de CONTINUO...\par}
\subsection{\texttt{cook2002concepts}}
\begin{itemize}\small
\item \textbf{[PARCIAL]} tesis/capitulos/02\_marco\_teorico.tex:70 --- \textquotedbl{}La tensión plana modela cuerpos delgados cargados en su plano, en los que la tensión normal al plano es despreciable ($\sigma$z = 0)\textquotedbl{}.
\\ \campo{Cook define la condición de tensión plana por las componentes que se anulan ($\sigma$z = $\tau$yz = $\tau$zx = 0) y da la [E] correspondiente, pero no la caracteriza físicamente como \textquotedbl{}cuerpos delgados cargados en su plano\textquotedbl{} ni pone el ejemplo de la chapa o membrana. Esa lectura física corresponde al co-citado timoshenko1970elasticity.}
\item \textbf{[PARCIAL]} tesis/capitulos/02\_marco\_teorico.tex:176 --- \textquotedbl{}La condición det J \textgreater{} 0 en todo el elemento garantiza que el mapeo sea biyectivo y el elemento válido... Un determinante nulo o negativo indica un elemento degenerado o invertido; EduFEM impone un umbral mínimo para detectar y rechazar esas geometrías\textquotedbl{}.
\\ \campo{Cook sí respalda el criterio operativo (J negativo = elemento inaceptable, y el uso de umbrales de forma que llevan al software a rechazar la malla), pero NO enuncia det J \textgreater{} 0 en términos de biyectividad del mapeo. Advertencia: el epígrafe §6.12 \textquotedbl{}Validity of Isoparametric Elements\textquotedbl{} (p. 237) no trata este punto --- trata de la capacidad de representar estados de deformación constante ---, de modo que un lector que busque el respaldo por el título no lo...}
\item \textbf{[PARCIAL]} tesis/capitulos/02\_marco\_teorico.tex:251 y 408 --- \textquotedbl{}El elemento Q9, cuya interpolación bicuadrática sí reproduce el modo de curvatura, está prácticamente libre de este bloqueo\textquotedbl{} / \textquotedbl{}el elemento Q9, de orden superior, lo evita y converge con rapidez\textquotedbl{}.
\\ \campo{Cook muestra con datos que el elemento de nueve nodos con 2$\times$2 reproduce la flexión de una viga en voladizo con error prácticamente nulo (vC = 1.006 frente al exacto 1.000), y afirma en p. 236 que el elemento es capaz de modelar flexión pura exactamente si es rectangular. Pero no formula la frase \textquotedbl{}el Q9 está prácticamente libre de bloqueo por cortante\textquotedbl{}: esa lectura la aporta el co-citado hughes2000fem. Para el Q4 el contraste sí es explícito en p. 235,...}
\item \textbf{[PARCIAL]} tesis/capitulos/02\_marco\_teorico.tex:315 y tesis/capitulos/04\_resultados.tex:58 --- \textquotedbl{}En el elemento Q4 los puntos de la regla 2$\times$2 coinciden con los puntos de Barlow, ubicaciones donde las tensiones convergen con un orden superior al del resto del elemento (superconvergencia)\textquotedbl{}.
\\ \campo{Tres precisiones honestas. (1) El término \textquotedbl{}puntos de Barlow\textquotedbl{} NO aparece en esta obra: no figura en el índice de materias (p. 711, entre \textquotedbl{}Bar element\textquotedbl{} y \textquotedbl{}Basis, reduced\textquotedbl{}, no hay entrada Barlow). (2) El argumento de localización sí está y sí concluye que en elementos de orden bajo los puntos óptimos son los de la regla de Gauss de orden 2, a (h/2)/$\surd$3 del centro --- que es lo que la tesis llama puntos de Barlow. (3) Hay una tensión con la página anterior: en...}
\item \textbf{[PARCIAL]} tesis/capitulos/04\_resultados.tex:58 --- \textquotedbl{}Que el Q4 supere el orden O(h\textasciicircum{}1) del gradiente crudo refleja la superconvergencia del muestreo en los puntos de Barlow y del promediado nodal\textquotedbl{}.
\\ \campo{El respaldo cubre la mitad de la afirmación. Cook atribuye la ganancia de orden en el Q4 (hasta O(h²)) a la recuperación por parches (patch recovery, §9.9), no al promediado nodal simple que emplea EduFEM. Al contrario, en p. 231 Cook advierte para su ejemplo: \textquotedbl{}It is also too large at these locations in adjacent elements on either side, so that nodal averaging at shared nodes is of no benefit to accuracy.\textquotedbl{} El muestreo superconvergente en puntos de Gauss...}
\end{itemize}
{\small\campo{Sin respaldo localizable:} 1) tesis/capitulos/02\_marco\_teorico.tex:348 --- \textquotedbl{}las tensiones principales y la de von Mises, al ser funciones no lineales de ellas, se recalculan a partir de las componentes ya extrapoladas en lugar de extrapolarse de forma directa \textbackslash{}autocite\{cook2002concepts\}\textquotedbl{}. NO encontré respaldo para esta prescripción en Cook. Lo que sí dice la fuente, en la página impresa 232 (PDF 246), apunta más bien en la dirección contraria: \textquotedbl{}Examining the von Mises stress $\sigma$e (Eq. 3.12-2) at a location of stress concentration and using eight-node quadrilateral elements, Tenchev [6.18] found that the von Mises stress $\sigma$e calculated directly at a node had about one-third the error of $\sigma$e calculated at the node by extrapolation from four Gauss points\textquotedbl{}; y en p. 233 Cook recoge la fórmula empírica de Tenchev (Eq. 6.10-6) que combina ambos valores, $\sigma$*eN = ($\sigma$eN + $\sigma$eG)/(3 $\-$ $\sigma$eN/$\sigma$eG). Es decir, Cook compara \textquotedbl{}von Mises calculada directamente en el nodo\textquotedbl{} contra \textquotedbl{}von Mises extrapolada desde los puntos de Gauss\textquotedbl{}, pero no enuncia la regla de la tesis (extrapolar sólo las componentes cartesianas y recomponer después los invariantes). La decisión de EduFEM es defendible, pero no es de Cook: o se le busca otra fuente, o se presenta como criterio propio.

2) tesis/capitulos/02\_marco\_teorico.tex:408 y tesis/capitulos/04\_resultados.tex:150 --- el valor de referencia u\_y\textasciicircum{}ref = 23,96 de la membrana de Cook, y la caracterización del caso como \textquotedbl{}banco clásico para evaluar la sensibilidad a la distorsión\textquotedbl{}. NO figura en esta obra....\par}
\subsection{\texttt{csi2017sap2000}}
\begin{itemize}\small
\item \textbf{[PARCIAL]} tesis/capitulos/02\_marco\_teorico.tex:32 (fila de la tabla comparativa) --- «Software comercial (SAP2000 \textbackslash{}autocite\{csi2017sap2000\}, ANSYS, Abaqus) \& Inglés (localización parcial en algunos productos) \& 2D/3D \& Baja (caja negra) \& Comercial \& Amplios \& No (no didáctico)»
\\ \campo{Respaldan los atributos verificables: idioma inglés (todo el manual, 569 pp., está en inglés) y alcance 2D/3D (elemento Plane en tensión/deformación plana, p. 216; cáscaras tridimensionales, p. 178). NO respalda ---ni puede respaldar--- la celda «Baja (caja negra)» ni «No (no didáctico)»: son juicios del autor de la tesis sobre la experiencia de uso, no afirmaciones del manual. De hecho el manual documenta abiertamente la formulación (integración 2x2,...}
\item \textbf{[PARCIAL]} tesis/capitulos/02\_marco\_teorico.tex:11 --- «Por otro, los paquetes comerciales de análisis estructural \textbackslash{}autocite\{csi2017sap2000\}, que operan desde la perspectiva del estudiante como cajas negras: entregan resultados sin exponer las matrices intermedias que los produjeron.»
\\ \campo{La cita funciona como identificación del producto (encuadre), y para eso es correcta. La afirmación sustantiva ---«no exponen las matrices intermedias»--- no está en el manual y no puede estarlo: el manual describe la formulación pero no describe qué muestra la interfaz. Es una observación del autor sobre el uso del programa, no un dato citable de esta fuente. Sugerencia: o bien reformular como afirmación propia sin cita, o bien mantener la cita explicitando...}
\item \textbf{[PARCIAL]} tesis/capitulos/02\_marco\_teorico.tex:408 --- «comparando además contra un modelo de cáscara de SAP2000 \textbackslash{}autocite\{csi2017sap2000\}» (y, en paralelo, 05\_conclusiones.tex:19 --- «un modelo equivalente en el software comercial SAP2000 \textbackslash{}autocite\{timoshenko1970elasticity,csi2017sap2000\}: [...] frente a SAP2000, del 0,21 \% en $\sigma$x y de hasta...
\\ \campo{El manual respalda la existencia y naturaleza del elemento de cáscara empleado, es decir, identifica correctamente la herramienta y el tipo de elemento. Los VALORES numéricos citados junto a la clave (0,04 \%, 0,26 \%, 2,9 \%, 0,21 \%, 0,56 \%) NO provienen de esta fuente: son resultados propios del contraste hecho en la tesis. La cita es de atribución del software, no de los datos; está bien ubicada tras «SAP2000» y no tras las cifras, así que no induce a...}
\end{itemize}
{\small\campo{Sin respaldo localizable:} AFIRMACIÓN SIN RESPALDO EN ESTA FUENTE --- tesis/capitulos/03\_diseno\_implementacion.tex:169: «Todas las representaciones de campo emplean la paleta de color \textbackslash{}emph\{jet\}, el arcoíris clásico de los programas comerciales de análisis estructural \textbackslash{}autocite\{csi2017sap2000\}, elegida por su inmediato reconocimiento en el ámbito de la ingeniería.» respalda=NO.

Motivo: el CSI Analysis Reference Manual es un manual de FORMULACIÓN (teoría de elementos, cargas, casos de análisis), no un manual de la interfaz gráfica ni del post-proceso. Se recorrieron las 569 páginas con búsqueda de texto completo (capa de texto presente y legible) y NO hay una sola ocurrencia de «contour», «color», «colour», «rainbow», «palette», «legend», «spectrum» ni «plot». La única mención tangencial a la interfaz es «Graphical User Interface» en pp. impresas 38, 183 y 4, siempre referida al mallado automático o a la tolerancia de fusión de nodos, nunca a la representación cromática de resultados.

Es decir: la afirmación de que la paleta arcoíris/jet es la convención cromática de los programas comerciales NO es localizable en esta fuente y no puede serlo. Recomendación: (a) quitar la cita y dejar la afirmación como observación del autor, o (b) sustituirla por una fuente que sí trate el tema ---el manual de usuario/verification o la documentación de post-proceso de SAP2000, o bien literatura sobre mapas de color arcoíris (Borland \& Taylor, Moreland, Rogowitz \& Treinish), que además respaldaría de forma directa el...\par}
\subsection{\texttt{garciacordoba2005tecnologica}}
\begin{itemize}\small
\item \textbf{[PARCIAL]} \textquotedbl{}por su finalidad, el trabajo es de carácter propositivo ---la propuesta de un artefacto que resuelve un problema formativo---; por su naturaleza se inscribe en la investigación aplicada de tipo tecnológico, con enfoque cuantitativo en su fase de verificación y validación\textquotedbl{} --- c:/Tesis/Tesis...
\\ \campo{Respalda con claridad dos piezas: el carácter PROPOSITIVO de la investigación tecnológica (p. 91, término literal) y su definición como búsqueda de conocimiento operativo sobre una realidad concreta. Pero NO respalda la formulación 'investigación aplicada de tipo tecnológico': García Córdoba trata 'aplicada' y 'tecnológica' como cuatro tipos PARALELOS y las contrasta expresamente ('En contraste con la investigación aplicada, ésta parte de un...}
\end{itemize}
{\small\campo{Sin respaldo localizable:} 1) \textquotedbl{}con enfoque CUANTITATIVO en su fase de verificación y validación\textquotedbl{} (02b\_diseno\_metodologico.tex:4). No encontré respaldo en esta fuente y hay indicios en contra: la tesis central del libro, expuesta en el prólogo, es \textquotedbl{}Ni cuantitativo ni cualitativo: Tecnológico\textquotedbl{} (conferencia magistral del autor citada en la p. impresa 11, PDF 13: \textquotedbl{}si partimos de la suposición de que se investiga no por intereses científicos o por mera contemplación, sino porque hay un deseo de transformar, entonces, me parece que ni lo cuantitativo ni lo cualitativo es suficiente, es necesario lo tecnológico\textquotedbl{}). El respaldo del enfoque cuantitativo debe recaer sobre sampieri2018metodologia, no sobre esta clave. Revisé el índice completo (pp. impresas 17-23) y no hay apartado sobre enfoques cuantitativo/cualitativo fuera de esa mención del prólogo.

2) \textquotedbl{}investigación aplicada de tipo tecnológico\textquotedbl{} como género y especie: no respaldado; la fuente opone ambos tipos (ver comentario del respaldo PARCIAL, p. impresa 92).

3) Las dos citas del capítulo 1 (líneas 11 y 44 de 01\_introduccion.tex) y la del capítulo 2b van SIN localizador de página. Las tres son afirmaciones conceptuales concretas, no encuadre genérico, así que convendría añadir \textbackslash{}autocite[p.\textasciitilde{}83], \textbackslash{}autocite[p.\textasciitilde{}85] y \textbackslash{}autocite[p.\textasciitilde{}91] respectivamente.

Advertencia metodológica sobre el ejemplar: el PDF es un escaneo sin capa de texto (pdftotext no devuelve nada); todas las páginas se leyeron como imagen y el número impreso se verificó al pie de cada una....\par}
\subsection{\texttt{harris2020numpy}}
\begin{itemize}\small
\item \textbf{[PARCIAL]} «Para mallas grandes \$\textbackslash{}bm\{K\}\$ es dispersa (mayoritariamente nula), por lo que EduFEM la construye en formato de coordenadas (COO) y la convierte a filas comprimidas (CSR) para el solucionador, apoyándose en las bibliotecas científicas de Python \textbackslash{}autocite\{harris2020numpy,virtanen2020scipy\}» ---...
\\ \campo{El único respaldo que esta fuente ofrece para la parte dispersa es esa frase de la p. 360, y en ella el artículo atribuye las matrices dispersas a SciPy y PyData/Sparse, NO a NumPy: el propio texto aclara antes en esa misma página que «NumPy provides in-memory, multidimensional, homogeneously typed (that is, single-pointer and strided) arrays on CPUs». Los formatos COO y CSR, la conversión entre ellos y el solucionador disperso NO aparecen en ninguna de...}
\item \textbf{[PARCIAL]} «El núcleo numérico se apoya en NumPy y SciPy, que proveen el álgebra densa y dispersa, la cuadratura y el solucionador lineal del método \textbackslash{}autocite\{harris2020numpy,virtanen2020scipy\}» --- tesis/capitulos/03\_diseno\_implementacion.tex:25 (y la fila «Cálculo numérico \textbar{} NumPy, SciPy \textbar{} Álgebra, cuadratura, solucionador disperso» de...
\\ \campo{La afirmación es compuesta y la fuente cubre solo un tercio de ella. Álgebra DENSA acelerada: respaldada de forma literal en la p. 358 (BLAS vía OpenBLAS/Intel MKL). Álgebra DISPERSA, CUADRATURA y SOLUCIONADOR LINEAL: no figuran en el artículo de NumPy; el propio artículo los remite a SciPy («fundamental algorithms for scientific computing, including mathematical, scientific and engineering routines», p. 359) y a virtanen2020scipy, al que cita como su...}
\end{itemize}
{\small\campo{Sin respaldo localizable:} Ninguna cita queda enteramente sin respaldo, pero SÍ hay componentes concretos atribuidos a esta clave que NO existen en el artículo (verificado extrayendo el texto completo de las 6 páginas, impresas 357-362):

1) Formatos de matriz dispersa COO y CSR y la conversión entre ellos (02\_marco\_teorico.tex:270). No aparecen en ninguna página. El artículo solo menciona matrices dispersas una vez, en la p. 360, y para adjudicárselas a SciPy y PyData/Sparse, no a NumPy: «SciPy and PyData/Sparse both provide sparse arrays». Respalda = NO para esta parte.

2) Cuadratura (integración numérica / puntos de Gauss) (03\_diseno\_implementacion.tex:25 y tabla tab:stack). El artículo no menciona cuadratura ni integración numérica en ninguna de sus 6 páginas. Respalda = NO.

3) Solucionador lineal / solucionador disperso (03\_diseno\_implementacion.tex:25 y tabla tab:stack). El artículo no describe ningún solucionador de sistemas lineales; solo menciona «accelerated linear algebra» vía backends BLAS (p. 358), que es producto matricial y factorización densa, no un solver disperso. Respalda = NO.

Los tres puntos están cubiertos por la clave acompañante virtanen2020scipy, así que la oración de la tesis no es falsa, pero la atribución conjunta «NumPy y SciPy» carga sobre harris2020numpy afirmaciones que esta fuente no sostiene.

OBSERVACIONES ADICIONALES

- Ninguna de las 4 citas lleva localizador (\textbackslash{}autocite[p.\textasciitilde{}NN]). Para un artículo de revisión de 6 páginas es aceptable citarlo completo, y de hecho...\par}
\subsection{\texttt{hughes2000fem}}
\begin{itemize}\small
\item \textbf{[PARCIAL]} El integrando de la rigidez no es polinomico: como B depende de J\textasciicircum{}\{-1\}, que es racional en (xi,eta), la integral carece de primitiva elemental; por ello se aproxima mediante cuadratura de Gauss-Legendre. --- tesis/capitulos/02\_marco\_teorico.tex:235
\\ \campo{Hughes muestra que las derivadas globales (y por tanto B) llevan el factor 1/j, es decir que son racionales en (xi,eta), y afirma que la integracion exacta es 'practicamente imposible' para elementos isoparametricos salvo en las configuraciones mas simples. Lo que NO dice en ninguna parte, con esas palabras ni equivalentes, es que 'la integral carece de primitiva elemental'; esa formulacion es de la tesis. Sugiero atemperar la redaccion (p. ej. 'no admite...}
\item \textbf{[PARCIAL]} En la membrana de Cook el desplazamiento converge a 23,96, el Q4 con malla 8x8 da u\_y = 22,08 (7,9 \% por debajo) por shear-locking, y el Q9 lo evita y converge con rapidez. --- tesis/capitulos/02\_marco\_teorico.tex:408
\\ \campo{Hughes respalda unicamente la ultima clausula ('el Q9, de orden superior, lo evita y converge con rapidez'), y de forma cualitativa. Ni el caso de la membrana de Cook ni ninguna de las cifras (23,96 / 22,08 / 7,9 \%) aparecen en este libro: el indice alfabetico no tiene entrada 'Cook' ni 'membrane' fuera de 'Elastic membrane, 428' (que es un problema de valores propios). Esas cifras deben quedar atribuidas a cook2002concepts y a los propios resultados de...}
\end{itemize}
{\small\campo{Sin respaldo localizable:} 1) \textquotedbl{}la integral carece de primitiva elemental\textquotedbl{} (02\_marco\_teorico.tex:235): Hughes nunca lo dice asi. Lo mas proximo es la p. 191 (\textquotedbl{}all integrals need to be calculated exactly ... this would be virtually impossible for isoparametric elements in all but the simplest configurations\textquotedbl{}). Recomiendo reformular la frase de la tesis o acompanarla del localizador p. 191.

2) \textquotedbl{}El elemento Q9 ... evita el bloqueo [volumetrico] por la via del mayor orden de interpolacion\textquotedbl{} (02\_marco\_teorico.tex:257): esta afirmacion NO tiene respaldo en Hughes y la fuente apunta en sentido contrario para el caso volumetrico. La Figura 4.4.3 (p. 222) lista el cuadrilatero bicuadratico (nueve nodos) entre los elementos que, para el limite incompresible, requieren integracion selectiva (3x3 normal / 2x2 reducida) para ser equivalentes a la formulacion mixta; es decir, el Q9 completamente integrado tambien esta expuesto al bloqueo volumetrico. Para el bloqueo POR CORTANTE, en cambio, la nota al pie 10 de la p. 243 si respalda que el nueve nodos de Lagrange da excelentes resultados. Sugiero separar ambos casos en el texto, o no citar a Hughes para la clausula volumetrica.

3) La cifra \textquotedbl{}aproximadamente 90 \% de precision\textquotedbl{} no es de la tesis, pero conviene notar que el unico dato numerico de Hughes sobre el Q4 en flexion (p. 243, con nu = 0,3) es un 90 \% de exactitud, no el 7,9 \% de defecto que la tesis reporta para Cook; son casos distintos y no deben mezclarse en la misma cita.

4) Observacion bibliografica (no...\par}
\subsection{\texttt{lee2015eigenmodes}}
\begin{itemize}\small
\item \textbf{[PARCIAL]} «Lee desarrolla simuladores interactivos [...] de la visualización de los modos propios de vibración \textbackslash{}autocite\{lee2015eigenmodes\}» --- tesis/capitulos/02\_marco\_teorico.tex:11 (y, de refilón, «visualizadores de modos propios» en 04\_resultados.tex:223)
\\ \campo{PROBLEMA DE PRECISIÓN TÉCNICA. El artículo NO trata de modos propios de vibración: sus «eigenmodes» son los autovalores y autovectores de la MATRIZ DE RIGIDEZ (modos de desplazamiento y de energía de deformación: modo de cuerpo rígido, modo espurio de energía nula / hourglass, shear locking). Los modos de vibración exigen la matriz de masa y aparecen solo como extensión (ver siguiente entrada, p.\textasciitilde{}883). La redacción de 02\_marco\_teorico.tex:11 y...}
\item \textbf{[PARCIAL]} Matiz de la anterior: hasta qué punto la fuente cubre los modos de vibración (02\_marco\_teorico.tex:11)
\\ \campo{Esta es la única sección donde aparecen los modos de vibración, y explícitamente como EXTENSIÓN del caso estático. La p.\textasciitilde{}885 confirma que la comparación de modos de desplazamiento contra modos de vibración libre es un ejercicio asignado, no el objeto del artículo. Es decir: la afirmación de la tesis no es falsa de raíz, pero invierte el énfasis de la fuente.}
\item \textbf{[NO]} Fila de la tabla comparativa: «Simuladores interactivos \textbackslash{}autocite\{lee2015interactive,lee2015eigenmodes\} \& Inglés \& 2D \& Media-alta, en un eslabón \& Libre \& Continuo \& Parcial (un eslabón)» --- tesis/capitulos/02\_marco\_teorico.tex:30; columna DIMENSIÓN = «2D»
\\ \campo{CONTRADICE la columna «2D» de la tabla. La herramienta descrita opera con sólidos tridimensionales de 20 nodos, láminas degeneradas de 8 nodos y placas: p.\textasciitilde{}874 enumera los tipos «plane stress, plate bending, shell, etc.» y p.\textasciitilde{}882 compara mallas de elementos shell de 8 nodos. La celda debería decir «2D y 3D» al menos para esta referencia, o la fila debería desdoblarse.}
\item \textbf{[NO]} Misma fila de la tabla comparativa (02\_marco\_teorico.tex:30), columna LICENCIA = «Libre»
\\ \campo{INCONSISTENCIA INTERNA DE LA TESIS. La tabla clasifica estos simuladores como de licencia «Libre», pero el propio párrafo de 02\_marco\_teorico.tex:11, cuatro renglones antes, dice de VisualFEA: «al ser cerrados y de pago no se los pudo examinar de primera mano». Y el artículo deja claro que la simulación de modos propios está EMBEBIDA en VisualFEA (p.\textasciitilde{}872 y p.\textasciitilde{}876: «The educational tool was implemented based on the development framework of VisualFEA»), es...}
\item \textbf{[PARCIAL]} «En relación con las herramientas educativas reportadas en la literatura --simuladores de armaduras, visualizadores de modos propios o entornos de procesamiento interactivo de ecuaciones del MEF--, la contribución distintiva de EduFEM es la cobertura integral del canal de cálculo [...] y una memoria de cálculo con sustitución...
\\ \campo{La afirmación de brecha se sostiene en lo esencial ---el artículo cubre UN eslabón (los modos propios de la rigidez), no el canal completo, y no genera ningún documento de memoria de cálculo---, pero conviene un matiz de honestidad: la p.\textasciitilde{}885 muestra que VisualFEA SÍ despliega los valores numéricos de autovalores, autovectores y de la matriz de rigidez «in text format». La distinción de EduFEM no es exponer números intermedios (eso ya existe), sino que sean...}
\end{itemize}
{\small\campo{Sin respaldo localizable:} Dos celdas de la fila de la tabla comparativa (tesis/capitulos/02\_marco\_teorico.tex:30) no tienen respaldo en esta fuente y una de ellas contradice al propio texto de la tesis:

1) LICENCIA = «Libre». El artículo no consigna licencia, precio ni condiciones de distribución en ninguna de sus 15 páginas. La herramienta de simulación de modos propios está embebida en VisualFEA, programa comercial de los propios autores (p. 872, p. 876). Además el párrafo de 02\_marco\_teorico.tex:11 califica a VisualFEA de «cerrados y de pago», de modo que la tabla se contradice con su propio párrafo introductorio. Dato AUSENTE en la fuente.

2) DIMENSIÓN = «2D». Refutado por la fuente: p. 874 lista tipos de elemento «plane stress, plate bending, shell, etc.»; p. 881 y la Figura 7 documentan un elemento sólido TRIDIMENSIONAL de 20 nodos con integración 2x2x2 y 3x3x3; p. 882 usa láminas degeneradas de 8 nodos. La herramienta no es 2D.

3) Precisión terminológica (no es falta de respaldo, es error de atribución): «visualización de los modos propios de VIBRACIÓN» (02\_marco\_teorico.tex:11) y «visualizadores de modos propios» (04\_resultados.tex:223). El artículo trata los autovalores y autovectores de la MATRIZ DE RIGIDEZ ---modos de desplazamiento, de cuerpo rígido, espurios de energía nula (hourglass) y de shear locking---; los modos de vibración aparecen solo como sección de extensión en p. 883 y como ejercicio de comparación en p. 885. Recomiendo reescribir a «modos propios de la matriz de rigidez...\par}
\subsection{\texttt{lee2015interactive}}
\begin{itemize}\small
\item \textbf{[PARCIAL]} «La literatura sobre enseñanza del método coincide en que la simulación interactiva y la manipulación directa de los objetos matemáticos del MEF favorecen un aprendizaje más profundo que la mera exposición teórica o el uso de herramientas opacas» --- tesis/capitulos/01\_introduccion.tex:21
\\ \campo{La mitad afirmativa está bien respaldada: Lee sostiene expresamente que la simulación interactiva lleva a una mejor comprensión («better understanding», «more intuitive understanding»). Lo que NO está en la fuente es el término de comparación que la tesis introduce: Lee contrapone la simulación a la «manual calculation» y a la «computer programming» (p.157), nunca a «la mera exposición teórica» ni a «herramientas opacas» --- la expresión «caja negra» /...}
\item \textbf{[PARCIAL]} «las herramientas profesionales son cajas negras optimizadas para producir resultados, no para revelar el procedimiento intermedio, mientras que el desarrollo puramente manual no escala más allá de uno o dos elementos. La investigación educativa documenta esta brecha y propone el software específicamente didáctico como puente...
\\ \campo{La segunda mitad de la afirmación está respaldada casi palabra por palabra: el cálculo manual «no escala» («very limited in its extent») y por eso se desarrolla software didáctico como puente. La primera mitad NO: Lee nunca caracteriza al software comercial como caja negra; su par de opuestos es cálculo manual vs. programación, no cálculo manual vs. paquete comercial. Para la parte de «cajas negras» el respaldo correcto es perezsantiago2023fem, ya citado...}
\item \textbf{[NO]} «La segunda familia la integran los simuladores que profundizan en un eslabón aislado del método» --- tesis/capitulos/02\_marco\_teorico.tex:11 --- y la fila de la tabla comparativa «Simuladores interactivos ... Parcial (un eslabón)» --- 02\_marco\_teorico.tex:30
\\ \campo{Hallazgo importante: la fuente contradice esta caracterización. Lee cubre CUATRO eslabones consecutivos (funciones de forma y modelado del elemento $\rightarrow$ matriz de rigidez elemental por integración numérica $\rightarrow$ ensamblaje global $\rightarrow$ solución por eliminación de Gauss y sustitución regresiva), y reclama explícitamente ser «more comprehensive and integrated than that of previous works». Etiquetarlo como «un eslabón aislado» / «Parcial (un eslabón)» subestima la...}
\item \textbf{[NO]} Fila de la tabla comparativa: régimen de licencia «Libre» para los simuladores interactivos de Lee --- tesis/capitulos/02\_marco\_teorico.tex:30
\\ \campo{Hallazgo grave de coherencia interna. El artículo no dice nada sobre licencia, y lo que sí dice es que el simulador es una función de VisualFEA --- que la propia tesis describe seis líneas más arriba, en 02\_marco\_teorico.tex:11, como producto comercial: «al ser cerrados y de pago no se los pudo examinar de primera mano». La tabla afirma «Libre» para la misma herramienta que el texto del mismo capítulo llama de pago. Un tribunal puede detectarlo con solo...}
\item \textbf{[PARCIAL]} «El software educativo interactivo demuestra eficacia para consolidar conceptos abstractos del análisis estructural» --- tesis/capitulos/03\_diseno\_implementacion.tex:9
\\ \campo{El verbo «demuestra eficacia» es más fuerte de lo que la fuente sostiene. Lee dice expresamente que NO hay datos cuantitativos de eficacia y que su evaluación fue cualitativa y observacional, sobre sus propios cursos, sin grupo de control. Nótese que sus conclusiones usan «It has been demonstrated», pero apoyado en esa misma evaluación cualitativa. Es una debilidad citable por el tribunal, sobre todo porque la tesis en 01\_introduccion.tex:56 ya dice...}
\item \textbf{[PARCIAL]} «En relación con las herramientas educativas reportadas en la literatura ---simuladores de armaduras, visualizadores de modos propios o entornos de procesamiento interactivo de ecuaciones del MEF---, la contribución distintiva de EduFEM es la cobertura integral del canal de cálculo ... articulada con una generación automática de...
\\ \campo{La designación de Lee como «entorno de procesamiento interactivo de ecuaciones del MEF» es correcta. La afirmación de novedad de EduFEM se sostiene solo parcialmente frente a esta fuente: Lee no genera ningún documento de memoria de cálculo (lo más cercano son diapositivas armadas a mano con capturas, p.167), así que ese aspecto de la novedad SÍ resiste; pero la «cobertura integral del canal de cálculo» está más disputada de lo que la tesis admite, porque...}
\item \textbf{[NO]} «...todo ello con una orientación deliberadamente minimalista coherente con su propósito didáctico» --- tesis/capitulos/05\_conclusiones.tex:15
\\ \campo{Contradicción interna de la tesis, además de falta de respaldo: en 03\_diseno\_implementacion.tex:155 el propio texto declara que el minimalismo «es decisión de este trabajo, no un resultado tomado de la literatura», y aquí la conclusión se lo atribuye a la literatura. Solución: mover la cita en la oración de modo que respalde la selección bidireccional / deshacer / validación (todo eso sí es interactividad documentada), y dejar el minimalismo sin cita, o...}
\end{itemize}
{\small\campo{Sin respaldo localizable:} Dos afirmaciones citadas a esta clave carecen de respaldo localizable en el artículo, y una tercera es contradicha por él:

1. RÉGIMEN DE LICENCIA «Libre» en la fila «Simuladores interactivos» de \textbackslash{}autoref\{tab:comparativa\} (tesis/capitulos/02\_marco\_teorico.tex:30). El artículo no menciona licencia en ninguna de sus 13 páginas. Lo que sí consta es que el simulador es una función de VisualFEA (p.168 y biografía p.169), producto que la propia tesis califica de cerrado y de pago catorce líneas antes (02\_marco\_teorico.tex:11). La celda debería decir «Comercial».

2. ORIENTACIÓN MINIMALISTA atribuida a la literatura en tesis/capitulos/05\_conclusiones.tex:15. La palabra «minimal» no aparece en el artículo y la herramienta descrita es funcionalmente densa. Además contradice lo que la propia tesis declara en 03\_diseno\_implementacion.tex:155 («es decisión de este trabajo, no un resultado tomado de la literatura»).

3. CARACTERIZACIÓN COMO «UN ESLABÓN AISLADO» (02\_marco\_teorico.tex:11 y la celda «Parcial (un eslabón)» de la línea 30). No es que falte respaldo: la fuente afirma lo contrario en p.157-158 (cuatro etapas, de la función de forma a la solución del sistema, y reclamo explícito de ser «more comprehensive and integrated than that of previous works»).

Observaciones adicionales sobre el uso general de la clave:
- Ninguna de las 12 citas lleva localizador \textbackslash{}autocite[p.\textasciitilde{}NN]. Como esta fuente es de solo 13 páginas y muy específica, conviene añadirlo al menos en 04\_resultados.tex:207...\par}
\subsection{\texttt{oberkampf2010vv}}
\begin{itemize}\small
\item \textbf{[PARCIAL]} 02b\_diseno\_metodologico.tex:16 --- 'La evaluacion de la correccion se realizo mediante un esquema de verificacion y validacion: verificacion por el metodo de soluciones manufacturadas y VALIDACION contra la viga de Timoshenko y la membrana de Cook'.
\\ \campo{El esquema V\&V esta cubierto, pero la ETIQUETA no: para Oberkampf y Roy la validacion exige datos experimentales, de modo que Timoshenko (solucion analitica) es verificacion de codigo/solucion y SAP2000 es comparacion codigo a codigo. Reescritura: «La \textbackslash{}emph\{evaluacion de la correccion\} se realizo mediante un esquema de verificacion y validacion (\textbackslash{}autoref\{sec:verificacion-validacion\}): verificacion de codigo por el metodo de soluciones manufacturadas y...}
\item \textbf{[PARCIAL]} 05\_conclusiones.tex:19 --- 'El marco metodologico de verificacion y validacion adoptado sigue las practicas establecidas en la literatura de computacion cientifica'; y, en la misma linea, 'la validacion con la membrana de Cook'.
\\ \campo{La afirmacion sobre el marco metodologico queda cubierta con solo cambiar la clave a \textbackslash{}autocite\{oberkampf2010vv,salari2000mms\}. Lo que NO queda cubierto es llamar «validacion» al contraste con Timoshenko/SAP2000/Cook. Reescritura del tramo final del parrafo: «La validacion analitica se realizo...» -\textgreater{} «La contrastacion con referencias externas se realizo con la viga de Timoshenko...»; y «Finalmente, la validacion con la membrana de Cook...}
\end{itemize}
\subsection{\texttt{onate2009structural}}
\begin{itemize}\small
\item \textbf{[PARCIAL]} capitulos/05\_conclusiones.tex:47 --- «...y con cuadriláteros de transición, ofreciendo así una variedad que enriquezca el estudio comparativo de la convergencia y del comportamiento bajo distorsión de malla».
\\ \campo{PARCIAL por dos motivos. (1) Los \textquotedbl{}cuadriláteros de transición\textquotedbl{} aparecen en Oñate en una sola frase de paso (pág. 167) y con una valoración que la tesis no recoge: dice que son elementos \textquotedbl{}not very popular\textquotedbl{}. Oñate NO desarrolla una familia de cuadriláteros de transición. (2) El \textquotedbl{}comportamiento bajo distorsión de malla\textquotedbl{} sí está respaldado, pero en otras páginas: pág. impresa 193-194 (§6.2.2, el Q9 reproduce polinomios cuadráticos cartesianos en formas...}
\end{itemize}
{\small\campo{Sin respaldo localizable:} Ninguna afirmación quedó sin respaldo localizable: las tres citas de onate2009structural (01\_introduccion.tex:62, 02\_marco\_teorico.tex:101, 05\_conclusiones.tex:47) apuntan a contenido que existe en el volumen, y en dos de los tres casos el respaldo es literal y textualmente coincidente. Observaciones para el autor:

1) NINGUNA de las tres citas lleva localizador. Recomendado: 01\_introduccion.tex:62 -\textgreater{} \textbackslash{}autocite[p.\textasciitilde{}164, 183-184, 193]\{onate2009structural\}; 02\_marco\_teorico.tex:101 -\textgreater{} \textbackslash{}autocite[p.\textasciitilde{}187]\{onate2009structural\} (y p.\textasciitilde{}159-160 para el cuadrado natural, p.\textasciitilde{}161 y 164 para las funciones de forma del Q9); 05\_conclusiones.tex:47 -\textgreater{} \textbackslash{}autocite[p.\textasciitilde{}175-177]\{onate2009structural\} para triángulos y \textbackslash{}autocite[p.\textasciitilde{}214]\{onate2009structural\} para la distorsión de malla.

2) Único punto flojo: \textquotedbl{}cuadriláteros de transición\textquotedbl{} (05\_conclusiones.tex:47). En Oñate el respaldo es una frase incidental de la pág. impresa 167 que además los califica de \textquotedbl{}not very popular\textquotedbl{}. Si el autor quiere sostener esa línea de trabajo futuro, o bien la matiza, o bien la apoya en otra fuente; alternativamente puede reencuadrarla como \textquotedbl{}unión de elementos disímiles\textquotedbl{} (§9.3, pág. impresa 311), que sí es una sección propia del libro.

3) Verificación de paginación: el desfase impreso = pdf - 24 se comprobó leyendo el encabezado impreso de cada página citada (pdf 188 -\textgreater{} \textquotedbl{}164 Higher order 2D solid elements\textquotedbl{}; pdf 191 -\textgreater{} \textquotedbl{}Serendipity rectangular elements 167\textquotedbl{}; pdf 207 -\textgreater{} \textquotedbl{}...183\textquotedbl{}; pdf 208 -\textgreater{} \textquotedbl{}184 Higher order 2D solid...\par}
\subsection{\texttt{perezsantiago2023fem}}
{\small\campo{Sin respaldo localizable:} La afirmacion de que la carencia conceptual \textquotedbl{}dificulta diagnosticar resultados anomalos o evaluar la calidad de una malla\textquotedbl{} no aparece formulada asi en el articulo. Por eso la prosa se reescribio el 2026-09-09 a lo que la fuente si sostiene: que el egresado debe saber planificar, verificar y validar su analisis, no solo operarlo.\par}
\subsection{\texttt{reddy2006introduction}}
\begin{itemize}\small
\item \textbf{[PARCIAL]} Encuadre: \textquotedbl{}los textos clásicos desarrollan la teoría con un grado de abstracción considerable, donde las funciones de forma, la matriz Jacobiana, la matriz de deformación-desplazamiento y la integración de Gauss aparecen como pasos de una cadena algebraica cuya conexión con la geometría concreta de un elemento no siempre...
\\ \campo{Cita de mero encuadre (pila de textos clásicos): no se le atribuye a Reddy ninguna afirmación verificable, sino que se lo usa como ejemplar del género 'texto clásico formal'. El uso es legítimo --- el libro efectivamente presenta funciones de forma, jacobiano y cuadratura como cadena algebraica abstracta (cap. 9), como muestra el pasaje transcrito --- pero no hay proposición concreta que comprobar. No requiere localizador.}
\item \textbf{[PARCIAL]} \textquotedbl{}...ofreciendo así una variedad que enriquezca el estudio comparativo de la convergencia y del comportamiento bajo distorsión de malla\textquotedbl{} --- tesis/capitulos/05\_conclusiones.tex:47
\\ \campo{Respalda la parte de 'comportamiento bajo distorsión de malla' (§9.4.2 'Element Geometries', p. impresa 562-563, trata distorsión y ángulos inadmisibles por tipo de elemento, incluido el triangular cuadrático). NO encontré en Reddy un estudio comparativo de convergencia entre tipos de elemento triangulares vs. cuadriláteros; §14.5.3 'Convergence and Accuracy of Solutions' (p. impresa 745) trata convergencia en general, no la comparación por tipo de...}
\end{itemize}
{\small\campo{Sin respaldo localizable:} Ninguna afirmación quedó completamente sin respaldo. Matices y observaciones:

1) \textquotedbl{}estudio comparativo de la convergencia\textquotedbl{} (05\_conclusiones.tex:47) --- no localicé en Reddy una comparación de convergencia entre familias de elementos que sostenga esa mitad de la frase; solo el tratamiento general de convergencia y errores del cap. 14 (§14.5, p. impresa 743-750) y el de distorsión geométrica del §9.4.2 (p. impresa 562-563). La marco PARCIAL, no NO, porque la frase es prospectiva (\textquotedbl{}puede ampliarse... ofreciendo así\textquotedbl{}) y el resto sí está respaldado.

2) La cita de 01\_introduccion.tex:9 es de encuadre puro (Reddy como ejemplar de \textquotedbl{}texto clásico\textquotedbl{}); no hay proposición verificable que comprobar y no requiere localizador.

3) NINGUNA de las tres citas trae localizador (\textbackslash{}\textbackslash{}autocite[p.\textasciitilde{}NN]). Recomiendo agregarlos, ahora que están verificados:
   - 02\_marco\_teorico.tex:49 $\rightarrow$ \textbackslash{}\textbackslash{}autocite[p.\textasciitilde{}612]\{reddy2006introduction\} (o pp. 612-613)
   - 05\_conclusiones.tex:47 $\rightarrow$ \textbackslash{}\textbackslash{}autocite[pp.\textasciitilde{}525-526, 566]\{reddy2006introduction\}
   - 01\_introduccion.tex:9 $\rightarrow$ sin localizador (encuadre)

4) PROBLEMA DE FICHA (no de respaldo), coherente con el veredicto DISCREPA de la Tarea 1: la entrada .bib declara location = \{New York\} e isbn = \{978-0-07-246685-0\}, pero el ejemplar consultado es la \textquotedbl{}International Edition 2006\textquotedbl{} de McGraw-Hill Education (Asia), impresa en Singapur. El ISBN 978-0-07-246685-0 corresponde a la edición norteamericana, no al ejemplar en mano; no lo verifiqué impreso en este PDF. Si se quiere que la...\par}
\subsection{\texttt{salari2000mms}}
\begin{itemize}\small
\item \textbf{[PARCIAL]} El error se cuantifica con normas de error globales y las normas deben decrecer al refinar h con la tasa teorica, lo que confirma la correcta implementacion --- tesis/capitulos/02\_marco\_teorico.tex:396-406 y 04\_resultados.tex:16
\\ \campo{La fuente respalda el esquema general: norma L2 del error frente a la solucion manufacturada, serie de al menos dos mallas, y comparacion del orden observado contra el teorico. NO respalda las tasas concretas O(h\textasciicircum{}\{p+1\}) en L2 y O(h\textasciicircum{}p) en H1 para elementos de orden p --- esas la tesis las cita correctamente a strang2008analysis/hughes2000fem/zienkiewicz2013fem. Ademas, el reporte usa norma RMS/L2 discreta y norma infinito, no la seminorma H1 con ese nombre.}
\end{itemize}
{\small\campo{Sin respaldo localizable:} HALLAZGO PRINCIPAL --- 1 afirmacion sin respaldo en esta fuente:

(1) tesis/capitulos/02\_marco\_teorico.tex:406 --- \textquotedbl{}Las integrales de error se evaluan con un punto de Gauss mas por direccion que la rigidez (\$3\textbackslash{}times3\$ para Q4, \$4\textbackslash{}times4\$ para Q9) para no subestimar artificialmente el error \textbackslash{}autocite\{salari2000mms\}.\textquotedbl{}  respalda = NO.

Esta afirmacion NO tiene respaldo localizable en \textquotedbl{}Code Verification by the Method of Manufactured Solutions\textquotedbl{} (SAND2000-1444). Verificacion realizada sobre el texto completo de las 125 paginas del PDF (tiene capa de texto; extraido con pdftotext -layout):
  - grep -i \textquotedbl{}gauss\textquotedbl{}          -\textgreater{} 0 apariciones
  - grep -i \textquotedbl{}quadratur\textquotedbl{}      -\textgreater{} 0 apariciones
  - grep -i \textquotedbl{}integration point\textquotedbl{} -\textgreater{} 0 apariciones
  - grep -i \textquotedbl{}alias\textquotedbl{}          -\textgreater{} 0 apariciones
  - grep -i \textquotedbl{}integrat\textquotedbl{}       -\textgreater{} 3 apariciones, ninguna pertinente: p.\textasciitilde{}19 impresa (dificultades de integrar numericamente soluciones exactas con singularidades, en la critica al metodo MES) y p.\textasciitilde{}57 (Apendice B, integracion de la Ec. B11).
  - grep -i \textquotedbl{}finite element\textquotedbl{} -\textgreater{} 6 apariciones, todas de mera enumeracion de familias de metodos; el reporte no desarrolla ningun ejemplo en elementos finitos.

El reporte es de orientacion CFD/diferencias finitas: sus casos trabajados son conduccion de calor, ecuacion de Burgers 2D y Navier-Stokes incompresible y compresible sobre mallas estructuradas. No discute cuadratura de Gauss, ni orden de integracion de las integrales de error, ni el aliasing del error con la cuadratura de la...\par}
\subsection{\texttt{sampieri2018metodologia}}
\begin{itemize}\small
\item \textbf{[PARCIAL]} El trabajo tiene «enfoque cuantitativo en su fase de verificación y validación», entendido como establecer la calidad del artefacto «midiendo errores numéricos, tensiones, desplazamientos y tasas de convergencia frente a referencias de exactitud conocida» --- tesis/capitulos/02b\_diseno\_metodologico.tex:4 (cita conjunta...
\\ \campo{Sampieri respalda SOLO la mitad «cuantitativa» de la frase: datos numéricos, medición y prueba de hipótesis. Complemento en p. impresa 7 (PDF 48): «5. Ya que los datos son numéricos se deben analizar con métodos estadísticos» y «Algunas de las características esenciales del enfoque cuantitativo son: 1. Búsqueda de la mayor objetividad posible en todo el proceso o ruta». En cambio, los otros tres rótulos de la misma oración ---«propositivo», «investigación...}
\end{itemize}
{\small\campo{Sin respaldo localizable:} NINGUNA afirmación queda sin respaldo, con una salvedad parcial ya detallada arriba. Detalle:

(1) SALVEDAD en 02b\_diseno\_metodologico.tex:4 --- la oración atribuye a la pareja de fuentes cuatro rótulos: «propositivo», «investigación aplicada», «tecnológico» y «enfoque cuantitativo». Sampieri respalda solo el cuarto. Los tres primeros no existen en este libro: la búsqueda de las cadenas «investigación aplicada», «aplicada y básica» y «básica y aplicada» sobre el texto completo extraído del PDF (763 páginas, 47.453 líneas) devuelve cero coincidencias; el libro organiza la disciplina por rutas (cuantitativa/cualitativa/mixta) y por alcances (exploratorio, descriptivo, correlacional, explicativo), nunca por la dicotomía básica/aplicada ni por «investigación tecnológica». Esa carga la lleva íntegramente garciacordoba2005tecnologica. No es una cita falsa ---la cita es conjunta y Sampieri sí cubre lo cuantitativo---, pero si el tribunal pide localizador conviene desdoblarla.

(2) No busqué respaldo para el uso de «variable independiente / dependiente / controlada» de la línea 21 más allá de la operacionalización, porque la cita de esa oración está anclada explícitamente a la operacionalización («especifica para cada una las operaciones concretas con que se la mide») y ese punto sí está respaldado en p. 137.

(3) Ninguna de las tres citas trae localizador \textbackslash{}autocite[p.\textasciitilde{}NN] en el .tex, así que no hubo páginas declaradas que comprobar; las tres páginas que doy son hallazgos míos, leídas en...\par}
\subsection{\texttt{stimpson2007verdict}}
\begin{itemize}\small
\item \textbf{[PARCIAL]} «De la documentación de Verdict provienen los mínimos admisibles, \$q\_\{SJ\}\textbackslash{}ge 0\{,\}50\$ y \$Q\textbackslash{}ge 0\{,\}25\$» --- tesis/capitulos/02\_marco\_teorico.tex:376
\\ \campo{El minimo Q \textgreater{}= 0,25 esta impreso tal cual en la pag. 57. El minimo q\_SJ \textgreater{}= 0,50 NO figura para cuadrilateros: Verdict fija el Acceptable Range del scaled Jacobian del QUAD en [0.3, 1] (pag. 51). El 0,5 corresponde a otras formas: hexaedro [0.5, 1] (pag. 92) y triangulo [0.5, 2sqrt(3)/3] (pag. 29). Como EduFEM solo maneja Q4/Q9 (cuadrilateros), la cifra atribuible a esta fuente es 0,30, no 0,50. El codigo usa 0.5 (education/mod00\_*.py:91: \_SJ\_ACCEPT,...}
\item \textbf{[PARCIAL]} «el estrechamiento (taper) \$T\_R\$, que es nulo en todo paralelogramo y crece con la trapecialidad» --- tesis/capitulos/02\_marco\_teorico.tex:376 (parrafo posterior a eq:descomposicion-bilineal)
\\ \campo{Verdict respalda el concepto (taper = razon entre la derivada cruzada X12 y los ejes principales) y el hecho de que se anula en el cuadrado (q for unit square: 0; en general X12 = 0 en cualquier paralelogramo, por construccion). La formula concreta de EduFEM, T\_R = max(\textbar{}X3$\cdot$X1\textbar{}/\textbar{}\textbar{}X1\textbar{}\textbar{}\textasciicircum{}2, \textbar{}X3$\cdot$X2\textbar{}/\textbar{}\textbar{}X2\textbar{}\textbar{}\textasciicircum{}2), es una proyeccion normalizada distinta de la de Verdict (\textbar{}\textbar{}X12\textbar{}\textbar{}/min(\textbar{}\textbar{}X1\textbar{}\textbar{},\textbar{}\textbar{}X2\textbar{}\textbar{})). La tesis no la atribuye a Verdict y declara sus cortes (0,3 / 0,5)...}
\item \textbf{[PARCIAL]} «El módulo M0 atiende la calidad de malla en el pre-proceso, evaluando el Jacobiano escalado y la compacidad (stretch) con los umbrales bibliográficos de Verdict» --- tesis/capitulos/04\_resultados.tex:209
\\ \campo{Las dos metricas nombradas (scaled Jacobian §5.15 y stretch §5.21) existen en Verdict y el umbral de compacidad 0,25 es bibliografico. El umbral de Jacobiano escalado que implementa M0 (0,5) no es el de Verdict para cuadrilateros (0,3), asi que la frase \textquotedbl{}con los umbrales bibliograficos de Verdict\textquotedbl{} es exacta solo para uno de los dos. Mismo arreglo que en 02\_marco\_teorico.tex:376.}
\end{itemize}
{\small\campo{Sin respaldo localizable:} Una sola afirmacion queda sin respaldo localizable en esta fuente: el minimo admisible q\_SJ \textgreater{}= 0,50 para el Jacobiano escalado (tesis/capitulos/02\_marco\_teorico.tex:376, replicado como \textquotedbl{}umbrales bibliograficos de Verdict\textquotedbl{} en tesis/capitulos/04\_resultados.tex:209). Revise las tres secciones de Scaled Jacobian del informe: cuadrilatero §5.15 (pag. impresa 51) da Acceptable Range [0.3, 1]; triangulo §4.9 (pag. 29) da [0.5, 2sqrt(3)/3]; hexaedro §7.11 (pag. 92) da [0.5, 1]. No hay ninguna otra tabla ni pasaje del informe que prescriba 0,50 para cuadrilateros, ni en §1.4 Metric Ranges (pag. 12) ni en la introduccion del capitulo 5 (pags. 35-36). Como EduFEM solo trabaja con Q4/Q9, la cifra citable es 0,30. Recomendacion: o corregir el texto a \textquotedbl{}q\_SJ \textgreater{}= 0,30 y Q \textgreater{}= 0,25 \textbackslash{}\textbackslash{}autocite[pp.\textasciitilde{}51, 57]\{stimpson2007verdict\}\textquotedbl{}, o mantener el 0,50 del codigo declarandolo convencion mas estricta de EduFEM (como ya se hace con los cortes superiores 0,80 y 0,50) y ajustar tambien la frase de 04\_resultados.tex:209. Observacion adicional, no bloqueante: ninguna de las tres citas lleva localizador \textbackslash{}\textbackslash{}autocite[p.\textasciitilde{}NN]; conviene agregarlo (pag. 3 para el encuadre, 51 y 57 para las metricas, 35-36 para la descomposicion bilineal), dado que el informe tiene 108 paginas.\par}
\subsection{\texttt{strang2008analysis}}
\begin{itemize}\small
\item \textbf{[PARCIAL]} Respaldo secundario del mismo enunciado: la estimación de aproximación en normas de media cuadrática que sustenta las tasas h\textasciicircum{}\{k-s\} (Teorema 3.3)
\\ \campo{Es el teorema de interpolación que da la cota h\textasciicircum{}\{k-s\} de manera rigurosa, pero para el interpolante u\_I, no para la solución de Ritz u\textasciicircum{}h; el paso de uno a otro es lo que hace §2.2. Sirve como refuerzo, no como cita principal. Si se busca una sola página, la correcta es la 107.}
\end{itemize}
{\small\campo{Sin respaldo localizable:} Ninguna afirmación queda sin respaldo: las dos apariciones de la clave en los capítulos (02\_marco\_teorico.tex:406 y 04\_resultados.tex:16) atribuyen a esta fuente exactamente la misma afirmación ---las tasas asintóticas O(h\textasciicircum{}\{p+1\}) en L2 y O(h\textasciicircum{}p) en H1--- y ambas están sostenidas por §2.2 \textquotedbl{}Rates of Convergence\textquotedbl{}, pp. 106-107. En ninguno de los dos sitios hay localizador \textbackslash{}autocite[p.\textasciitilde{}NN], de modo que no hubo página que comprobar: la correcta es p. 107 (o el par 106-107).

ADVERTENCIAS QUE AFECTAN LA CITA, no el contenido:

1. DESFASE DE EDICIÓN (ya marcado como DISCREPA en la Tarea 1). El .bib declara edition = \{2\}, Wellesley-Cambridge Press, Wellesley MA, 2008. El ejemplar disponible en tesis/bibliografia/ es la PRIMERA edición: Prentice-Hall, Inc., serie Prentice-Hall Series in Automatic Computation (editor George Forsythe), 1973, sin mención de edición impresa. TODAS las páginas que doy arriba son de ese ejemplar de 1973. NO puedo verificar que la paginación de la edición de 2008 coincida, porque no tengo ese ejemplar. Por tanto: si la tesis mantiene la referencia a la edición de 2008, no debe citarse \textquotedbl{}p. 107\textquotedbl{} sin antes cotejar ese ejemplar; lo consistente es o bien citar la edición de 1973 que efectivamente se consultó, o bien conseguir la de 2008 y verificar la página allí.

2. Ninguna de las dos citas es de mero encuadre, pese a ir en pila con hughes2000fem / zienkiewicz2013fem: en ambos casos se atribuye una afirmación concreta (dos exponentes). Strang y Fix la sostienen,...\par}
\subsection{\texttt{suarez1998edelas2d}}
\begin{itemize}\small
\item \textbf{[PARCIAL]} ED-Elas2D está «disponible en español» y comparte con EduFEM «el idioma»; celda de tabla «Idioma: Español» (02\_marco\_teorico.tex:13 y :28)
\\ \campo{El artículo NUNCA afirma en prosa que el programa esté en español ni menciona idiomas o versiones lingüísticas. La evidencia es solo indirecta y mixta, procedente de las capturas: en español aparecen la barra del visor de ayuda (p. 246), el rótulo «Puntos de Gauss 1» dentro del diálogo de cálculo (p. impresa 249, PDF 8) y el título de ventana «Nivel3 - ED-Elas2D Postprocessing» (p. impresa 253, PDF 12); en inglés aparecen todos los menús y botones del...}
\item \textbf{[PARCIAL]} «software académico cerrado»; celda de tabla «Licencia: Académico» (02\_marco\_teorico.tex:13, :28 y :30)
\\ \campo{El artículo no declara licencia alguna, ni «cerrado», ni «propietario», ni «código abierto». Lo único impreso es el origen académico/institucional (UPC, CIMNE, financiación COMETT y Leonardo da Vinci) y que se distribuye como producto de CIMNE. «Académico» es una inferencia razonable de esa procedencia; «cerrado» es una inferencia del silencio (no hay mención de disponibilidad del código) y no está sostenida por ningún pasaje.}
\end{itemize}
{\small\campo{Sin respaldo localizable:} Ninguna de las tres citas del .tex lleva localizador (`\textbackslash{}autocite\{suarez1998edelas2d\}` sin `[p.\textasciitilde{}NN]`), así que no hubo páginas declaradas que comprobar; las páginas de esta tabla son las que localicé yo.

Afirmaciones sin respaldo localizable, o con respaldo solo indirecto, en esta fuente:

1. «disponible en español» / «comparte [...] también el idioma» (02\_marco\_teorico.tex:13) y la celda «Idioma: Español» de la tabla (02\_marco\_teorico.tex:28) --- el artículo, escrito íntegramente en inglés, no dice en ninguna parte que el programa esté en español. Solo se ve castellano en tres detalles de capturas (barra del visor de ayuda de Windows en p. 246; «Puntos de Gauss 1» en p. 249; «Nivel3» en el título de ventana de p. 253), mientras que los menús del pre y post-proceso y el índice teórico aparecen en inglés. Marcado PARCIAL arriba: la afirmación tal como está escrita («disponible en español», sin matices) excede lo que el artículo sostiene.

2. «software académico cerrado» y la celda «Licencia: Académico» (02\_marco\_teorico.tex:13, :28, :30) --- no hay ningún pasaje sobre licencia, distribución, precio ni disponibilidad del código fuente. Solo consta la filiación UPC/CIMNE y la financiación COMETT y Leonardo da Vinci (agradecimientos, p. 254). Marcado PARCIAL: «académico» se infiere de la procedencia, «cerrado» no tiene respaldo textual.

3. Nota metodológica sobre el PDF: solo las páginas PDF 1, 2 y 13 tienen capa de texto (páginas impresas 243 y 254). Las páginas PDF 3-12 (impresas...\par}
\subsection{\texttt{timoshenko1970elasticity}}
\begin{itemize}\small
\item \textbf{[PARCIAL]} A partir de las componentes cartesianas se obtienen las tensiones principales y la equivalente de von Mises, la medida escalar con que habitualmente se resume un estado tensional biaxial. tesis/capitulos/02\_marco\_teorico.tex:298
\\ \campo{Se separan dos cosas. (a) Tensiones principales: respaldadas, art. 9, p. 14, ec. (14) --- pero Timoshenko las define por la condicion tau = 0 y las construye con el circulo de Mohr; NO escribe la forma cerrada sigma\_\{1,2\} = (sigma\_x+sigma\_y)/2 +- sqrt(((sigma\_x-sigma\_y)/2)\textasciicircum{}2 + tau\_xy\textasciicircum{}2) que da la tesis (en p. 15, ec. (15), da tau\_max = (sigma\_x - sigma\_y)/2 ya en ejes principales). (b) von Mises: la 2.a ed. NO usa el nombre \textquotedbl{}von Mises\textquotedbl{} ni escribe una...}
\end{itemize}
{\small\campo{Sin respaldo localizable:} 1) La expresion explicita de la tension equivalente de von Mises, sigma\_VM = sqrt(1/2[(s1-s2)\textasciicircum{}2+(s2-sz)\textasciicircum{}2+(sz-s1)\textasciicircum{}2)]) (02\_marco\_teorico.tex:298-312): NO aparece en esta fuente. La 2.a ed. no nombra a von Mises en el cuerpo del texto (solo en una nota al pie de la p. 149) ni define ninguna tension equivalente escalar; lo mas cercano es la energia de distorsion, ec. (88), p. 149, y la tension cortante octaedrica, art. 71, p. 219 (\textquotedbl{}This shear stress is called the 'octahedral shear stress' [...] It occurs frequently in the theory of plasticity\textquotedbl{}). Tampoco respalda la afirmacion de que sea \textquotedbl{}la medida escalar con que habitualmente se resume un estado tensional biaxial\textquotedbl{}: eso es uso moderno de MEF, propio de cook2002concepts. 2) La forma cerrada de sigma\_\{1,2\} con tau\_xy (misma ecuacion): no figura como tal; Timoshenko llega a las direcciones principales por tan 2alpha = 2 tau\_xy/(sigma\_x - sigma\_y), ec. (14), p. 14, y a las tensiones por construccion grafica del circulo de Mohr, p. 14-15. 3) Las cifras de error del software (0,04 \%, 0,26 \%, 2,9 \%, etc.) son resultados propios de la tesis y no se le atribuyen a la fuente; no requieren respaldo aqui. 4) ADVERTENCIA DE EDICION, critica para los localizadores: el ejemplar disponible es la 2.a edicion (1951) y todas las paginas impresas de este informe son de ese ejemplar, mientras que referencias.bib declara edition = \{3\}, year = \{1970\}. La paginacion de la 3.a ed. NO coincide con la de la 2.a (la 3.a intercala el capitulo de...\par}
\subsection{\texttt{virtanen2020scipy}}
\begin{itemize}\small
\item \textbf{[NO]} El sistema se resuelve \textquotedbl{}con un ordenamiento de columnas de minimo grado que aprovecha la simetria de la matriz para reducir el llenado (fill-in)\textquotedbl{}. tesis/capitulos/02\_marco\_teorico.tex:296
\\ \campo{El articulo de SciPy no menciona en ningun lugar el ordenamiento de columnas de minimo grado ni la reduccion de fill-in. Lo mas cercano es la mencion de SuperLU en la p. 264 ('SuperLU63,64 for solving sparse linear systems') y en la p. 265 ('SuperLU63 was updated to version...'), pero el detalle del ordenamiento COLAMD pertenece a la documentacion de SuperLU (referencias 63 y 64 del propio articulo: Li XS et al., SuperLU Users' Guide, Report LBNL-44289),...}
\item \textbf{[PARCIAL]} NumPy y SciPy \textquotedbl{}proveen el algebra densa y dispersa, la cuadratura y el solucionador lineal del metodo\textquotedbl{}. tesis/capitulos/03\_diseno\_implementacion.tex:25
\\ \campo{La p. 263 (Box 2) respalda plenamente el algebra DENSA ('basic solver for Ax = b', descomposiciones LU/Cholesky/QR) y, en la entrada 'sparse' de la misma pagina, el algebra DISPERSA y el solucionador lineal. Lo que NO queda respaldado en el sentido que la tesis necesita es 'la cuadratura': el articulo solo habla de integracion en terminos genericos --- p. 261, resumen: 'SciPy includes algorithms for optimization, integration, interpolation, eigenvalue...}
\end{itemize}
{\small\campo{Sin respaldo localizable:} Una sola afirmacion queda sin respaldo localizable en esta fuente:

1) tesis/capitulos/02\_marco\_teorico.tex:296 --- \textquotedbl{}con un ordenamiento de columnas de minimo grado que aprovecha la simetria de la matriz para reducir el llenado (fill-in)\textquotedbl{}. Busque \textquotedbl{}minimum degree\textquotedbl{}, \textquotedbl{}COLAMD\textquotedbl{}, \textquotedbl{}fill-in\textquotedbl{}, \textquotedbl{}permutation\textquotedbl{} y \textquotedbl{}reorder\textquotedbl{} sobre las 15 paginas del PDF: cero ocurrencias. El articulo solo dice que SciPy envuelve SuperLU para resolver sistemas lineales dispersos (p. 264) y que SuperLU fue actualizado de version (p. 265), sin entrar en la estrategia de ordenamiento. Ese detalle pertenece a las referencias 63/64 del propio articulo (Li XS et al., SuperLU Users' Guide, Report LBNL-44289, Lawrence Berkeley National Laboratory) o a la documentacion de scipy.sparse.linalg.splu. Recomendacion: cambiar la fuente de esa frase o eliminarla del alcance de \textbackslash{}autocite\{virtanen2020scipy\}.

Y una afirmacion parcialmente sin respaldo:

2) tesis/capitulos/03\_diseno\_implementacion.tex:25 --- el termino \textquotedbl{}cuadratura\textquotedbl{} atribuido a NumPy/SciPy. El articulo respalda \textquotedbl{}integration\textquotedbl{} en sentido generico (integrales definidas y EDOs, p. 261 y p. 263), no la cuadratura de Gauss-Legendre del elemento isoparametrico, que EduFEM implementa en fem/ sin usar scipy.integrate. La unica aparicion literal de \textquotedbl{}quadrature\textquotedbl{} (p. 262) es historica y se refiere a paquetes anteriores a SciPy.

Observaciones generales:
- Ninguna de las 5 citas a esta clave lleva localizador (\textbackslash{}autocite[p.\textasciitilde{}NN]). Si el tribunal exige localizador en citas de...\par}
\subsection{\texttt{zienkiewicz2013fem}}
\begin{itemize}\small
\item \textbf{[PARCIAL]} ADVERTENCIA GLOBAL SOBRE LAS PAGINAS DE ESTE INFORME (afecta a todas las filas). El .bib declara edicion 7 (2013), ISBN 978-1-85617-633-0; el PDF verificado es la 6.a edicion (portada: 'Sixth edition'; creditos p. PDF 5: 'Sixth edition 2005'; ISBN 0 7506 6320 0; Elsevier Butterworth-Heinemann, con la cooperacion de CIMNE)....
\\ \campo{Nota metodologica, no una afirmacion de la tesis. La incluyo primero porque condiciona la lectura de todas las paginas siguientes.}
\item \textbf{[PARCIAL]} \textquotedbl{}El integrando no es polinomico: como B depende de J\textasciicircum{}\{-1\}, que es racional en (xi,eta), la integral carece de primitiva elemental\textquotedbl{} y por eso se aproxima con cuadratura de Gauss-Legendre (02\_marco\_teorico.tex:235; cita conjunta con hughes2000fem).
\\ \campo{La conclusion (hay que integrar numericamente porque la integracion algebraica no es viable) esta respaldada literalmente en la p. 147, y la p. 146 introduce explicitamente J\textasciicircum{}\{-1\} en B. Lo que NO dice Zienkiewicz con esas palabras es el argumento intermedio de la tesis ('J\textasciicircum{}\{-1\} es racional en (xi,eta)' y 'carece de primitiva elemental'): el libro solo afirma que la forma explicita de Gbar 'defies our mathematical skill'. Si se quiere el argumento de...}
\item \textbf{[PARCIAL]} \textquotedbl{}En el elemento Q4 los puntos de la regla 2x2 coinciden con los puntos de Barlow, ubicaciones donde las tensiones convergen con un orden superior al del resto del elemento (superconvergencia)\textquotedbl{} (02\_marco\_teorico.tex:315; repetido en 04\_resultados.tex:58 como explicacion del orden O(h\textasciicircum{}\{1,54\}) del Q4). Cita conjunta con...
\\ \campo{ATENCION - discrepancia de fondo. El concepto general (superconvergencia en los puntos de Gauss, atribuido a Barlow) SI esta respaldado en pp. 461-462. Pero la identificacion concreta que hace la tesis para el Q4 (p = 1) NO lo esta: segun la regla de la p. 464 ('the Gauss-Legendre points numbering p') a p = 1 le corresponde UN punto, el centro del elemento, no los cuatro de la regla 2x2. Lo confirman las figuras: la Fig. 13.6 (p. 465, leida como imagen)...}
\item \textbf{[PARCIAL]} El promediado nodal \textquotedbl{}suaviza el campo y produce un contorno continuo\textquotedbl{} y \textquotedbl{}el salto entre los valores sin promediar es un indicador util del error de discretizacion -es el principio en que se apoyan los estimadores de error por recuperacion-\textquotedbl{} (02\_marco\_teorico.tex:357).
\\ \campo{El respaldo es solido para: (i) el campo crudo es discontinuo entre elementos y el promediado nodal lo suaviza (p. 207), y (ii) los estimadores por recuperacion miden la diferencia entre el campo recuperado y el campo de elementos finitos (pp. 476-477, Eq. 13.39-13.40). El matiz: Zienkiewicz define el estimador como \textbar{}\textbar{}sigma* - sigma\_hat\textbar{}\textbar{} (recuperado menos crudo), no literalmente como 'el salto entre los valores sin promediar' de elementos vecinos; el...}
\end{itemize}
{\small\campo{Sin respaldo localizable:} 1) IDENTIFICACION Q4 \textless{}-\textgreater{} PUNTOS DE BARLOW EN LA REGLA 2x2 (02\_marco\_teorico.tex:315 y 04\_resultados.tex:58). No encontre respaldo en esta fuente y ademas la contradice: en la p. 464 la regla es \textquotedbl{}the Gauss-Legendre points numbering p\textquotedbl{}, y la Fig. 13.6 (p. 465) y la Fig. 13.8 (p. 467, recuadro \textquotedbl{}4-node elements\textquotedbl{}) dibujan UN solo punto de muestreo superconvergente, en el centro del cuadrilatero bilineal. Zienkiewicz solo llega a los \textquotedbl{}dos puntos de Gauss\textquotedbl{} en el caso CUADRATICO unidimensional (p. 461, Fig. 13.4). Recomendacion: dejar cook2002concepts (o el articulo original de Barlow 1976, Int. J. Numer. Meth. Eng. 10:243-251) como unico soporte de esa frase, o reformularla. Esto tiene efecto colateral sobre 04\_resultados.tex:58, donde el orden O(h\textasciicircum{}\{1,54\}) del Q4 se explica \textquotedbl{}por la superconvergencia del muestreo en los puntos de Barlow\textquotedbl{}: la explicacion sigue siendo defendible, pero no con esta fuente.

2) \textquotedbl{}el campo recuperado converge de todos modos a O(h\textasciicircum{}2)\textquotedbl{} para el Q9 (02\_marco\_teorico.tex:315) y las cifras O(h\textasciicircum{}\{2,00\}) / O(h\textasciicircum{}\{1,54\}) de 04\_resultados.tex: son resultados propios de la tesis, no atribuidos a esta fuente. Solo la tasa TEORICA generica (O(h\textasciicircum{}\{p+1\}) / O(h\textasciicircum{}p)) tiene respaldo aqui (p. 38).

3) \textquotedbl{}carece de primitiva elemental\textquotedbl{} porque \textquotedbl{}J\textasciicircum{}\{-1\} es racional en (xi,eta)\textquotedbl{} (02\_marco\_teorico.tex:235): el razonamiento explicito no aparece; el libro (p. 147) solo dice que la forma explicita de Gbar no es simple y que \textquotedbl{}algebraic integration usually defies our mathematical skill\textquotedbl{}.

4)...\par}
\section{Catálogo de respaldos}
\noindent\small Para cada fuente: el ejemplar consultado y, por cada afirmación de la tesis, la página impresa y el pasaje que la sostiene.\par\vspace{5pt}
\begin{pasaje}\footnotesize\textbf{Nota de vigencia.} Los comentarios de este catálogo se redactaron durante la auditoría y describen el estado de la tesis \emph{en ese momento}. Varias de las observaciones que mencionan ya fueron aplicadas el 9 de septiembre de 2026: las nueve entradas que declaraban una edición distinta de la del ejemplar están corregidas, y las citas a Roache 1998, Pérez-Santiago 2023 y Cook 1974 —fuentes de las que no se dispone— fueron sustituidas. Donde un comentario diga «el .bib declara» o «DISCREPA», se refiere al estado anterior; la sección 1 de este documento refleja el estado vigente.\end{pasaje}
\subsection{\texttt{alvarez\_metodologia} --- Metodologia de la Investigacion Cientifica}
{\footnotesize\campo{Ejemplar:} Dr. Carlos Alvarez de Zayas. \campo{Archivo:} \texttt{Alvarez de Zayas (sin fecha) - Metodologia de la Investigacion Cientifica.pdf}\par}\vspace{4pt}
\vspace{3pt}\noindent{\small\textbf{p.\ 7} \quad \campo{La insuficiencia de una herramienta con esos atributos es la CONTRADICCION que origina el trabajo. (tesis/capitulos/01\_introduccion.tex:15, cita sin localizador)}}\par
\begin{pasaje}\footnotesize El problema, (el por que), de la investigacion, lo podemos definir como la situacion propia de un objeto, que provoca una necesidad en un sujeto, el cual desarrollara una actividad para transformar la situacion mencionada y resolver el problema. [...] El problema se manifiesta externamente en el objeto y es consecuencia precisamente del desconocimiento de elementos y relaciones que existen en el mismo. El planteamiento del problema cientifico es la expresion de los limites del conocimiento cientifico actual que genera la insatisfaccion en el sujeto: la necesidad.\end{pasaje}
\vspace{3pt}\noindent{\small\textbf{p.\ 8} \quad \campo{Uso de la categoria OBJETO DE ESTUDIO como parte de la realidad sobre la que actua el investigador para resolver el problema. (tesis/capitulos/01\_introduccion.tex:17, cita sin localizador)}}\par
\begin{pasaje}\footnotesize El objeto de la Investigacion Cientifica (el que?) es aquella parte de la realidad objetiva, sobre la cual actua el investigador en el proceso de la Investigacion Cientifica con vista a la solucion del problema y que es construido idealmente por este, como sujeto activo de dicho proceso, sobre bases teoricas cientificamente fundamentadas y que se encuentra condicionado por el escenario historico cultural donde se desarrolle el mismo. El objeto de la investigacion es aquella parte de la realidad que se abstrae como consecuencia de agrupar, en forma sistemica, un conjunto de fenomenos, hechos o procesos, que el investigador presupone afines, a partir del problema. Es decir, el problema es la manifestacion externa del objeto en cuestion, lo que implica que cuando se va precisando el problema se hace a la vez la determinacion del objeto.\end{pasaje}
\vspace{3pt}\noindent{\small\textbf{p.\ 13} \quad \campo{El CAMPO DE ACCION es 'la parte de ese objeto que el trabajo se propone mejorar'. (tesis/capitulos/01\_introduccion.tex:17, cita sin localizador)}}\par
\begin{pasaje}\footnotesize La precision del objeto. Campo de Accion. El campo de accion o materia de estudio es aquella parte del objeto conformado por el conjunto de aspectos, propiedades y relaciones que se abstraen del objeto, en la actividad practica del sujeto, con un objetivo determinado, a partir de ciertas condiciones y situaciones. Observese que en el proceso de investigacion, del problema al objeto, hay un primer paso de abstraccion; y del objeto al campo, un segundo. El campo de accion es un concepto mas estrecho que el de objeto, es una parte del mismo, por ejemplo, el objeto puede ser el proceso de ensenanza-aprendizaje de la asignatura y el campo de accion son los contenidos de dicha materia.\end{pasaje}
\vspace{3pt}\noindent{\small\textbf{p.\ 21} \quad \campo{El modelo de investigacion es 'una representacion ideal del objeto en la que el investigador abstrae y sistematiza las relaciones que considera esenciales'. (tesis/capitulos/02b\_diseno\_metodologico.tex:14, cita sin localizador)}}\par
\begin{pasaje}\footnotesize El modelo es una representacion ideal del objeto a investigar, donde el sujeto (el investigador) abstrae todos aquellos elementos y relaciones que el considera esenciales y lo sistematiza, en el objeto modelado. El modelo teorico contiene las relaciones esenciales del objeto y del campo de accion, las que se estructuran para cumplimentar el objetivo que se propone alcanzar, es decir; el modelo trata de reflejar la realidad, pero de acuerdo con la intencion del investigador. El modelo sustituye al objeto en determinadas etapas de la investigacion.\end{pasaje}
\vspace{3pt}\noindent{\small\textbf{p.\ 22} \quad \campo{Respaldo COLATERAL, no citado explicitamente ahi: 'La variable causa es la forma en que la herramienta expone ese canal [...]; la variable efecto es la observabilidad y contrastabilidad del procedimiento intermedio' (tesis/capitulos/01\_introduccion.tex:15). El \textbackslash{}autocite cierra la oracion siguiente, pero el par causa/efecto...}}\par
\begin{pasaje}\footnotesize La variable es en principio, un concepto que determina una propiedad o cualidad del objeto que puede variar de una o mas maneras en el tiempo y que sintetiza conceptualmente lo que se quiere conocer acerca del objeto de investigacion. [...] las variables se pueden clasificar en dependiente e independiente; la dependiente se corresponde con aquella propiedad (efecto) que varia como resultado de la influencia o relacion de otra propiedad (causa) sobre ella. La propiedad o variable causa es la denominada variable independiente.\end{pasaje}
\vspace{6pt}
\subsection{\texttt{atkinson2000learning} --- Learning from Examples: Instructional Principles from the Worked...}
{\footnotesize\campo{Ejemplar:} Robert K. Atkinson, Sharon J. Derry, Alexander Renkl, Donald Wortham. American Educational Research Association (Review of Educational Research 70(2):181-214). 2000. \campo{Archivo:} \texttt{Atkinson et al 2000 - Learning from Examples.pdf}\par}\vspace{4pt}
\vspace{3pt}\noindent{\small\textbf{p.\ 181} \quad \campo{02\_marco\_teorico §1.2, tercer principio, y 04\_resultados matriz de principios: un ejemplo resuelto es la solucion de un experto que el alumno estudia; los ejemplos eficaces tienen componentes integrados, estructura rotulada o segmentada y varios ejemplos por tipo de problema. \textbackslash{}autocite[p.\textasciitilde{}181]\{atkinson2000learning\}}}\par
\begin{pasaje}\footnotesize Worked examples are instructional devices that provide an expert's problem solution for a learner to study. Worked-examples research is a cognitive-experimental program that has relevance to classroom in- struction and the broader educational research community. A frame- work for organizing the findings of this research is proposed, leading to instructional design principles. For instance, one instructional de- sign\end{pasaje}
\vspace{3pt}\noindent{\small\textbf{p.\ 187} \quad \campo{02\_marco\_teorico §1.2, tercer principio: separar el texto de la figura o de las ecuaciones divide la atencion y perjudica el aprendizaje. \textbackslash{}autocite[p.\textasciitilde{}187]\{atkinson2000learning\}}}\par
\begin{pasaje}\footnotesize the split-attention effect and hypothesized that it interfered with the student's acquisition of schemas representing the basic domain concepts and principles that students should learn from examples. Tarmizi and Sweller (1988) further hypothesized that requiring students to split their attention between multiple sources of information during example- based geometry instruction would decrease subsequent problem-solvi\end{pasaje}
\vspace{6pt}
\subsection{\texttt{bathe2014fem} --- Finite Element Procedures}
{\footnotesize\campo{Ejemplar:} Klaus-Jürgen Bathe (Professor of Mechanical Engineering, Massachusetts Institute of Technology). Prentice Hall (Prentice-Hall, Inc. / Simon \&amp; Schuster, A Viacom Company). 1996. \campo{Archivo:} \texttt{Bathe 1996 - Finite Element Procedures.pdf}\par}\vspace{4pt}
\vspace{3pt}\noindent{\small\textbf{p.\ 345 (y 339)} \quad \campo{«EduFEM emplea elementos isoparametricos cuadrilateros, en los que la misma familia de funciones de forma interpola tanto la geometria como el campo de desplazamientos» --- tesis/capitulos/02\_marco\_teorico.tex:101 (\textbackslash{}autocite\{bathe2014fem,onate2009structural\}, sin localizador)}}\par
\begin{pasaje}\footnotesize p. 345: «In the isoparametric formulation the element displacements are interpolated in the same way as the geometry; i.e., we use [ (5.19): u = sum h\_i u\_i ; v = sum h\_i v\_i ; w = sum h\_i w\_i ] where u, v, and w are the local element displacements at any point of the element and u\_i, v\_i, and w\_i, i = 1, ..., q, are the corresponding element displacements at its nodes.» \textbar{}\textbar{} p. 339: «The interpolation of the element coordinates and element displacements using the same interpolation functions, which are defined in a natural coordinate system, is the basis of the isoparametric finite element formulation.»\end{pasaje}
\vspace{3pt}\noindent{\small\textbf{p.\ 346 y 347} \quad \campo{«La condicion det J \textgreater{} 0 en todo el elemento garantiza que el mapeo sea biyectivo y el elemento valido [...] Un determinante nulo o negativo indica un elemento degenerado o invertido» --- tesis/capitulos/02\_marco\_teorico.tex:176 (\textbackslash{}autocite\{bathe2014fem,cook2002concepts\}, sin localizador)}}\par
\begin{pasaje}\footnotesize p. 346: «[ (5.25): d/dx = J\textasciicircum{}-1 d/dr ] which requires that the inverse of J exists. This inverse exists provided that there is a one-to-one (i.e., unique) correspondence between the natural and the local coordinates of the element, as expressed in (5.20) and (5.21). [...] However, in cases where the element is much distorted or folds back upon itself, as in Fig. 5.6, the unique relation between the coordinate systems does not exist (see also Section 5.3.2 for singularities in the Jacobian transformation, Example 5.17).» \textbar{}\textbar{} p. 347, rotulo de la Fig. 5.6: «For J nonsingular, all interior angles must be smaller than 180 degrees.» y «Figure 5.6 Elements with possible singular Jacobian»; en el cuerpo: «dV = det J dr ds dt (5.28) where det J is the determinant of the Jacobian operator in (5.24)».\end{pasaje}
\vspace{3pt}\noindent{\small\textbf{p.\ 469, 470 (Tabla 5.9) y 472} \quad \campo{«Una subintegracion puede dejar la matriz de rigidez con rango deficiente y admitir modos de deformacion de energia nula [...] Una sobreintegracion (por ejemplo 3x3 en Q4) no aporta precision adicional y solo encarece el calculo. La eleccion 2x2/3x3 equilibra exactitud y estabilidad» --- tesis/capitulos/02\_marco\_teorico.tex:244...}}\par
\begin{pasaje}\footnotesize p. 469: «We recommend that full numerical integration always be used for a displacement-based or mixed finite element formulation, where we define \textquotedbl{}full\textquotedbl{} numerical integration as the order that gives the exact matrices (i.e., the analytically integrated values) when the elements are geometrically undistorted. Table 5.9 lists this order for elements used in two-dimensional analyses.» \textbar{}\textbar{} p. 470, TABLE 5.9 «Recommended full Gauss numerical integration orders for the evaluation of isoparametric displacement-based element matrices»: «4-node ... 2 x 2 \textbar{} 4-node distorted ... 2 x 2 \textbar{} 8-node ... 3 x 3 \textbar{} 9-node ... 3 x 3 \textbar{} 9-node distorted ... 3 x 3 \textbar{} 16-node ... 4 x 4». \textbar{}\textbar{} p. 472: «This integration order yields an element stiffness matrix with one spurious zero energy mode (see Exercise 5.56); that is, the element matrix not only has three zero eigenvalues (corresponding to the physical rigid...\end{pasaje}
\vspace{3pt}\noindent{\small\textbf{p.\ 282 y 283 (y 277)} \quad \campo{«El bloqueo volumetrico aparece en deformacion plana cuando el material se aproxima a la incompresibilidad (nu -\textgreater{} 0,5) [...] El resultado es [...] una respuesta artificialmente rigida, acompanada a menudo de oscilaciones espurias en el campo de tensiones» --- tesis/capitulos/02\_marco\_teorico.tex:253...}}\par
\begin{pasaje}\footnotesize p. 282: «when nu = 0.499, the same meshes of nine-node displacement-based elements result into poor stress predictions [see Fig. 4.20(c)]. Large stress fluctuations are seen in the individual elements of the coarse mesh and the fine mesh. Hence, in summary, we see here that the displacement-based element used in the analysis is effective when nu = 0.3, but as nu approaches 0.5, the stress prediction becomes very inaccurate.» \textbar{}\textbar{} p. 283: «We refer to finite element formulations with this desirable behavior as nonlocking, whereas otherwise the finite elements are locking. The term \textquotedbl{}locking\textquotedbl{} is based upon experiences in the analysis of beams, plates, and shells [...] where an inappropriate formulation -- one that locks -- results in displacements very much smaller than those intuitively expected for a given mesh [...] In the analysis of almost incompressible behavior, using a formulation...\end{pasaje}
\vspace{3pt}\noindent{\small\textbf{p.\ 472 y 473} \quad \campo{«Los modos espurios (hourglass) [...] surgen al subintegrar la rigidez [...] la matriz elemental queda con rango deficiente y admite modos de deformacion de energia nula, patrones no fisicos que contaminan la solucion. EduFEM los evita empleando la integracion completa de cada elemento (2x2 para Q4 y 3x3 para Q9)» ---...}}\par
\begin{pasaje}\footnotesize p. 472: «Figure 5.40 shows a very simple analysis case using a single eight-node element with 2 x 2 Gauss integration in which the model is unstable; that is, if the solution is obtained, the calculated nodal point displacements are very large and have no resemblance to the correct solution. [...] However, in a large, complex model, in general, elements with spurious zero energy modes in an uncontrolled manner improve the overall solution results, introduce large errors, or result in an unstable solution.» y «Such modes of no physical reality -- which we refer to also as \textquotedbl{}phantom\textquotedbl{} modes -- can introduce uncontrolled errors into a dynamic step-by-step solution». \textbar{}\textbar{} p. 473: «For these reasons any element with a spurious zero energy mode should not be used in engineering practice, in linear or in nonlinear analysis, and we therefore do not discuss such elements in this book.»\end{pasaje}
\vspace{3pt}\noindent{\small\textbf{p.\ 476 (y 477)} \quad \campo{«Existen tecnicas correctivas [...] entre ellas la integracion reducida selectiva (SRI), que subintegra solo los terminos responsables del bloqueo, y la formulacion B-bar, que proyecta la parte volumetrica o de corte de la deformacion» --- tesis/capitulos/02\_marco\_teorico.tex:257 y, con la misma redaccion,...}}\par
\begin{pasaje}\footnotesize p. 476, seccion 5.5.6 «Reduced and Selective Integration»: «When such a reduction in the order of the numerical integration is used, we refer to the procedure as reduced integration. For example, the use of 2 x 2 Gauss integration (although not recommended for use in practice; see Section 5.5.5) for the nine-node isoparametric element stiffness matrix corresponds to a reduced integration. In addition to merely using a reduced integration order, selective integration may also be considered, in which case different strain terms are integrated with different orders of integration. In these cases of reduced and selective numerical integration the specific integration scheme should be regarded as an integral part of the element formulation.» \textbar{}\textbar{} p. 477, Ejemplo 5.44: «This integration evaluates the stiffness matrix terms corresponding to bending exactly, whereas the terms corresponding to the...\end{pasaje}
\vspace{3pt}\noindent{\small\textbf{p.\ 726 y 727} \quad \campo{«Un modelo bien restringido -- al menos tres GDL que impidan los movimientos de cuerpo rigido -- hace que K\_ff sea definida positiva y, por tanto, no singular» --- tesis/capitulos/02\_marco\_teorico.tex:291 (\textbackslash{}autocite\{bathe2014fem\}, cita unica, sin localizador)}}\par
\begin{pasaje}\footnotesize p. 726, seccion 8.2.5: «In the Gauss elimination discussed so far, we implicitly assumed that the stiffness matrix K is positive definite, that is, that the structure considered is properly restrained and stable. [...] Note that the stiffness matrix of a finite element is not positive definite unless the element has been properly restrained, i.e., the rigid body motions have been suppressed. Instead, the stiffness matrix of an unrestrained finite element is positive semidefinite, U\textasciicircum{}T K U \textgreater{}= 0 (8.32), where U\textasciicircum{}T K U = 0 when U corresponds to a rigid body mode. [...] The stiffness matrix of the structure is then rendered positive definite by eliminating the rows and columns that correspond to the restrained degrees of freedom, i.e., by eliminating the possibility for the structure to undergo rigid body motions.» \textbar{}\textbar{} p. 727, EXAMPLE 8.11: «The plane stress element has three rigid body modes:...\end{pasaje}
\vspace{3pt}\noindent{\small\textbf{p.\ 348 (y 345--347)} \quad \campo{«Ademas, la formulacion isoparametrica, la cuadratura de Gauss y el ensamblaje son directamente generalizables a tres dimensiones» --- tesis/capitulos/01\_introduccion.tex:64 (\textbackslash{}autocite\{bathe2014fem,zienkiewicz2013fem\})}}\par
\begin{pasaje}\footnotesize p. 348: «These matrices are evaluated using numerical integration, as indicated for the stiffness matrix K in (5.30). [...] This formulation was for one-, two-, or three-dimensional elements. We shall now consider some specific cases and demonstrate the details of the calculation of element matrices.»\end{pasaje}
\vspace{3pt}\noindent{\small\textbf{p.\ 476 (parcial) --- no localizado para el...} \quad \campo{«El elemento Q4 exhibe el fenomeno de bloqueo por cortante con claridad [...] comportamiento caracteristico de un elemento que resulta excesivamente rigido cuando debe representar flexion mediante una interpolacion bilineal incapaz de reproducir el modo de curvatura pura» --- tesis/capitulos/04\_resultados.tex:172...}}\par
\begin{pasaje}\footnotesize p. 476: «we recall that the displacement formulation of finite element analysis yields a strain energy smaller than the exact strain energy of the mathematical/mechanical model being considered, and physically, a displacement formulation results in overestimating the system stiffness.» \textbar{}\textbar{} p. 275 (sobre el elemento de VIGA de dos nodos, no sobre el Q4 plano): «the element predictive capability of the pure displacement-based formulation is drastically different, displaying a behavior that is much too stiff when the element is thin [...] the pure displacement-based element shows an erroneous shear contribution.»\end{pasaje}
\vspace{3pt}\noindent{\small\textbf{p.\ 345--347} \quad \campo{«la cadena cinematica N -\textgreater{} J -\textgreater{} B precede a la constitutiva, y la deformacion eps = B u antecede a la tension sigma = D eps, en consonancia con la presentacion clasica de la literatura» --- tesis/capitulos/04\_resultados.tex:209 (\textbackslash{}autocite\{bathe2014fem,cook2002concepts\})}}\par
\begin{pasaje}\footnotesize p. 345: «To be able to evaluate the stiffness matrix of an element, we need to calculate the strain-displacement transformation matrix.» \textbar{}\textbar{} p. 346: «Using (5.19) and (5.25), we evaluate du/dx, du/dy, du/dz, dv/dx, ..., dw/dz and can therefore construct the strain-displacement transformation matrix B, with eps = B u\textasciicircum{} (5.26)» \textbar{}\textbar{} p. 347: «The element stiffness matrix corresponding to the local element degrees of freedom is then K = integral\_V B\textasciicircum{}T C B dV (5.27)»\end{pasaje}
\vspace{6pt}
\subsection{\texttt{bishay2020teaching} --- Teaching the finite element method fundamentals to undergraduate...}
{\footnotesize\campo{Ejemplar:} Peter L. Bishay (PhD, Assistant Professor, Department of Mechanical Engineering, California State University, Northridge, California). Firma en la cabecera de pagina: \textquotedbl{}BISHAY\textquotedbl{}. En el bloque \textquotedbl{}How to cite this article\textquotedbl{} figura como \textquotedbl{}Bishay PL\textquotedbl{}.. Wiley Periodicals LLC (revista Computer Applications in Engineering Education, John Wiley \& Sons). No figura ciudad impresa.. 2020. \campo{Archivo:} \texttt{Bishay 2020 - Teaching the Finite Element Method Fundamentals.pdf}\par}\vspace{4pt}
\vspace{3pt}\noindent{\small\textbf{p.\ 1 (pie de pagina)} \quad \campo{NOTA PREVIA SOBRE PAGINACION (afecta a todos los renglones): el .bib declara pages = \{1007--1027\} (referencias.bib:148), pero el ejemplar PDF disponible imprime la paginacion 1-21. Encabezados: 'BISHAY \textbar{}3', '4\textbar{}', ... '\textbar{} 21'; pie de la p.1: 'Comput Appl Eng Educ. 2020;1-21.  wileyonlinelibrary.com/journal/cae'. Todas las paginas...}}\par
\begin{pasaje}\footnotesize Comput Appl Eng Educ. 2020;1-21.        wileyonlinelibrary.com/journal/cae   \textbar{} (c) 2020 Wiley Periodicals LLC        1\end{pasaje}
\vspace{3pt}\noindent{\small\textbf{p.\ 6 (constructor/analizador) y 5...} \quad \campo{02\_marco\_teorico.tex:11 - 'Bishay propone un constructor y un analizador de reticulados con los que el estudiante arma estructuras de barras y observa el ensamblaje y la resolucion del sistema, complementados con miniproyectos generados por los propios alumnos'}}\par
\begin{pasaje}\footnotesize [p. 6] 'To allow students perform sweep studies by running the code multiple times with varying any specific parameter to understand its effect on any given output, user-defined functions called \textquotedbl{}Analyze\_2D\_Truss\textquotedbl{} and \textquotedbl{}Analyze\_3D\_Truss\textquotedbl{} were developed. [...] Main functions, called \textquotedbl{}Main\_2D\_Truss\_Analyzer\textquotedbl{} and \textquotedbl{}Main\_3D\_Truss\_Analyzer\textquotedbl{}, as shown in Figure 2, are used to pass all inputs to these analyzer functions and obtain the outputs.' / 'To allow students create and analyze complicated trusses, another main code called \textquotedbl{}Main\_Truss\_Builder\textquotedbl{} was also provided to them. The code uses the nodes and elements of a unit cell to do any of the following: Generate duplicates at any distance in a certain direction, generate duplicates at any angle around any axis, or generate N duplicates on a circle around any axis.' --- [p. 5] 'The concept of finite element assembly to obtain the global...\end{pasaje}
\vspace{3pt}\noindent{\small\textbf{p.\ 17} \quad \campo{02\_marco\_teorico.tex:11 y :37 - las herramientas de Bishay 'no abordan el continuo elastico bidimensional' / 'se limita a estructuras de barras y reticulados'}}\par
\begin{pasaje}\footnotesize 'Multiple extensions or modifications can be integrated in the proposed approach. For example, the truss analysis code can be extended to consider 2D or 3D frames made of 2D or 3D beam elements that can carry transverse loads and bend accordingly. The Truss Builder code will still be working with this extension, but the Truss Analyzer code will change, as now there will be additional rotational degrees of freedom at each node (one for 2D analysis and three for 3D analysis).'\end{pasaje}
\vspace{3pt}\noindent{\small\textbf{p.\ 2} \quad \campo{02\_marco\_teorico.tex:9 - 'las herramientas profesionales son cajas negras optimizadas para producir resultados, no para revelar el procedimiento intermedio [...] La investigacion educativa documenta esta brecha y propone el software especificamente didactico como puente entre ambos extremos'}}\par
\begin{pasaje}\footnotesize 'As most engineering programs do not require coverage of FEM theory in the undergraduate curriculum, students deal with FEA software as a black box that provides them with the problem solution magically based on the given inputs, without any understanding of at least the important fundamentals of the theory.' [y mas abajo, misma pagina] 'Although teaching the use of FEA commercial software to perform different studies is much easier than teaching the FEM theory, introducing the fundamentals of the theory is still very important. Zamani [36] highlighted the challenges in teaching FEA to undergraduate students due to the widespread availability of commercial FEA software in the academic institutions and the pressure from industry to include it in the modern undergraduate curriculum.'\end{pasaje}
\vspace{3pt}\noindent{\small\textbf{p.\ 17 (y 2 para la parte de 'herramientas...} \quad \campo{01\_introduccion.tex:21 - 'La literatura sobre ensenanza del metodo coincide en que la simulacion interactiva y la manipulacion directa de los objetos matematicos del MEF favorecen un aprendizaje mas profundo que la mera exposicion teorica o el uso de herramientas opacas'}}\par
\begin{pasaje}\footnotesize 'Teaching FEM using the proposed approach of building a useful code with the students step by step was found to be very effective in helping students understand the process of solving an FEA problem. The concepts were strengthened through the mini-projects approach that allowed the students to generate their own assignments including the description of the problem and various types of questions, some of which promote critical thinking.'\end{pasaje}
\vspace{3pt}\noindent{\small\textbf{p.\ 13 (dato estadistico) y 17 (atribucion...} \quad \campo{03\_diseno\_implementacion.tex:9 - 'El software educativo interactivo demuestra eficacia para consolidar conceptos abstractos del analisis estructural'}}\par
\begin{pasaje}\footnotesize [p. 13] 'Testing the null hypothesis that the performance index data of the two cohorts are independent random samples from normal distributions with equal means and unknown and unequal variances, using the two-sample t-test and 5\% significance level gave h = 1 which indicates that the t-test rejects the null hypothesis. The p value is .048. This shows that those who worked on the mini-project (Cohort 2) made significant improvement in their performance compared to the students who worked on a regular HW assignment instead (Cohort 1).' --- [p. 17] 'Accordingly, the performance improvement is believed to happen mainly because of the assignment-generation mini-projects idea, which was facilitated by using the Truss Builder and Truss Analyzer computational tools.'\end{pasaje}
\vspace{3pt}\noindent{\small\textbf{p.\ 1 (resumen) y 17} \quad \campo{03\_diseno\_implementacion.tex:155 - 'El diseno de estos modulos parte de la evidencia sobre la interactividad como motor del aprendizaje en ingenieria \textbackslash{}autocite\{lee2015interactive,bishay2020teaching\}'}}\par
\begin{pasaje}\footnotesize [p. 1, Abstract] 'The anonymous students' written comments also supported the hypothesis that this new approach is more engaging to students compared to traditional approaches.' --- [p. 17] 'Being a part of generating an assignment or selecting the problem to be solved was found to be interesting to students as it allows them to think more critically and consider different scenarios. It also makes them more engaged and interested to learn all details of the topic in hand.'\end{pasaje}
\vspace{3pt}\noindent{\small\textbf{p.\ 4 y 5} \quad \campo{02\_marco\_teorico.tex:29 (tabla comparativa) - columna 'Calculo paso a paso: Media-alta' para el constructor/analizador de reticulados}}\par
\begin{pasaje}\footnotesize [p. 4] 'the author supplemented the course with a series of lectures on writing a simple computer code on MATLAB for 2D and 3D truss analysis. The code is introduced step by step as every lecture focuses on one or few concepts, as shown in Figure 1.' --- [p. 5] 'This series of lectures explained every single line in the truss analysis code.'\end{pasaje}
\vspace{3pt}\noindent{\small\textbf{p.\ 5} \quad \campo{02\_marco\_teorico.tex:29 (tabla comparativa) - columna 'Idioma: Ingles' y columna 'Elementos: Barras (reticulado)'}}\par
\begin{pasaje}\footnotesize 'ki = (Ei Ai / Li) [ G  -G ; -G  G ], where G = [c\textasciicircum{}2  cs ; cs  s\textasciicircum{}2] for 2D truss  (1)' [y] 'The last code for the full analysis of 2D trusses is extended to analyze 3D trusses where now each node has three translational degrees of freedom and the 6 x 6 version of the element stiffness matrix, ki, in Equations (1)-(3) is to be used.'\end{pasaje}
\vspace{3pt}\noindent{\small\textbf{p.\ 17} \quad \campo{05\_conclusiones.tex:51 - recomendacion de 'un estudio con estudiantes [...] mediante una comparacion antes-despues e instrumentos de percepcion'; 'los antecedentes revisados \textbackslash{}autocite\{perezsantiago2023fem,bishay2020teaching,lee2015interactive\} indican, en cambio, que conceptos conviene poner a prueba'}}\par
\begin{pasaje}\footnotesize 'The study focused on only two cohorts, 27 students each, taking a 7-week summer course; one was taught FEM traditionally and the other was taught FEM with the integration of the student-generated assignments mini-projects. Assessment was done using two methods: performance in a midterm exam taking into consideration the academic standing of each student in the prerequisite course, and student self-reporting surveys.'\end{pasaje}
\vspace{6pt}
\subsection{\texttt{chi2014icap} --- The ICAP Framework: Linking Cognitive Engagement to Active...}
{\footnotesize\campo{Ejemplar:} Michelene T. H. Chi, Ruth Wylie (Arizona State University). Routledge/Taylor \& Francis (Educational Psychologist 49(4):219-243). 2014. \campo{Archivo:} \texttt{Chi y Wylie 2014 - The ICAP Framework.pdf}\par}\vspace{4pt}
\vspace{3pt}\noindent{\small\textbf{p.\ 220} \quad \campo{02\_marco\_teorico §1.2, segundo principio, y 04\_resultados matriz de principios: el marco ICAP clasifica las conductas en pasivo, activo, constructivo e interactivo y postula que el aprendizaje crece al pasar de un modo al siguiente (I\textgreater{}C\textgreater{}A\textgreater{}P). \textbackslash{}autocite[p.\textasciitilde{}220]\{chi2014icap\}}}\par
\begin{pasaje}\footnotesize greater than the Con- structive mode, which is greater than the Active mode, which in turn is greater than the Passive mode (I\textgreater{}C\textgreater{}A\textgreater{}P). Thus, the ICAP hypothesis predicts different levels of learning for different modes of overt behaviors. Higher levels imply learning with deeper understanding. Although a preliminary version of this framework has been presented in the literature (Chi, 2009), this article expands the f\end{pasaje}
\vspace{3pt}\noindent{\small\textbf{p.\ 222} \quad \campo{02\_marco\_teorico §1.2, segundo principio: lo constructivo exige que el alumno produzca algo nuevo con la informacion, no solo que la manipule. \textbackslash{}autocite[p.\textasciitilde{}222]\{chi2014icap\}}}\par
\begin{pasaje}\footnotesize Our taxonomy defines constructive behaviors as those in which learners generate or produce additional externalized outputs or products beyond what was provided in the learning materi- als. Thus, a characteristic descriptor of the constructive mode is generative. To meet the criteria for constructive, the outputs of generative behaviors should contain new ideas that go beyond the information given; otherwise such beha\end{pasaje}
\vspace{6pt}
\subsection{\texttt{cook2002concepts} --- Concepts and Applications of Finite Element Analysis}
{\footnotesize\campo{Ejemplar:} Robert D. Cook; David S. Malkus; Michael E. Plesha; Robert J. Witt (University of Wisconsin-Madison). John Wiley \&amp; Sons, Inc.. 2002 (linea de copyright: \textquotedbl{}Copyright (c) 1974, 1981, 1989, 2002 by John Wiley \&amp; Sons Inc.\textquotedbl{}; el registro LC lleva \textquotedbl{}TA646 .C66 2001\textquotedbl{} y numero de control 2001033009). \campo{Archivo:} \texttt{Cook et al 2002 - Concepts and Applications of Finite Element Analysis.pdf}\par}\vspace{4pt}
\vspace{3pt}\noindent{\small\textbf{p.\ 218} \quad \campo{tesis/capitulos/01\_introduccion.tex:64 --- En 3D \textquotedbl{}la matriz constitutiva crece a 6x6, la de deformación-desplazamiento a 6x24 en un hexaedro de ocho nodos y la rigidez elemental a 24x24\textquotedbl{}.}}\par
\begin{pasaje}\footnotesize Equation 6.2-10 is expanded so that the square matrix is 9 by 9, and Eq. 6.2-11 is expanded so that the rectangular matrix is 9 by 24. The stiffness matrix of the eight-node element is given by

  [k]      = $\int$(-1..1)$\int$(-1..1)$\int$(-1..1) [B]\textasciicircum{}T  [E]  [B]  J d$\xi$ d$\eta$ d$\zeta$        (6.5-6)
 24$\times$24                                 24$\times$6  6$\times$6  6$\times$24

where J is the determinant of [J] in Eq. 6.5-2.\end{pasaje}
\vspace{3pt}\noindent{\small\textbf{p.\ 79} \quad \campo{tesis/capitulos/02\_marco\_teorico.tex:70 --- \textquotedbl{}La tensión plana modela cuerpos delgados cargados en su plano, en los que la tensión normal al plano es despreciable ($\sigma$z = 0)\textquotedbl{}.}}\par
\begin{pasaje}\footnotesize Constitutive matrix [E] is symmetric; E\_ij = E\_ji. [E] can represent isotropic or anisotropic material properties. For isotropy and plane stress conditions ($\sigma$z = $\tau$yz = $\tau$zx = 0), [E] and its inverse are  [E] = E/(1$\-$$\nu$²) [1 $\nu$ 0; $\nu$ 1 0; 0 0 (1$\-$$\nu$)/2] ...\end{pasaje}
\vspace{3pt}\noindent{\small\textbf{p.\ 94} \quad \campo{tesis/capitulos/02\_marco\_teorico.tex:70 --- \textquotedbl{}La deformación plana modela cuerpos prismáticos largos en los que la deformación longitudinal se anula ($\varepsilon$z = 0)\textquotedbl{}.}}\par
\begin{pasaje}\footnotesize Plane strain exists when strain normal to the analysis plane is constrained to be zero. When $\varepsilon$z = 0, matrix [E] of Eq. 3.1-3 must be replaced by

  [E] = E/((1+$\nu$)(1$\-$2$\nu$)) [1$\-$$\nu$  $\nu$  0; $\nu$  1$\-$$\nu$  0; 0  0  (1$\-$2$\nu$)/2]        (3.4-11)\end{pasaje}
\vspace{3pt}\noindent{\small\textbf{p.\ 236} \quad \campo{tesis/capitulos/02\_marco\_teorico.tex:176 --- \textquotedbl{}La condición det J \textgreater{} 0 en todo el elemento garantiza que el mapeo sea biyectivo y el elemento válido... Un determinante nulo o negativo indica un elemento degenerado o invertido; EduFEM impone un umbral mínimo para detectar y rechazar esas geometrías\textquotedbl{}.}}\par
\begin{pasaje}\footnotesize Extreme shape distortions are shown in Fig. 6.11-2. In each of these two examples, a Gauss point lies outside the element, and the Jacobian J is negative there. An acceptably shaped element should not display negative J at any Gauss point location. Indeed, in a well-shaped element J does not change much from one location to another. [...] Commercial software commonly checks the geometry of each element [6.21]. Before analysis begins, elongation and other geometric distortions can be examined and each assigned a numerical value. [...] the software might refuse to continue with analysis until shapes it considers unacceptable are corrected. The \textquotedbl{}poor shape\textquotedbl{} threshold may be high, so that software may accept elements having great enough distortion to introduce significant errors.\end{pasaje}
\vspace{3pt}\noindent{\small\textbf{p.\ 223} \quad \campo{tesis/capitulos/02\_marco\_teorico.tex:244 y 255 --- \textquotedbl{}Una subintegración puede dejar la matriz de rigidez con rango deficiente y admitir modos de deformación de energía nula: los modos espurios (hourglass), patrones de deformación no físicos que el elemento no puede resistir\textquotedbl{}.}}\par
\begin{pasaje}\footnotesize Underintegration and Spurious Modes. Use of an integration rule of less than full order is called \textquotedbl{}underintegration\textquotedbl{} or \textquotedbl{}reduced\textquotedbl{} integration. [...] But underintegration introduces the defect variously known as spurious mode, singular mode, zero-energy deformation mode, hourglass mode, kinematic mode, instability, and mechanism. An element whose stiffness matrix incorporates a spurious mode has no resistance to nodal loads that tend to activate the mode. Mathematically, an underintegrated [k] is called rank-deficient, which means that its order minus the number of possible rigid-body motions is greater than its rank.\end{pasaje}
\vspace{3pt}\noindent{\small\textbf{p.\ 223} \quad \campo{tesis/capitulos/02\_marco\_teorico.tex:244 y 255 --- \textquotedbl{}La elección 2$\times$2 / 3$\times$3 equilibra exactitud y estabilidad\textquotedbl{}; \textquotedbl{}EduFEM los evita empleando la integración completa de cada elemento (2$\times$2 para Q4 y 3$\times$3 para Q9)\textquotedbl{}.}}\par
\begin{pasaje}\footnotesize Full Integration. We define \textquotedbl{}full integration\textquotedbl{} as a quadrature rule of sufficient accuracy to exactly integrate all stiffness coefficients k\_ij of an undistorted element. [...] Then, if thickness is uniform and [E] is constant, the stiffness matrix integrand [B]\textasciicircum{}T[E][B] t J contains squares of $\xi$ and $\eta$. Therefore an order 2 Gauss rule (four points) integrates [k] exactly. [...] The [B] matrix of an undistorted eight- or nine-node plane element contains squares of $\xi$ and $\eta$, so fourth powers appear in the integrand and a 3 by 3 Gauss rule provides exact results.\end{pasaje}
\vspace{3pt}\noindent{\small\textbf{p.\ 98} \quad \campo{tesis/capitulos/02\_marco\_teorico.tex:251 y tesis/capitulos/04\_resultados.tex:172 --- \textquotedbl{}El bloqueo por cortante afecta al elemento bilineal Q4 en estados dominados por flexión... el campo de desplazamientos bilineal del Q4 no puede representar esa curvatura sin generar deformaciones de corte espurias... Esas deformaciones parásitas...}}\par
\begin{pasaje}\footnotesize Element Defects. Like the CST, the Q4 element cannot exhibit pure bending. When bent, it displays shear strain as well as the expected bending strain. This parasitic shear absorbs strain energy, so that if a given bending deformation is prescribed, the bending moment needed to produce it is larger than the correct value. In other words, the Q4 element exhibits shear locking behavior.\end{pasaje}
\vspace{3pt}\noindent{\small\textbf{p.\ 235} \quad \campo{tesis/capitulos/02\_marco\_teorico.tex:251 y 408 --- \textquotedbl{}El elemento Q9, cuya interpolación bicuadrática sí reproduce el modo de curvatura, está prácticamente libre de este bloqueo\textquotedbl{} / \textquotedbl{}el elemento Q9, de orden superior, lo evita y converge con rapidez\textquotedbl{}.}}\par
\begin{pasaje}\footnotesize TABLE 6.11-1 STRESSES AND DEFLECTIONS IN TWO-ELEMENT CANTILEVER BEAMS UNDER TRANSVERSE TIP LOAD P. EXACT VALUES BY BEAM THEORY ARE UNITY (NEGLECTING TRANSVERSE SHEAR DEFORMATION). [...] 9-node, 2$\times$2 $\rightarrow$ rectangular elements: $\sigma$xB = 1.000, vC = 1.006; trapezoidal: 1.125, 1.109; curved sides: 0.958, 0.955. [p. 236:] ... we expect the element to be capable of modeling pure bending exactly, and such is indeed the case if elements are rectangular.\end{pasaje}
\vspace{3pt}\noindent{\small\textbf{p.\ 113} \quad \campo{tesis/capitulos/02\_marco\_teorico.tex:272 --- \textquotedbl{}Para una arista lineal de Q4 con carga q constante y longitud L, el reparto exacto es F1 = (L/6)(2qs + qe) y F2 = (L/6)(qs + 2qe)\textquotedbl{}.}}\par
\begin{pasaje}\footnotesize Hence, nodal loads associated with the linear side in Fig. 3.11-1a are

  \{F4; F3\} = $\int$(-a..a) [N]\textasciicircum{}T[N] dx \{q4; q3\} + [N\_\{a/2\}]\textasciicircum{}T F = (a/3)[2 1; 1 2]\{q4; q3\} + (F/4)\{1; 3\}        (3.11-4)\end{pasaje}
\vspace{3pt}\noindent{\small\textbf{p.\ 113} \quad \campo{tesis/capitulos/02\_marco\_teorico.tex:272 --- \textquotedbl{}para la arista cuadrática de Q9 con q constante la distribución resulta L/6, 4L/6, L/6, concentrando la mayor parte en el nodo medio\textquotedbl{}.}}\par
\begin{pasaje}\footnotesize Hence, nodal loads associated with a quadratic variation on a quadratic side are

  \{F4; F7; F3\} = $\int$(-a..a) [N]\textasciicircum{}T[N] dx \{q4; q7; q3\} = (a/15)[4 2 $\-$1; 2 16 2; $\-$1 2 4]\{q4; q7; q3\}        (3.11-6)

If loading is uniform, then q4 = q7 = q3, and Eq. 3.11-6 says that of the total force on the side, one-sixth is allocated to each end node and two-thirds to the midside node.\end{pasaje}
\vspace{3pt}\noindent{\small\textbf{p.\ 117} \quad \campo{tesis/capitulos/02\_marco\_teorico.tex:298 --- \textquotedbl{}se obtienen las tensiones principales y la equivalente de von Mises, la medida escalar con que habitualmente se resume un estado tensional biaxial\textquotedbl{}.}}\par
\begin{pasaje}\footnotesize $\sigma$e = (1/$\surd$2)[($\sigma$x $\-$ $\sigma$y)² + ($\sigma$y $\-$ $\sigma$z)² + ($\sigma$z $\-$ $\sigma$x)² + 6($\tau$xy² + $\tau$yz² + $\tau$zx²)]\textasciicircum{}\{1/2\}        (3.12-2a)

$\sigma$e = (1/$\surd$2)[($\sigma$1 $\-$ $\sigma$2)² + ($\sigma$2 $\-$ $\sigma$3)² + ($\sigma$3 $\-$ $\sigma$1)²]\textasciicircum{}\{1/2\}        (3.12-2b)

[p. 116:] Some stress quantities are invariant; that is, they have the same numerical value regardless of the coordinate system in which they are computed. One such quantity is the von Mises or \textquotedbl{}effective\textquotedbl{} stress $\sigma$e, which is used to predict the onset of yielding when material behaves in a ductile fashion. [p. 117:] The von Mises stress $\sigma$e is often computed, and its element-by-element contours plotted, because it is a scalar measure of the intensity of the entire state of stress.\end{pasaje}
\vspace{3pt}\noindent{\small\textbf{p.\ 232} \quad \campo{tesis/capitulos/02\_marco\_teorico.tex:348 --- \textquotedbl{}con a $\approx$ 1,866 y c $\approx$ 0,134: cada nodo recibe el peso a del punto de Gauss más próximo, c del diagonalmente opuesto y b de los otros dos\textquotedbl{} (matriz de extrapolación E del Q4).}}\par
\begin{pasaje}\footnotesize For example, let us calculate stress $\sigma$x at corner A from the $\sigma$x values at Gauss points 1,2,3,4. We substitute r = s = $\-$$\surd$3, and obtain from Eq. 6.10-3

  $\sigma$xA = 1.866$\sigma$x1 $\-$ 0.500$\sigma$x2 + 0.134$\sigma$x3 $\-$ 0.500$\sigma$x4        (6.10-5)\end{pasaje}
\vspace{3pt}\noindent{\small\textbf{p.\ 231} \quad \campo{tesis/capitulos/02\_marco\_teorico.tex:315 --- \textquotedbl{}en el Q9 los puntos óptimos son también los de la regla 2$\times$2, y no los nueve de la regla 3$\times$3 con que se integra su rigidez\textquotedbl{}.}}\par
\begin{pasaje}\footnotesize In Q8 and Q9 elements, which are commonly integrated using four-point and nine-point rules respectively, stresses computed by applying Eq. 6.10-1 at locations of the four-point rule may have only one-tenth the error of stresses computed from the same formula at other locations [6.17]. [p. 230:] This observation does not imply that stiffness matrix integration and stress calculation must use the same order rule. Stresses may be calculated at Gauss points of a lower-order rule than used for integration.\end{pasaje}
\vspace{3pt}\noindent{\small\textbf{p.\ 231} \quad \campo{tesis/capitulos/02\_marco\_teorico.tex:315 y tesis/capitulos/04\_resultados.tex:58 --- \textquotedbl{}En el elemento Q4 los puntos de la regla 2$\times$2 coinciden con los puntos de Barlow, ubicaciones donde las tensiones convergen con un orden superior al del resto del elemento (superconvergencia)\textquotedbl{}.}}\par
\begin{pasaje}\footnotesize In brief, an analytical argument that locates optimal stress points is as follows [3.3,6.16,6.17]. Impose on an element a complete polynomial displacement field \{f\_h\} whose degree is one higher than the degree of the highest-order complete polynomial in the actual displacement field. [...] Now seek locations in the element where strains of field \{f\_h\} are the same as strains [B]\{d\_h\} [...] These locations are found to be Gauss point locations in lower-order elements, but slightly different locations in higher-order elements (which are seldom used). As an example, consider Eqs. 4.9-4. Locations at which du/dx = dū/dx are at distances (h/2)/$\surd$3 from the element center. These are Gauss point locations of an order 2 rule.\end{pasaje}
\vspace{3pt}\noindent{\small\textbf{p.\ 325} \quad \campo{tesis/capitulos/04\_resultados.tex:58 --- \textquotedbl{}Que el Q4 supere el orden O(h\textasciicircum{}1) del gradiente crudo refleja la superconvergencia del muestreo en los puntos de Barlow y del promediado nodal\textquotedbl{}.}}\par
\begin{pasaje}\footnotesize It happens that stresses $\sigma$* provided by patch recovery are superconvergent; that is, their convergence rate is at least O(h\textasciicircum{}\{p+1\}) [9.21]. For linear elements Q4 and CST in Fig. 9.9-3, the rate is O(h²); for quadratic elements Q8 and LST, the rate is O(h³) for points on the boundary and O(h?) for points internal to the mesh.\end{pasaje}
\vspace{6pt}
\subsection{\texttt{csi2017sap2000} --- CSI Analysis Reference Manual --- For SAP2000, ETABS, SAFE and...}
{\footnotesize\campo{Ejemplar:} Computers \& Structures, Inc. (autoría corporativa; la portada dice \textquotedbl{}A PRODUCT OF COMPUTERS \& STRUCTURES, INC.\textquotedbl{}, sin autores personales). Computers \& Structures, Inc.. 2017 (portadilla: \textquotedbl{}November 2017\textquotedbl{}; copyright \textquotedbl{}Computers \& Structures, Inc., 1978-2017\textquotedbl{}). \campo{Archivo:} \texttt{Computers and Structures 2017 - CSI Analysis Reference Manual.pdf}\par}\vspace{4pt}
\vspace{3pt}\noindent{\small\textbf{p.\ 179 (y 192)} \quad \campo{tesis/capitulos/04\_resultados.tex:99 --- «un modelo equivalente resuelto en SAP2000 \textbackslash{}autocite\{csi2017sap2000\} con elementos de cáscara (\textbackslash{}emph\{shell\}) de formulación Mindlin-Reissner. Esa formulación no coincide exactamente con el continuo bidimensional de tensión plana de EduFEM ---difiere en la cinemática y en la discretización---»}}\par
\begin{pasaje}\footnotesize p. 179: «Plate-bending behavior includes two-way, out-of-plane, plate rotational stiffness components and a translational stiffness component in the direction normal to the plane of the element. You may choose a thin-plate (Kirchhoff) formulation that neglects transverse shearing deformation, or a thick-plate (Mindlin/Reissner) formulation which includes the effects of transverse shearing deformation. Out-of-plane displacements are cubic.» --- p. 192: «Two thickness formulations are available, which determine whether or not transverse shearing deformations are included in the plate-bending behavior of a plate or shell element: $\cdot$ The thick-plate (Mindlin/Reissner) formulation, which includes the effects of transverse shear deformation $\cdot$ The thin-plate (Kirchhoff) formulation, which neglects transverse shearing deformation [...] It is generally recommended that you use the thick-plate...\end{pasaje}
\vspace{3pt}\noindent{\small\textbf{p.\ 216 (2D) y 178 (3D)} \quad \campo{tesis/capitulos/02\_marco\_teorico.tex:32 (fila de la tabla comparativa) --- «Software comercial (SAP2000 \textbackslash{}autocite\{csi2017sap2000\}, ANSYS, Abaqus) \& Inglés (localización parcial en algunos productos) \& 2D/3D \& Baja (caja negra) \& Comercial \& Amplios \& No (no didáctico)»}}\par
\begin{pasaje}\footnotesize p. 216 (capítulo del elemento Plane): «The Plane element is a three- or four-node element for modeling two-dimensional solids of uniform thickness. [...] Structures that can be modeled with this element include: $\cdot$ Thin, planar structures in a state of plane stress $\cdot$ Long, prismatic structures in a state of plane strain» --- p. 178 (capítulo del elemento Shell): «The Shell element is a three- or four-node formulation that combines membrane and plate-bending behavior. [...] Structures that can be modeled with this element include: [...] $\cdot$ Three-dimensional curved shells, such as tanks and domes»\end{pasaje}
\vspace{3pt}\noindent{\small\textbf{p.\ 178} \quad \campo{tesis/capitulos/02\_marco\_teorico.tex:11 --- «Por otro, los paquetes comerciales de análisis estructural \textbackslash{}autocite\{csi2017sap2000\}, que operan desde la perspectiva del estudiante como cajas negras: entregan resultados sin exponer las matrices intermedias que los produjeron.»}}\par
\begin{pasaje}\footnotesize «A four-point numerical integration formulation is used for the Shell stiffness. Stresses and internal forces and moments, in the element local coordinate system, are evaluated at the 2-by-2 Gauss integration points and extrapolated to the joints of the element. An approximate error in the element stresses or internal forces can be estimated from the difference in values calculated from different elements attached to a common joint.»\end{pasaje}
\vspace{3pt}\noindent{\small\textbf{p.\ 178} \quad \campo{tesis/capitulos/02\_marco\_teorico.tex:408 --- «comparando además contra un modelo de cáscara de SAP2000 \textbackslash{}autocite\{csi2017sap2000\}» (y, en paralelo, 05\_conclusiones.tex:19 --- «un modelo equivalente en el software comercial SAP2000 \textbackslash{}autocite\{timoshenko1970elasticity,csi2017sap2000\}: [...] frente a SAP2000, del 0,21 \% en $\sigma$x y de hasta...}}\par
\begin{pasaje}\footnotesize «The Shell element is a three- or four-node formulation that combines membrane and plate-bending behavior. The four-joint element does not have to be planar. [...] Structures that can be modeled with this element include: $\cdot$ Floor systems $\cdot$ Wall systems $\cdot$ Bridge decks $\cdot$ Three-dimensional curved shells, such as tanks and domes $\cdot$ Detailed models of beams, columns, pipes, and other structural members»\end{pasaje}
\vspace{6pt}
\subsection{\texttt{garciacordoba2005tecnologica} --- La investigacion tecnologica: Investigar, idear e innovar en...}
{\footnotesize\campo{Ejemplar:} Fernando Garcia-Cordoba (ficha catalografica: \textquotedbl{}Garcia-Cordoba, Fernando\textquotedbl{}; la cubierta exterior lo imprime sin guion: \textquotedbl{}Garcia Cordoba\textquotedbl{}). Prologuista: Angel Eduardo Vargas Garza (prol.), no coautor.. Editorial Limusa, S.A. de C.V. --- Grupo Noriega Editores. 2005. \campo{Archivo:} \texttt{Garcia-Cordoba 2005 - La investigacion tecnologica.pdf}\par}\vspace{4pt}
\vspace{3pt}\noindent{\small\textbf{p.\ 83} \quad \campo{\textquotedbl{}en la investigación tecnológica el problema no se elige libremente, sino que resulta de interpretar y diagnosticar una realidad concreta\textquotedbl{} --- c:/Tesis/Tesis Software/SoftwareED/tesis/capitulos/01\_introduccion.tex:11 (cierra el párrafo que define la 'situación problemática' del trabajo). Cita sin localizador:...}}\par
\begin{pasaje}\footnotesize 3.2.1. El problema --- \textquotedbl{}El problema en la investigación tecnológica, a diferencia del problema de la investigación científica, no se escoge libremente respondiendo a los intereses del investigador. En lo tecnológico el problema es producto de una interpretación de la realidad y requiere de un diagnóstico en el cual convergen numerosas apreciaciones relativas a diversos aspectos, entre los que se incluyen eventos, recursos y participantes.\textquotedbl{} (segundo párrafo del apartado 3.2.1). En la misma página, primer párrafo: \textquotedbl{}Un problema de investigación tecnológico posee un carácter práctico y concreto, señala un obstáculo o una necesidad e intenta modificar o crear una situación\textquotedbl{}.\end{pasaje}
\vspace{3pt}\noindent{\small\textbf{p.\ 85} \quad \campo{\textquotedbl{}en la investigación tecnológica la hipótesis es la solución tentativa a un problema concreto y su criterio de veracidad es la efectividad en la práctica, no la contrastación de un enunciado teórico\textquotedbl{} --- c:/Tesis/Tesis Software/SoftwareED/tesis/capitulos/01\_introduccion.tex:44 (justifica que el trabajo plantee una hipótesis de...}}\par
\begin{pasaje}\footnotesize 3.2.3. La hipótesis --- \textquotedbl{}La hipótesis en la investigación tecnológica es la solución tentativa a un problema concreto. Es información que determina las acciones a seguir para modificar la realidad en el sentido deseado. No es una afirmación teórica tentativa que debe de ser sometida a verificación, sino la enumeración y descripción de acciones y recursos que se prueban y modifican durante el proceso de investigación para determinar su utilidad y obtener el conocimiento operativo acertado.\textquotedbl{} Y en el párrafo siguiente: \textquotedbl{}En lo tecnológico la hipótesis posee una lógica distinta a la de la investigación científica, ya que no es el conocimiento teórico de una realidad lo que se busca, lo importante es el saber operativo. El criterio de veracidad es su efectividad en la práctica concreta y se parte de una hipótesis conformada, según el caso, por recursos, acciones y participantes, que se aprueban...\end{pasaje}
\vspace{3pt}\noindent{\small\textbf{p.\ 90-92 (apartado 3.4.1 'De acuerdo con...} \quad \campo{\textquotedbl{}por su finalidad, el trabajo es de carácter propositivo ---la propuesta de un artefacto que resuelve un problema formativo---; por su naturaleza se inscribe en la investigación aplicada de tipo tecnológico, con enfoque cuantitativo en su fase de verificación y validación\textquotedbl{} --- c:/Tesis/Tesis...}}\par
\begin{pasaje}\footnotesize p. 90: \textquotedbl{}Respecto al uso de la información producto de la investigación, se tienen cuatro tipos: pura, aplicada, tecnológica y de desarrollo. Esta clasificación, más que conformar subdivisiones de la investigación tecnológica, permite distinguirla de entre otras modalidades (fig. 3.5). [...] Investigación aplicada: es la que a partir de descubrimientos y conocimientos obtenidos en la investigación pura los emplea para alcanzar resultados que puedan ser explotables como productos, procedimientos, maquinaria, herramientas, etcétera.\textquotedbl{} --- p. 91: \textquotedbl{}Investigación tecnológica: es la búsqueda de conocimientos de carácter operativo cuyo fin es beneficiar a los sectores de extracción, transformación o servicios. Proporciona soluciones o productos que generan recursos económicos en circunstancias concretas. Son quehaceres de tipo propositivo en los cuales no sólo se establece cuál es la solución,...\end{pasaje}
\vspace{6pt}
\subsection{\texttt{harris2020numpy} --- Array programming with NumPy}
{\footnotesize\campo{Ejemplar:} Charles R. Harris; K. Jarrod Millman; Stefan J. van der Walt; Ralf Gommers; Pauli Virtanen; David Cournapeau; Eric Wieser; Julian Taylor; Sebastian Berg; Nathaniel J. Smith; Robert Kern; Matti Picus; Stephan Hoyer; Marten H. van Kerkwijk; Matthew Brett; Allan Haldane; Jaime Fernandez del Rio; Mark Wiebe; Pearu Peterson; Pierre Gerard-Marchant; Kevin Sheppard; Tyler Reddy; Warren Weckesser; Hameer Abbasi; Christoph Gohlke; Travis E. Oliphant (26 autores, mismo orden que el .bib). Springer Nature (revista Nature; \textquotedbl{}Publisher's note Springer Nature...\textquotedbl{} y (c) The Author(s) 2020, acceso abierto CC BY 4.0). 2020. \campo{Archivo:} \texttt{Harris et al 2020 - Array programming with NumPy.pdf}\par}\vspace{4pt}
\vspace{3pt}\noindent{\small\textbf{p.\ 357 (PDF pág. 1) y 361 (PDF pág. 5)} \quad \campo{«Un motor desarrollado con bibliotecas numéricas de código abierto \textbackslash{}autocite\{harris2020numpy,virtanen2020scipy\} permite exhibir las matrices y operaciones intermedias sin las restricciones de un producto cerrado» --- tesis/capitulos/01\_introduccion.tex:25. Lo que se le atribuye a la fuente es el carácter de código abierto (y por...}}\par
\begin{pasaje}\footnotesize p. 357: «NumPy is a community-developed, open-source library, which provides a multidimensional Python array object along with array-aware functions that operate on it.» \textbar{}\textbar{} p. 361: «NumPy was initially developed by students, faculty and researchers to provide an advanced, open-source array programming library for Python, which was free to use and unencumbered by license servers and software protection dongles.» y «Owing to NumPy's simple memory model, it is easy to write low-level, hand-optimized code, usually in C or Fortran, to manipulate NumPy arrays and pass them back to Python.»\end{pasaje}
\vspace{3pt}\noindent{\small\textbf{p.\ 360 (PDF pág. 4)} \quad \campo{«Para mallas grandes \$\textbackslash{}bm\{K\}\$ es dispersa (mayoritariamente nula), por lo que EduFEM la construye en formato de coordenadas (COO) y la convierte a filas comprimidas (CSR) para el solucionador, apoyándose en las bibliotecas científicas de Python \textbackslash{}autocite\{harris2020numpy,virtanen2020scipy\}» ---...}}\par
\begin{pasaje}\footnotesize «SciPy and PyData/Sparse both provide sparse arrays, which typically contain few non-zero values and store only those in memory for efficiency.»\end{pasaje}
\vspace{3pt}\noindent{\small\textbf{p.\ 358 (PDF pág. 2) para el álgebra densa;...} \quad \campo{«El núcleo numérico se apoya en NumPy y SciPy, que proveen el álgebra densa y dispersa, la cuadratura y el solucionador lineal del método \textbackslash{}autocite\{harris2020numpy,virtanen2020scipy\}» --- tesis/capitulos/03\_diseno\_implementacion.tex:25 (y la fila «Cálculo numérico \textbar{} NumPy, SciPy \textbar{} Álgebra, cuadratura, solucionador disperso» de...}}\par
\begin{pasaje}\footnotesize p. 358: «It provides extensive support for generating pseudorandom numbers, includes an assortment of probability distributions, and performs accelerated linear algebra, using one of several backends such as OpenBLAS18,19 or Intel MKL optimized for the CPUs at hand (see Supplementary Methods for more details).» \textbar{}\textbar{} p. 359: «SciPy provides fundamental algorithms for scientific computing, including mathematical, scientific and engineering routines. [...] The combination of NumPy, SciPy and Matplotlib, together with an advanced interactive environment such as IPython20 or Jupyter21, provides a solid foundation for array programming in Python.»\end{pasaje}
\vspace{3pt}\noindent{\small\textbf{p.\ 359 (PDF pág. 3); complemento en 357...} \quad \campo{«La elección de un entorno íntegro en Python sobre las bibliotecas científicas de uso extendido \textbackslash{}autocite\{harris2020numpy,virtanen2020scipy\} favorece, además, la inspección y extensión del código por parte de docentes y estudiantes» --- tesis/capitulos/04\_resultados.tex:223. Dos afirmaciones: (a) uso extendido de esas...}}\par
\begin{pasaje}\footnotesize p. 359: «Python is an open-source, general-purpose interpreted programming language well suited to standard programming tasks such as cleaning data, interacting with web resources and parsing text. Adding fast array operations and linear algebra enables scientists to do all their work within a single programming language--one that has the advantage of being famously easy to learn and teach, as witnessed by its adoption as a primary learning language in many universities.» y, más abajo en la misma página: «NumPy and its ecosystem are commonly taught in university courses, boot camps and summer schools, and are the focus of community conferences and workshops worldwide. NumPy and its API have become truly ubiquitous.» \textbar{}\textbar{} p. 357: «Now, 15 years later, NumPy underpins almost every Python library that does scientific or numerical computation8-11, including SciPy12, Matplotlib13, pandas14,...\end{pasaje}
\vspace{6pt}
\subsection{\texttt{hughes2000fem} --- The Finite Element Method: Linear Static and Dynamic Finite...}
{\footnotesize\campo{Ejemplar:} Thomas J. R. Hughes. Prentice-Hall, Inc. (A Division of Simon \& Schuster). 1987. \campo{Archivo:} \texttt{Hughes 1987 - The Finite Element Method.pdf}\par}\vspace{4pt}
\vspace{3pt}\noindent{\small\textbf{p.\ 77-79} \quad \campo{La forma fuerte del equilibrio exige mayor continuidad; el MEF parte de la forma debil obtenida multiplicando por desplazamientos virtuales admisibles (nulos sobre Gamma\_u) e integrando por partes con el teorema de la divergencia, lo que conduce al principio de los trabajos virtuales. --- tesis/capitulos/02\_marco\_teorico.tex:49}}\par
\begin{pasaje}\footnotesize p. 77 (forma fuerte): \textquotedbl{}(S) Given f\_i: Omega -\textgreater{} R, q\_i: Gamma\_\{q\_i\} -\textgreater{} R, and h\_i: Gamma\_\{h\_i\} -\textgreater{} R, find u\_i: Omega-bar -\textgreater{} R such that sigma\_\{ij,j\} + f\_i = 0 in Omega (equilibrium equations) (2.7.8); u\_i = q\_i on Gamma\_\{q\_i\} (2.7.9); sigma\_\{ij\} n\_j = h\_i on Gamma\_\{h\_i\} (2.7.10)\textquotedbl{}. p. 78 (forma debil y trabajos virtuales): \textquotedbl{}The weak formulation goes as follows: (W) Given ... find u\_i in S\_i such that for all w\_i in V\_i, Int\_Omega w\_\{(i,j)\} sigma\_\{ij\} dOmega = Int\_Omega w\_i f\_i dOmega + Sum\_\{i=1\}\textasciicircum{}\{n\_sd\} ( Int\_\{Gamma\_\{h\_i\}\} w\_i h\_i dGamma ) (2.7.11)\textquotedbl{} y \textquotedbl{}Remarks 5. In the solid mechanics literature, (W) is sometimes referred to as the principle of virtual work, or principle of virtual displacements, w\_i being the virtual displacements.\textquotedbl{} p. 79 (obtencion por partes): \textquotedbl{}0 = Int\_Omega w\_i (sigma\_\{ij,j\} + f\_i) dOmega = -Int\_Omega w\_\{i,j\} sigma\_\{ij\} dOmega + Int\_Gamma w\_i sigma\_\{ij\} n\_j dGamma +...\end{pasaje}
\vspace{3pt}\noindent{\small\textbf{p.\ 147 y 191} \quad \campo{El integrando de la rigidez no es polinomico: como B depende de J\textasciicircum{}\{-1\}, que es racional en (xi,eta), la integral carece de primitiva elemental; por ello se aproxima mediante cuadratura de Gauss-Legendre. --- tesis/capitulos/02\_marco\_teorico.tex:235}}\par
\begin{pasaje}\footnotesize p. 147: \textquotedbl{}The matrix (3.9.6) is the inverse of the matrix in (3.9.3), i.e., [xi\_\{,x\} xi\_\{,y\}; eta\_\{,x\} eta\_\{,y\}] = (x\_\{,xi\})\textasciicircum{}\{-1\} = (1/j)[y\_\{,eta\} -x\_\{,eta\}; -y\_\{,xi\} x\_\{,xi\}] (3.9.9) where j = det(x\_\{,xi\}) = x\_\{,xi\} y\_\{,eta\} - x\_\{,eta\} y\_\{,xi\} (3.9.10)\textquotedbl{} (y p. 148, recuadro: \textquotedbl{}N\_\{a,x\} = [N\_\{a,xi\} y\_\{,eta\} - N\_\{a,eta\} y\_\{,xi\}] / j\textquotedbl{}). p. 191 (Effect of Numerical Quadrature): \textquotedbl{}The previously developed theory assumes that we rigorously adhere to the Galerkin recipe. In particular, this means that all integrals need to be calculated exactly. As indicated in Chapter 3, this would be virtually impossible for isoparametric elements in all but the simplest configurations. Because numerically integrated stiffnesses represent the rule rather than the exception, it is important to understand the accuracy implications of approximate numerical integration.\textquotedbl{}\end{pasaje}
\vspace{3pt}\noindent{\small\textbf{p.\ 141-142} \quad \campo{La regla de Gauss de m puntos integra exactamente polinomios de grado 2m-1; los puntos son xi = +-1/raiz(3) para 2 puntos y xi = +-raiz(3/5), 0 con pesos 5/9, 8/9, 5/9 para 3 puntos. --- tesis/capitulos/02\_marco\_teorico.tex:242}}\par
\begin{pasaje}\footnotesize p. 141: \textquotedbl{}Gaussian Quadrature. In one dimension the Gauss quadrature formulas are optimal. Accuracy of order 2 n\_int is achieved by n\_int integration points. The locations of the quadrature points and values of associated weights are determined to attain maximum accuracy.\textquotedbl{} Reglas: \textquotedbl{}1. n\_int = 1: xi\_1 = 0, W\_1 = 2\textquotedbl{}; p. 142: \textquotedbl{}2. n\_int = 2: xi\_1 = -1/raiz3, xi\_2 = 1/raiz3, W\_1 = W\_2 = 1\textquotedbl{}; \textquotedbl{}3. n\_int = 3: xi\_1 = -raiz(3/5), xi\_2 = 0, xi\_3 = raiz(3/5), W\_1 = W\_3 = 5/9, W\_2 = 8/9\textquotedbl{}. p. 143: \textquotedbl{}Gaussian Quadrature in Several Dimensions. Gaussian rules for integrals in several dimensions are constructed by employing one-dimensional Gaussian rules on each coordinate separately\textquotedbl{} (producto tensorial, ec. 3.8.13).\end{pasaje}
\vspace{3pt}\noindent{\small\textbf{p.\ 221-222} \quad \campo{EduFEM usa cuadratura 2x2 para Q4 y 3x3 para Q9, ordenes que integran exactamente la rigidez de los respectivos elementos no distorsionados. --- tesis/capitulos/02\_marco\_teorico.tex:242}}\par
\begin{pasaje}\footnotesize p. 221, clave de la Tabla 4.4.1: \textquotedbl{}U2  2 x 2 uniform integration (= exact in present case)\textquotedbl{}. p. 222, Figura 4.4.3 (Some equivalent elements), columna 'Normal Gaussian quadrature (k-barra)': fila 'Bilinear / Constant' -\textgreater{} \textquotedbl{}2 x 2\textquotedbl{}; fila 'Biquadratic / Bilinear' -\textgreater{} \textquotedbl{}3 x 3\textquotedbl{}; fila 'Trilinear / Constant' -\textgreater{} \textquotedbl{}2 x 2 x 2\textquotedbl{}.\end{pasaje}
\vspace{3pt}\noindent{\small\textbf{p.\ 243 (con Fig. 4.7.1 en p. 244)} \quad \campo{El bloqueo por cortante afecta al Q4 bilineal en estados dominados por flexion: el campo bilineal no representa la curvatura sin generar deformaciones de corte espurias, que absorben energia y vuelven al elemento artificialmente rigido, de modo que subestima los desplazamientos. El Q9, de interpolacion bicuadratica, esta...}}\par
\begin{pasaje}\footnotesize p. 243: \textquotedbl{}It may be recalled that the standard four-node quadrilaterals employed in the beam-bending elasticity example of Sec. 4.4 ... attained a level of accuracy of approximately 90\% in the case of Poisson's ratio = 0.3. ... For 'bending' cases like this, the accuracy level attained by simple elements tends to be somewhat disappointing. If only one element is used through the thickness, the accuracy degrades further, leading to worthless results. ... They noted that if a bending moment is applied to a single element, the element responds in shear rather than bending, and this spurious shearing is responsible for the overly stiff behavior (see Fig. 4.7.1).\textquotedbl{} Nota al pie 10 de la misma pagina: \textquotedbl{}Clearly, another possibility would be simply to employ higher-order elements, such as the eight-node serendipity, or nine-node Lagrange, quadrilateral. As might be anticipated, excellent results...\end{pasaje}
\vspace{3pt}\noindent{\small\textbf{p.\ 192, 208 y 220-221} \quad \campo{El bloqueo volumetrico aparece en deformacion plana cuando nu -\textgreater{} 0,5: el factor (1-2nu) tiende a cero, la restriccion de incompresibilidad sobre-condiciona el elemento, la malla no tiene GDL suficientes para satisfacer eps\_x + eps\_y -\textgreater{} 0 sin rigidizarse, y el resultado es una respuesta artificialmente rigida. ---...}}\par
\begin{pasaje}\footnotesize p. 192: \textquotedbl{}In isotropic linear elasticity the condition of incompressibility may be expressed in terms of Poisson's ratio nu. As nu approaches 1/2, resistance to volume change is greatly increased ... B/mu = 2(1+nu)/(3(1-2nu)) (4.2.1). Clearly, as nu -\textgreater{} 1/2, the ratio approaches infinity. ... Thus the Lame parameter, lambda [= 2 nu mu/(1-2nu)], also becomes unbounded in the incompressible limit\textquotedbl{}; y antes, en la Discussion: \textquotedbl{}for typical finite elements which we have considered so far, does for constrained media problems such as ones involving incompressible, or nearly incompressible, behavior. In these cases the approximation may be so poor that it is practically useless.\textquotedbl{} p. 208: \textquotedbl{}identical reasoning may be used to conclude that every node in the entire mesh must have zero displacement. Thus the only possible incompressible displacement is u\textasciicircum{}h = 0. This result holds no matter how many...\end{pasaje}
\vspace{3pt}\noindent{\small\textbf{p.\ 221 y 232-234} \quad \campo{Existen tecnicas correctivas: la integracion reducida selectiva (SRI), que subintegra solo los terminos responsables del bloqueo, y la formulacion B-barra, que proyecta la parte volumetrica o de corte de la deformacion. --- tesis/capitulos/02\_marco\_teorico.tex:257 y tesis/capitulos/05\_conclusiones.tex:45}}\par
\begin{pasaje}\footnotesize p. 221: \textquotedbl{}Definitions. Selective reduced integration refers to the case in which reduced integration is used only on the lambda-term, whereas normal integration is used on the mu-term. Uniform reduced integration refers to the case when both terms are integrated with a reduced rule. ... Selective integration retains the correct rank of the element stiffness, and therefore the global stiffness also possesses correct rank.\textquotedbl{} p. 232: \textquotedbl{}4.5.2 Strain Projection: The B-bar approach.\textquotedbl{} p. 233: \textquotedbl{}These expressions are standard but must be modified to be successful in application to nearly incompressible cases. Let B\_a\textasciicircum{}dil denote the dilatational part of B\_a ... The deviatoric part of B\_a is then defined by B\textasciicircum{}dev = B\_a - B\_a\textasciicircum{}dil (4.5.7) ... To achieve an effective formulation for nearly incompressible applications, B\_a\textasciicircum{}dil needs to be replaced by an 'improved' dilatational contribution, which we...\end{pasaje}
\vspace{3pt}\noindent{\small\textbf{p.\ 189-190} \quad \campo{Al refinar la malla, las normas de error deben decrecer con las tasas asintoticas teoricas O(h\textasciicircum{}\{p+1\}) en L2 y O(h\textasciicircum{}p) en H1, con p el orden polinomico del elemento, segun la teoria de aproximacion del MEF. --- tesis/capitulos/02\_marco\_teorico.tex:406 y tesis/capitulos/04\_resultados.tex:16}}\par
\begin{pasaje}\footnotesize p. 189: \textquotedbl{}\textbar{}\textbar{}u - U\textasciicircum{}h\textbar{}\textbar{}\_m \textless{}= c h\textasciicircum{}alpha \textbar{}\textbar{}u\textbar{}\textbar{}\_r (4.1.6) where c is a constant independent of u and h, alpha = min(k + 1 - m, r - m), k is the degree of complete polynomial appearing in the element shape functions, and h is the mesh parameter\textquotedbl{}; y el teorema: \textquotedbl{}\textbar{}\textbar{}e\textbar{}\textbar{}\_m \textless{}= c-barra h\textasciicircum{}alpha \textbar{}\textbar{}u\textbar{}\textbar{}\_r\textquotedbl{}. p. 190: \textquotedbl{}Remarks 2. Assume u is smooth in the sense that u in H\textasciicircum{}\{k+1\}. Then the error satisfies \textbar{}\textbar{}e\textbar{}\textbar{}\_m \textless{}= c h\textasciicircum{}\{k+1-m\} \textbar{}\textbar{}u\textbar{}\textbar{}\_\{k+1\}. This is sometimes referred to as the standard error estimate. ... 3. ... For example, let k = m = 1. Then \textbar{}\textbar{}e\textbar{}\textbar{}\_0 \textless{}= c h\textasciicircum{}2 \textbar{}\textbar{}u\textbar{}\textbar{}\_2 ; \textbar{}\textbar{}e\textbar{}\textbar{}\_1 \textless{}= c h \textbar{}\textbar{}u\textbar{}\textbar{}\_2. Using mechanics terminology, the former relation gives the convergence rate in the L\_2-norm for 'displacements,' and the latter gives the convergence rate in the L\_2-norm for 'displacement gradients' (i.e., strains or stresses).\textquotedbl{}\end{pasaje}
\vspace{3pt}\noindent{\small\textbf{p.\ 118-119, 137-138} \quad \campo{La discretizacion con cuadrilateros isoparametricos Q4 y Q9 cubre la interpolacion bilineal y la cuadratica completa y permite ilustrar los efectos del orden de interpolacion sobre la precision. --- tesis/capitulos/01\_introduccion.tex:62}}\par
\begin{pasaje}\footnotesize p. 118: \textquotedbl{}3.3 ISOPARAMETRIC ELEMENTS ... If the element interpolation function u\textasciicircum{}h can be written as u\textasciicircum{}h(xi) = Sum N\_a(xi) d\_a\textasciicircum{}e (3.3.2) the element is said to be isoparametric. The key point to observe in the definition is that the shape functions which define (3.3.1) also serve to define (3.3.2). The bilinear quadrilateral, derived in the previous section, is an example of an isoparametric element.\textquotedbl{} p. 119: \textquotedbl{}Definition 4. ... The determinant of the derivative, denoted by j = det(dx/dxi), is called the Jacobian determinant\textquotedbl{} y \textquotedbl{}3. In actual computations the sign of the Jacobian determinant is monitored at special points to ensure condition (iv) is satisfied. If a zero or negative value is encountered, computations are terminated. Generally this is an indication of an input-data error or an overly distorted element.\textquotedbl{} p. 137: \textquotedbl{}The degree of completeness is generally considered to be the...\end{pasaje}
\vspace{3pt}\noindent{\small\textbf{p.\ 243 (nota al pie 10)} \quad \campo{En la membrana de Cook el desplazamiento converge a 23,96, el Q4 con malla 8x8 da u\_y = 22,08 (7,9 \% por debajo) por shear-locking, y el Q9 lo evita y converge con rapidez. --- tesis/capitulos/02\_marco\_teorico.tex:408}}\par
\begin{pasaje}\footnotesize p. 243, nota al pie 10: \textquotedbl{}Clearly, another possibility would be simply to employ higher-order elements, such as the eight-node serendipity, or nine-node Lagrange, quadrilateral. As might be anticipated, excellent results may be obtained in this fashion.\textquotedbl{}\end{pasaje}
\vspace{6pt}
\subsection{\texttt{lee2015eigenmodes} --- Interactive Simulation of Eigenmodes for Instruction of Finite...}
{\footnotesize\campo{Ejemplar:} Jae Young Lee (Department of Regional Construction Engineering, Chonbuk National University, Jeonju, South Korea); Hee-Ryong Ryu (Rural Development Administration, Haman, South Korea). Wiley Periodicals, Inc. --- revista Computer Applications in Engineering Education (Comput Appl Eng Educ), vol. 23, pp. 872-886. 2015. \campo{Archivo:} \texttt{Lee y Ryu 2015 - Interactive Simulation of Eigenmodes.pdf}\par}\vspace{4pt}
\vspace{3pt}\noindent{\small\textbf{p.\ 872} \quad \campo{«La literatura sobre enseñanza del método coincide en que la simulación interactiva y la manipulación directa de los objetos matemáticos del MEF favorecen un aprendizaje más profundo que la mera exposición teórica o el uso de herramientas opacas» --- tesis/capitulos/01\_introduccion.tex:21 (cita sin localizador, compartida con...}}\par
\begin{pasaje}\footnotesize The visualization is updated instantly in response to the change of model type, element order, integration rule, and boundary constraints, etc. Such interactivity is an essential feature of the simulation process for the experimental study of eigenmodes and associated model behaviors. [...] While a number of finite element packages with application of eigenmode analysis to various problems have been presented, none of these offer educational features that disclose the basis of the computational procedures and related concepts.\end{pasaje}
\vspace{3pt}\noindent{\small\textbf{p.\ 874} \quad \campo{Refuerzo de la misma afirmación de 01\_introduccion.tex:21: la «manipulación directa» de los objetos matemáticos del MEF}}\par
\begin{pasaje}\footnotesize The object of eigenmode visualization can be specified as a single element or multiple elements, depending on which part of the finite element model is selected by user's mouse clicking. [...] These factors as well as the selection of the object can be altered interactively while the display of the eigenmodes is updated in response to the user's action. The interactivity of the process is a crucial feature of the educational tool which allows observation of the finite element characteristics through experimental manipulation of various control factors.\end{pasaje}
\vspace{3pt}\noindent{\small\textbf{p.\ 885} \quad \campo{Evidencia empírica implícita en «favorecen un aprendizaje más profundo» (01\_introduccion.tex:21)}}\par
\begin{pasaje}\footnotesize The use of the software significantly facilitated the explanation to the students and the students' understanding of complex concepts related to finite element formulation. Students have shown more active response to the lectures and attained a higher degree of understanding of the subjects through their experiences from the assigned exercises using the program.\end{pasaje}
\vspace{3pt}\noindent{\small\textbf{p.\ 872} \quad \campo{«Lee desarrolla simuladores interactivos [...] de la visualización de los modos propios de vibración \textbackslash{}autocite\{lee2015eigenmodes\}» --- tesis/capitulos/02\_marco\_teorico.tex:11 (y, de refilón, «visualizadores de modos propios» en 04\_resultados.tex:223)}}\par
\begin{pasaje}\footnotesize The eigenvalues and eigenvectors of the stiffness matrix are interpreted in terms of displacement and strain energy modes representing the basic characteristics of an element or an assembly of elements. [...] They are designated as eigenmodes in this paper, and are associated with various types of modes such as deformation mode, rigid body mode, spurious zero energy mode, and shear locking mode. The displacement modes are the primary eigenmodes corresponding to the eigenvalues and eigenvectors of the stiffness matrix in the finite element equations.\end{pasaje}
\vspace{3pt}\noindent{\small\textbf{p.\ 883} \quad \campo{Matiz de la anterior: hasta qué punto la fuente cubre los modos de vibración (02\_marco\_teorico.tex:11)}}\par
\begin{pasaje}\footnotesize Extension to the Study of Dynamic Analysis and Structural Stability Analysis. Understanding of eigenmodes in static analysis can be extended to the study of dynamic analysis and stability analysis. [...] [K]\{d\} = l[M]\{d\}, for free vibration, or [K]\{d\} = l[G]\{d\}, for bucking failure makes a generalized eigenvalue problem, where [M] and [G] are the mass matrix and geometric stiffness matrix respectively.\end{pasaje}
\vspace{3pt}\noindent{\small\textbf{p.\ 881} \quad \campo{Fila de la tabla comparativa: «Simuladores interactivos \textbackslash{}autocite\{lee2015interactive,lee2015eigenmodes\} \& Inglés \& 2D \& Media-alta, en un eslabón \& Libre \& Continuo \& Parcial (un eslabón)» --- tesis/capitulos/02\_marco\_teorico.tex:30; columna DIMENSIÓN = «2D»}}\par
\begin{pasaje}\footnotesize The phenomenon of spurious zero energy modes can also be observed in the case of three-dimensional solid elements as shown in Figure 7. In this case, the 2 x 2 x 2 integration rule produces a number of spurious modes, while the 3 x 3 x 3 rule suppresses the mechanism. The contour image inside the three-dimensional solid is rendered using the volume visualization technique of VisualFEA [16]. [...] Figure 7 Display of a spurious zero energy mode for a 20 node solid element with 2 x 2 x 2 integration rule.\end{pasaje}
\vspace{3pt}\noindent{\small\textbf{p.\ AUSENTE} \quad \campo{Misma fila de la tabla comparativa (02\_marco\_teorico.tex:30), columna LICENCIA = «Libre»}}\par
\begin{pasaje}\footnotesize (no existe pasaje: el artículo no consigna en ninguna de sus 15 páginas el régimen de licencia, precio, disponibilidad ni condiciones de distribución de VisualFEA; solo afirma en p.\textasciitilde{}872 que «has been implemented as an educational tool embedded in a finite element analysis program, VisualFEA, which was developed by the authors»)\end{pasaje}
\vspace{3pt}\noindent{\small\textbf{p.\ 872} \quad \campo{«Este enfoque materializa el aporte pedagógico central del trabajo: hacer visible y manipulable lo que en la enseñanza tradicional permanece como una caja negra entre los datos de entrada y los resultados» --- tesis/capitulos/05\_conclusiones.tex:17 (cita compartida con perezsantiago2023fem)}}\par
\begin{pasaje}\footnotesize While a number of finite element packages with application of eigenmode analysis to various problems have been presented, none of these offer educational features that disclose the basis of the computational procedures and related concepts. [...] One of the advantages of the tool is that the educational functions are operated in the course of an actual finite element analysis process integrated with the preprocessing and postprocessing stages as schematized in Figure 1. Thus, the concept of eigenmodes and their implication in element behaviors can be understood in association with the practical application of finite element procedures.\end{pasaje}
\vspace{3pt}\noindent{\small\textbf{p.\ 885} \quad \campo{«En relación con las herramientas educativas reportadas en la literatura --simuladores de armaduras, visualizadores de modos propios o entornos de procesamiento interactivo de ecuaciones del MEF--, la contribución distintiva de EduFEM es la cobertura integral del canal de cálculo [...] y una memoria de cálculo con sustitución...}}\par
\begin{pasaje}\footnotesize Students are provided with the educational program as a tool to fulfill assignments of eigenmode visualization for various case studies and numerical experiments. [...] Identify the meaning of the eigenvalues and eigenvector in relationship to element behavior. Students are asked to identify the relationships shown in Table 1, using numerical values of the eigenvalues, eigenvectors and stiffness matrix displayed in text format.\end{pasaje}
\vspace{6pt}
\subsection{\texttt{lee2015interactive} --- Interactive Simulation of Finite Element Equation Processing for...}
{\footnotesize\campo{Ejemplar:} Jae Young Lee (autor unico; Regional Construction Engineering, Chonbuk National University, Jeonju, South Korea). Wiley Periodicals, Inc. --- revista Computer Applications in Engineering Education (Comput Appl Eng Educ), vol. 23, pp. 157-169. 2015 (ano del volumen impreso: \textquotedbl{}Comput Appl Eng Educ 23:157-169, 2015\textquotedbl{}; recibido 21 enero 2013, aceptado 21 septiembre 2013, copyright \textquotedbl{}2013 Wiley Periodicals, Inc.\textquotedbl{}). \campo{Archivo:} \texttt{Lee 2015 - Interactive Simulation of Finite Element Equation Processing.pdf}\par}\vspace{4pt}
\vspace{3pt}\noindent{\small\textbf{p.\ 157} \quad \campo{«La literatura sobre enseñanza del método coincide en que la simulación interactiva y la manipulación directa de los objetos matemáticos del MEF favorecen un aprendizaje más profundo que la mera exposición teórica o el uso de herramientas opacas» --- tesis/capitulos/01\_introduccion.tex:21}}\par
\begin{pasaje}\footnotesize «Students may experience and explore each stage of the computation using the interactive simulation, and attain a better understanding of the concepts and procedures of the stiffness method, bypassing manual calculation or computer programming.» (Resumen, p.157). Y en la Introducción, misma p.157: «VisualFEA provides users with more control of the equation system, more active involvement in computational processes, and a more graphical interpretation of concepts; it thus leads to a more intuitive understanding of the system.»\end{pasaje}
\vspace{3pt}\noindent{\small\textbf{p.\ 157} \quad \campo{«Fundamentada en la literatura sobre enseñanza del MEF, que documenta el valor de la simulación interactiva y de la manipulación directa de los objetos matemáticos del método, esa transparencia constituye un apoyo plausible al aprendizaje de sus fundamentos» --- tesis/capitulos/01\_introduccion.tex:44}}\par
\begin{pasaje}\footnotesize «An interactive simulation with graphical visualization is used as a tool to teach and learn the concepts and procedures related to the construction and solution of finite element stiffness equations. The tool encompasses the stages of processing finite element equations from element modeling to equation solving. Any operation of this educational tool is not an emulation of computational processes, but an actual execution of finite element processing.» (Resumen, p.157)\end{pasaje}
\vspace{3pt}\noindent{\small\textbf{p.\ 168} \quad \campo{«los criterios de diseño pedagógico se fundamentan cualitativamente en la literatura sobre enseñanza del MEF» --- tesis/capitulos/01\_introduccion.tex:56}}\par
\begin{pasaje}\footnotesize «The educational tool presented in this article has been used in the above-mentioned courses for the last six years. In the preceding years, the author gave lectures in a conventional way without using the software. There are no quantitative data measuring the efficiency of the lectures in the two periods. However, the author has qualitatively evaluated the degree of students' interest and understanding on the subjects in the two courses through examinations, home work and semester projects. The educational tool was observed to be conducive to enhancing the students' overall understanding and interest in the subjects.» (sección EVALUATION OF THE EDUCATIONAL EFFICIENCY, p.168)\end{pasaje}
\vspace{3pt}\noindent{\small\textbf{p.\ 157} \quad \campo{«las herramientas profesionales son cajas negras optimizadas para producir resultados, no para revelar el procedimiento intermedio, mientras que el desarrollo puramente manual no escala más allá de uno o dos elementos. La investigación educativa documenta esta brecha y propone el software específicamente didáctico como puente...}}\par
\begin{pasaje}\footnotesize «These computational processes can be understood through number crunching exercises based on either manual calculation work or computer programming. Manual calculation work is very rigorous and tedious and thus is very limited in its extent. Computer programming is not feasible for students who are not familiar with the programming language. Therefore, neither of these two conventional approaches would be attractive for most students in an introductory course in the finite element method. For this reason, a number of computer programs have been developed and used for the computer-aided education of the finite element method and the stiffness matrix equations in particular.» (Introducción, p.157)\end{pasaje}
\vspace{3pt}\noindent{\small\textbf{p.\ 157 (título y resumen) y 168...} \quad \campo{«Lee desarrolla simuladores interactivos del procesamiento de las ecuaciones de elementos finitos» --- tesis/capitulos/02\_marco\_teorico.tex:11}}\par
\begin{pasaje}\footnotesize Título, p.157: «Interactive Simulation of Finite Element Equation Processing for Educational Purposes». Conclusiones, p.168: «The educational tool for teaching and learning the system of finite element equations is one of the features implemented in VisualFEA.»\end{pasaje}
\vspace{3pt}\noindent{\small\textbf{p.\ 157 y 158} \quad \campo{«La segunda familia la integran los simuladores que profundizan en un eslabón aislado del método» --- tesis/capitulos/02\_marco\_teorico.tex:11 --- y la fila de la tabla comparativa «Simuladores interactivos ... Parcial (un eslabón)» --- 02\_marco\_teorico.tex:30}}\par
\begin{pasaje}\footnotesize Resumen, p.157: «The tool encompasses the stages of processing finite element equations from element modeling to equation solving.» Introducción, p.157: «However, the scope of the present work is more comprehensive and integrated than that of previous works, which have been chiefly concerned with the display and simple operations of stiffness matrices using computer graphics.» Y p.158: «The principal aim of these instructional functions is to help students understand the overall procedures of computing and solving the finite element equations in the following four steps. / Shape function and element modeling / Computation of the element equations / Assemblage of the global system equations / Solution of the global system equations.»\end{pasaje}
\vspace{3pt}\noindent{\small\textbf{p.\ AUSENTE} \quad \campo{Fila de la tabla comparativa: régimen de licencia «Libre» para los simuladores interactivos de Lee --- tesis/capitulos/02\_marco\_teorico.tex:30}}\par
\begin{pasaje}\footnotesize No existe. Búsqueda de «licen», «free», «open source», «commercial» en el texto completo de las 13 páginas: cero coincidencias pertinentes. Lo único que consta es la pertenencia del simulador al producto: «one of the features implemented in VisualFEA» (p.168) y «He is the author of VisualFEA» (biografía, p.169).\end{pasaje}
\vspace{3pt}\noindent{\small\textbf{p.\ 168} \quad \campo{«El software educativo interactivo demuestra eficacia para consolidar conceptos abstractos del análisis estructural» --- tesis/capitulos/03\_diseno\_implementacion.tex:9}}\par
\begin{pasaje}\footnotesize «There are no quantitative data measuring the efficiency of the lectures in the two periods. However, the author has qualitatively evaluated the degree of students' interest and understanding on the subjects in the two courses through examinations, home work and semester projects. The educational tool was observed to be conducive to enhancing the students' overall understanding and interest in the subjects.» (p.168). Conclusiones, misma p.168: «It has been demonstrated that the diverse functions of the program facilitate the teaching and understanding of the concepts and methods related to the computation of the finite element equations.»\end{pasaje}
\vspace{3pt}\noindent{\small\textbf{p.\ 158 y 166} \quad \campo{«El diseño de estos módulos parte de la evidencia sobre la interactividad como motor del aprendizaje en ingeniería» --- tesis/capitulos/03\_diseno\_implementacion.tex:155}}\par
\begin{pasaje}\footnotesize p.158: «The numerical value of each shape function is computed and displayed using the natural coordinates at the position of the screen cursor. The values are continuously updated while the screen cursor moves over the element.» p.166: «Another purpose of the manual assembly option is for students to become more actively involved in the assembly procedure and thus to explore the computational processes in a similar way to the actual experience.»\end{pasaje}
\vspace{3pt}\noindent{\small\textbf{p.\ 163 y 166} \quad \campo{«Cada módulo es una capa superpuesta interactiva que ilumina sobre la malla real el punto donde ocurre el cálculo, evitando la abstracción de una ventana desconectada del modelo» --- tesis/capitulos/04\_resultados.tex:207}}\par
\begin{pasaje}\footnotesize p.163: «...the assembled destination in the global stiffness matrix is highlighted when the source entry of the element stiffness matrix is designated, as shown in Figure 5. At the same time, the corresponding degree of freedom is drawn over the model image...» p.166: «When an entry of the global stiffness matrix is clicked, its contributing elements are highlighted in the model image. Alternatively, the entries assembled by the element stiffness matrix are highlighted by clicking the image of an element in the model view.» y «When a node is selected from the view of the model, all entries associated with the nodal degrees of freedom are highlighted.»\end{pasaje}
\vspace{3pt}\noindent{\small\textbf{p.\ 157, 158, 161, 167} \quad \campo{«En relación con las herramientas educativas reportadas en la literatura ---simuladores de armaduras, visualizadores de modos propios o entornos de procesamiento interactivo de ecuaciones del MEF---, la contribución distintiva de EduFEM es la cobertura integral del canal de cálculo ... articulada con una generación automática de...}}\par
\begin{pasaje}\footnotesize Alcance declarado, p.157: «The tool encompasses the stages of processing finite element equations from element modeling to equation solving.» Etapas, p.158: los cuatro pasos (función de forma y modelado, ecuaciones del elemento, ensamblaje, solución). Detalle numérico, p.161: «All the variables and intermediate calculations are arranged in a hierarchical tree view following the sequence of computation... The numerical values are updated in real time as the user changes the geometry or the properties of the element.» Materiales docentes, p.167: «Preparation of teaching materials: teaching materials, such as computer presentation slides, have been conveniently prepared using image clips or animation frames created by the program.»\end{pasaje}
\vspace{3pt}\noindent{\small\textbf{p.\ AUSENTE} \quad \campo{«...todo ello con una orientación deliberadamente minimalista coherente con su propósito didáctico» --- tesis/capitulos/05\_conclusiones.tex:15}}\par
\begin{pasaje}\footnotesize No existe. Búsqueda de «minimal» en el texto completo: cero coincidencias. Lo más próximo, p.158: «VisualFEA is unique in its intuitive and user-friendly features of graphical visualization and simulation for learning the concepts and procedures in structural mechanics through drill-and-practice». La fuente presenta, de hecho, una herramienta de gran densidad funcional (cuatro algoritmos de ensamblaje-solución, vistas de árbol expandibles, animación, ensamblaje manual), no minimalista.\end{pasaje}
\vspace{3pt}\noindent{\small\textbf{p.\ 158 y 168} \quad \campo{«los antecedentes revisados indican, en cambio, qué conceptos conviene poner a prueba» (en una futura validación pedagógica con estudiantes) --- tesis/capitulos/05\_conclusiones.tex:51}}\par
\begin{pasaje}\footnotesize p.158: «The principal aim of these instructional functions is to help students understand the overall procedures of computing and solving the finite element equations in the following four steps. / Shape function and element modeling / Computation of the element equations / Assemblage of the global system equations / Solution of the global system equations.» p.168: «...the author has qualitatively evaluated the degree of students' interest and understanding on the subjects in the two courses through examinations, home work and semester projects.»\end{pasaje}
\vspace{6pt}
\subsection{\texttt{oberkampf2010vv} --- Verification and Validation in Scientific Computing}
{\footnotesize\campo{Ejemplar:} William L. Oberkampf; Christopher J. Roy. Cambridge University Press. 2010. \campo{Archivo:} \texttt{Oberkampf y Roy 2010 - Verification and Validation in Scientific Computing.pdf}\par}\vspace{4pt}
\vspace{3pt}\noindent{\small\textbf{p.\ 309-312 (Seccion 8.3, 'Richardson...} \quad \campo{02b\_diseno\_metodologico.tex:14 --- 'la incertidumbre numerica se cuantifica con las tasas de convergencia y, cuando la referencia no es exacta, con la extrapolacion de Richardson de la propia secuencia de mallas'.}}\par
\begin{pasaje}\footnotesize p.309: «The basic concept behind Richardson extrapolation (Richardson, 1911, 1927) is as follows. If one knows the formal rate of convergence of a discretization method with mesh refinement, and if discrete solutions on two systematically refined meshes are available, then one can use this information to obtain an estimate of the exact solution to the mathematical model. Depending on the level of confidence one has in this estimate, it can be used either to correct the fine mesh solution or to provide a discretization error estimate for it.» // p.311 Eq. (8.62), extrapolacion generalizada de orden p con factor de refinamiento r. // p.322: «the much more common case is when the formal order does not match the observed order. In this case, the error estimate is much less reliable and should generally be converted into a numerical uncertainty.»\end{pasaje}
\vspace{3pt}\noindent{\small\textbf{p.\ 13 (Seccion 1.2.3.3); 14 (aforismo de...} \quad \campo{02b\_diseno\_metodologico.tex:16 --- 'La evaluacion de la correccion se realizo mediante un esquema de verificacion y validacion: verificacion por el metodo de soluciones manufacturadas y VALIDACION contra la viga de Timoshenko y la membrana de Cook'.}}\par
\begin{pasaje}\footnotesize p.13: «Verification is the process of assessing software correctness and numerical accuracy of the solution to a given mathematical model. Validation is the process of assessing the physical accuracy of a mathematical model based on comparisons between computational results and EXPERIMENTAL DATA.» // p.14: «In verification, the primary benchmarks are highly accurate solutions to specific mathematical models. In validation, the benchmarks are high-quality experimental measurements.» // p.174 (sobre SAP2000): «code comparisons do not provide substantive evidence that software is functioning correctly (Trucano et al., 2003). Thus code-to-code comparisons should not be used as a substitute for rigorous code verification assessments.»\end{pasaje}
\vspace{3pt}\noindent{\small\textbf{p.\ 208 (encabezado del Cap. 6, 'Exact...} \quad \campo{02b\_diseno\_metodologico.tex:74 --- la seleccion INTENCIONAL de casos por valor probatorio 'es la practica establecida en la verificacion y validacion de codigos de calculo'.}}\par
\begin{pasaje}\footnotesize p.208: «When used for code verification, the ability of this exact solution to EXERCISE ALL TERMS in the mathematical model is more important than any physical realism of the solution. In fact, realistic exact solutions are often avoided for code verification due to the presence of singularities and/or discontinuities.» // p.190-191: «When conducting code verification studies, it is recommended that the MOST GENERAL MESH TOPOLOGY that will be employed for solving the problems of interest be used for the code verification studies.» // p.172: «There are a number of different criteria that can be used for verifying a scientific computing code. In order of increasing rigor, these criteria are: 1 simple tests, 2 code-to-code comparisons, 3 discretization error quantification, 4 convergence tests, and 5 order-of-accuracy tests.» // p.33: «Oberkampf and Trucano (2007) divided the types of...\end{pasaje}
\vspace{3pt}\noindent{\small\textbf{p.\ 13-14 (Seccion 1.2.3.3, 'Verification...} \quad \campo{02\_marco\_teorico.tex:394 --- 'Un programa de calculo debe distinguir entre verificacion ---comprobar que las ecuaciones se resuelven correctamente--- y validacion ---comprobar que las ecuaciones representan la realidad fisica---'.}}\par
\begin{pasaje}\footnotesize p.13-14: «Blottner (1990) captured the essence of each in the compact expressions: \textquotedbl{}Verification is solving the equations right;\textquotedbl{} \textquotedbl{}Validation is solving the right equations.\textquotedbl{}» y, en la misma pagina, «Verification is the process of assessing software correctness and numerical accuracy of the solution to a given mathematical model. Validation is the process of assessing the physical accuracy of a mathematical model based on comparisons between computational results and experimental data. [...] In verification, the association or relationship of the simulation to the real world is not an issue. In validation, the relationship between the mathematical model and the real world (experimental data) is the issue.» Ver tambien pp. 32-35 (definiciones ASME de code verification, solution verification y model validation).\end{pasaje}
\vspace{3pt}\noindent{\small\textbf{p.\ 219-222 (Seccion 6.3, 'Method of...} \quad \campo{02\_marco\_teorico.tex:396 --- procedimiento del MMS: elegir a priori un campo exacto, sustituirlo en las ecuaciones de equilibrio para deducir la fuerza de cuerpo, e inyectar esa carga con las condiciones de Dirichlet no homogeneas correspondientes.}}\par
\begin{pasaje}\footnotesize p.220: «The general concept behind MMS is to choose a solution a priori, then operate the governing PDEs onto the chosen solution, thereby generating additional analytic source terms that require no discretization. [...] MMS involves the solution to the backward problem: given an original set of equations and a chosen solution, find a modified set of equations that the chosen solution will satisfy.» // p.221, Steps 1-4: «Step 1 Establish the mathematical model in the form L(u)=0 [...] Step 2 Choose the analytic form of the manufactured solution u\textasciicircum{}. Step 3 Operate the mathematical model L(.) onto the manufactured solution u\textasciicircum{} to obtain the analytic source term s=L(u\textasciicircum{}). Step 4 Obtain the modified form of the mathematical model by including the analytic source term L(u)=s.» y, para las restricciones de Dirichlet: «The initial and boundary conditions can be obtained directly from the...\end{pasaje}
\vspace{3pt}\noindent{\small\textbf{p.\ 13-14; 32-37 (Seccion 2.2, 'Primary...} \quad \campo{04\_resultados.tex:5 --- 'siguiendo la distincion metodologica clasica entre verificacion y validacion establecida en el capitulo de marco teorico'.}}\par
\begin{pasaje}\footnotesize Mismo pasaje de p.13-14 citado arriba. Ademas p.32: «This book will use the definitions of V\&V as given by the ASME Guide. Also defined and discussed in this section are the terms code verification, solution verification, predictive capability, calibration, certification, and accreditation.» --- el libro es precisamente la exposicion sistematica de esa distincion clasica.\end{pasaje}
\vspace{3pt}\noindent{\small\textbf{p.\ 309-312 (Seccion 8.3, incluida la forma...} \quad \campo{04\_resultados.tex:150 (nota al pie) --- la extrapolacion de Richardson de la secuencia Q9 (23,925 -\textgreater{} 23,949 -\textgreater{} 23,961) apunta a \textasciitilde{}23,97; el orden observado, proximo a uno, refleja la singularidad de tensiones en las esquinas reentrantes; la incertidumbre resultante es del orden del 0,05 \%.}}\par
\begin{pasaje}\footnotesize p.311, Eq. (8.62): «f\_barra = f\_h + (f\_h - f\_rh)/(r\textasciicircum{}p - 1)» --- extrapolacion de Richardson generalizada a esquemas de orden p y factor de refinamiento r arbitrario, exactamente la forma que exige una secuencia Q9 con orden observado \textasciitilde{}1. // p.312: «There are five basic assumptions required for Richardson extrapolation to provide reliable estimates [...] (1) that both discrete solutions are in the asymptotic range, (2) that the meshes have a uniform (Cartesian) spacing [...] (4) that the solutions are smooth, and (5) that the other sources of numerical error are small.» // p.182: «the formal order of accuracy can be reduced in the presence of DISCONTINUITIES AND SINGULARITIES in the solution.» --- respalda literalmente la explicacion del orden proximo a uno. // p.322: «when the formal order does not match the observed order [...] the error estimate is much less reliable and should generally be...\end{pasaje}
\vspace{3pt}\noindent{\small\textbf{p.\ 13-14; 171-175; 219-222; 309-312 (el...} \quad \campo{05\_conclusiones.tex:19 --- 'El marco metodologico de verificacion y validacion adoptado sigue las practicas establecidas en la literatura de computacion cientifica'; y, en la misma linea, 'la validacion con la membrana de Cook'.}}\par
\begin{pasaje}\footnotesize Contratapa/preliminares del libro: «This book provides a comprehensive and systematic development of the basic concepts, principles, and procedures for verification and validation of models and simulations.» // p.174: «code comparisons do not provide substantive evidence that software is functioning correctly [...] code-to-code comparisons should not be used as a substitute for rigorous code verification assessments.»\end{pasaje}
\vspace{3pt}\noindent{\small\textbf{p.\ 320 (Seccion 8.5, requisitos de...} \quad \campo{02\_marco\_teorico.tex:406 --- 'Las integrales de error se evaluan con un punto de Gauss mas por direccion que la rigidez (3x3 para Q4, 4x4 para Q9) para no subestimar artificialmente el error, porque el error de la propia cuadratura puede interferir con el error de discretizacion y degradar el orden observado'. ATRIBUCION...}}\par
\begin{pasaje}\footnotesize p.320: «This additional requirement pertains to the order of accuracy of any numerical approximations used to compute the system response quantity. When the system response quantity is an integral, a derivative, or an average, then the numerical approximations used in its evaluation must be of at least the same order of accuracy as the underlying discretization scheme.» // y a continuacion: «However, in some cases the errors due to the numerical quadrature can interact with the numerical errors in the discrete solution and adversely impact the observed order of accuracy computation.» --- es exactamente el motivo de integrar las normas de error con una regla mas fina que la de la rigidez.\end{pasaje}
\vspace{6pt}
\subsection{\texttt{onate2009structural} --- Structural Analysis with the Finite Element Method. Linear...}
{\footnotesize\campo{Ejemplar:} Eugenio Onate (Onate, Eugenio). CIMNE (International Center for Numerical Methods in Engineering) y Springer --- logotipos conjuntos en portada interior; el prologo agradece \textquotedbl{}the joint publication of the book by CIMNE and Springer\textquotedbl{}; el copyright es solo de CIMNE. 2009. \campo{Archivo:} \texttt{Onate 2009 - Structural Analysis with the Finite Element Method Vol 1.pdf}\par}\vspace{4pt}
\vspace{3pt}\noindent{\small\textbf{p.\ 193} \quad \campo{capitulos/01\_introduccion.tex:62 --- «La discretización se realiza con elementos cuadriláteros isoparamétricos de cuatro nodos (Q4) y de nueve nodos (Q9), que cubren respectivamente la interpolación bilineal y la cuadrática completa [...] \textbackslash{}autocite\{hughes2000fem,onate2009structural\}». Parte 1: el Q4 es de interpolación BILINEAL.}}\par
\begin{pasaje}\footnotesize The straight-sided element geometry is exactly approximated by the bilinear (subparametric) interpolation  x = sum\_\{i=1\}\textasciicircum{}\{4\} N\_i x\_i  (6.22)  where N\_i = 1/4 (1 + xi xi\_i)(1 + eta eta\_i) are the shape functions for the 4-noded rectangle.\end{pasaje}
\vspace{3pt}\noindent{\small\textbf{p.\ 164} \quad \campo{capitulos/01\_introduccion.tex:62 --- Parte 2: el Q9 cubre la interpolación «cuadrática completa».}}\par
\begin{pasaje}\footnotesize The 9-noded Lagrange rectangle contains all the terms of a complete quadratic polynomial plus three additional terms of the cubic and quartic polynomials (xi\textasciicircum{}2 eta, xi eta\textasciicircum{}2, xi\textasciicircum{}2 eta\textasciicircum{}2). Therefore, the approximation is simply of quadratic order.\end{pasaje}
\vspace{3pt}\noindent{\small\textbf{p.\ 183-184} \quad \campo{capitulos/01\_introduccion.tex:62 --- Parte 3: los elementos Q4/Q9 «permiten ilustrar de manera contrastada los efectos del orden de interpolación sobre la precisión».}}\par
\begin{pasaje}\footnotesize 5.7 GENERAL PERFORMANCE OF TRIANGULAR AND RECTANGULAR ELEMENTS. We present next two examples which lead us to draw some conclusions on the behaviour of rectangular and triangular elements. [...] Different meshes of 3 and 6-noded triangles and 4 and 9-noded rectangles are used for the analysis. Numerical results [...] show that the 3-noded triangle is the less accurate of all elements studied. [...] The accuracy increases notably for the same number of DOFs when 6-noded triangles are used and, even more, when either the 4- or the 9-noded rectangles are used, as shown in Figure 5.16.\end{pasaje}
\vspace{3pt}\noindent{\small\textbf{p.\ 187} \quad \campo{capitulos/02\_marco\_teorico.tex:101 --- «EduFEM emplea elementos isoparamétricos cuadriláteros, en los que la misma familia de funciones de forma interpola tanto la geometría como el campo de desplazamientos \textbackslash{}autocite\{bathe2014fem,onate2009structural\}».}}\par
\begin{pasaje}\footnotesize 6.2 ISOPARAMETRIC QUADRILATERAL ELEMENTS. We recall that the term \textquotedbl{}isoparametric\textquotedbl{} means that the displacement shape functions are used to interpolate the element geometry in terms of the nodal coordinates. Thus, the geometry of a 2D isoparametric quadrilateral with n nodes is expressed as  x = sum\_\{i=1\}\textasciicircum{}\{n\} N\_i(xi,eta) x\_i ; y = sum\_\{i=1\}\textasciicircum{}\{n\} N\_i(xi,eta) y\_i   (6.1)\end{pasaje}
\vspace{3pt}\noindent{\small\textbf{p.\ 159-160} \quad \campo{capitulos/02\_marco\_teorico.tex:101 --- «Las funciones se definen sobre un elemento de referencia (el cuadrado natural \$(\textbackslash{}xi,\textbackslash{}eta)\textbackslash{}in[-1,1]\textasciicircum{}2\$), y un mapeo isoparamétrico las lleva a la geometría física».}}\par
\begin{pasaje}\footnotesize A local coordinate system xi, eta is defined for each element in order to facilitate the derivation of the shape functions. Such a natural or intrinsic coordinate system is normalized so that the element sides are located at xi = $\pm$1 and eta = $\pm$1 as shown in Figure 5.2. [...] Note that in the natural coordinate system the rectangular element becomes a square of side equal to two.\end{pasaje}
\vspace{3pt}\noindent{\small\textbf{p.\ 161 y 164} \quad \campo{capitulos/02\_marco\_teorico.tex:113 (mismo bloque citado en la línea 101) --- «El elemento Q9 [...] alcanzando interpolación bicuadrática como producto tensorial de polinomios de Lagrange unidimensionales», con las funciones de forma de la Ec. eq:N-q9.}}\par
\begin{pasaje}\footnotesize p.161: \textquotedbl{}The shape functions for 2D Lagrange elements can be obtained by the simple product of the two normalized 1D Lagrange polynomials corresponding to the natural coordinates xi and eta of the node. [...] N\_i(xi,eta) = l\_I\textasciicircum{}i(xi) l\_J\textasciicircum{}i(eta)  (5.10)\textquotedbl{} --- p.164: \textquotedbl{}5.3.2 Nine-noded quadratic Lagrange rectangle. The shape functions for the 9-noded Lagrange rectangle (Figure 5.5) are obtained by the product of two normalized 1D quadratic polynomials in xi and eta. [...] a) Corner nodes N\_i = 1/4 (xi\textasciicircum{}2 + xi xi\_i)(eta\textasciicircum{}2 + eta eta\_i), i = 1,3,5,7 (5.15) [...] c) Central node N\_9(xi,eta) = (1 - xi\textasciicircum{}2)(1 - eta\textasciicircum{}2) (5.17)\textquotedbl{}\end{pasaje}
\vspace{3pt}\noindent{\small\textbf{p.\ 175-176} \quad \campo{capitulos/05\_conclusiones.tex:47 --- «La biblioteca de elementos puede ampliarse con elementos triangulares (lineales y cuadráticos) [...] \textbackslash{}autocite\{onate2009structural,reddy2006introduction\}».}}\par
\begin{pasaje}\footnotesize p.175: \textquotedbl{}5.5.3 Shape functions for the 3-noded linear triangle. The shape functions for the 3-noded triangle are linear polynomials (M = 1).\textquotedbl{} --- p.176: \textquotedbl{}5.5.4 Shape functions for the six-noded quadratic triangle. The shape functions for this element are complete quadratic polynomials (M = 2).\textquotedbl{}\end{pasaje}
\vspace{3pt}\noindent{\small\textbf{p.\ 167 (transición); 193-194 y 214...} \quad \campo{capitulos/05\_conclusiones.tex:47 --- «...y con cuadriláteros de transición, ofreciendo así una variedad que enriquezca el estudio comparativo de la convergencia y del comportamiento bajo distorsión de malla».}}\par
\begin{pasaje}\footnotesize p.167: \textquotedbl{}Lagrange elements can have different number of nodes in each local direction xi or eta (Figure 5.7). [...] Therefore, the degree of approximation of the element does not change by simply increasing the number of nodes in one of the two local directions only. This explains why these elements are not very popular and they are only occasionally used as a transition between elements of two different orders.\textquotedbl{} --- p.194: \textquotedbl{}We conclude that the 9-noded element can better represent quadratic cartesian polynomials on linearly distorted shapes and therefore is generally preferable for modelling smooth solutions.\textquotedbl{} --- p.214: \textquotedbl{}The distortion of rectangular elements to quadrilateral shapes leads to stiffer results [CMPW].\textquotedbl{}\end{pasaje}
\vspace{6pt}
\subsection{\texttt{perezsantiago2023fem} --- FEM education in undergraduate studies: Industry-informed research}
{\footnotesize\campo{Ejemplar:} Rogelio Perez-Santiago; Esmeralda Campos (Tecnologico de Monterrey). Computer Applications in Engineering Education (Wiley). 2023. \campo{Archivo:} \texttt{Perez-Santiago y Campos 2023 - FEM Education in Undergraduate Studies.pdf}\par}\vspace{4pt}
\vspace{3pt}\noindent{\small\textbf{p.\ 1160} \quad \campo{La ensenanza del MEF forma parte del nucleo curricular de las carreras de ingenieria (01\_introduccion.tex, primer parrafo).}}\par
\begin{pasaje}\footnotesize Furthermore, this work may also support other undergraduate degree programs where FEM/FEA is part of the curriculum.\end{pasaje}
\vspace{3pt}\noindent{\small\textbf{p.\ 1160} \quad \campo{El software comercial opera, desde la perspectiva del estudiante, como una caja negra (01\_introduccion.tex, Planteamiento del problema).}}\par
\begin{pasaje}\footnotesize Other authors are concerned about learners using commercial software like a \textquotedbl{}black box\textquotedbl{} [28]. Thus, some people prefer to cover not only FEM theory but also the principles of strength of materials [31] or combine theory and practice with solid mechanics experiments [16].\end{pasaje}
\vspace{3pt}\noindent{\small\textbf{p.\ 1159 (resumen)} \quad \campo{La formacion privilegia las destrezas practicas por encima de los conceptos fundamentales (02\_marco\_teorico.tex:9 y 04\_resultados.tex).}}\par
\begin{pasaje}\footnotesize The results suggest that specialists agree on the importance of including a mixture of theoretical and applied topics in the syllabus but prefer practical skills over fundamental concepts. Besides learning from defeaturing to postprocessing a model, engineering students must know how to plan, verify, and validate their finite element studies.\end{pasaje}
\vspace{3pt}\noindent{\small\textbf{p.\ 1168} \quad \campo{Conectar los fundamentos teoricos con la operacion del software es el objetivo formativo (sostiene la orientacion de los modulos educativos, 04\_resultados.tex).}}\par
\begin{pasaje}\footnotesize This would allow the student to connect the theoretical fundamentals with the operation of FEA codes while the time required for alternative approaches, like potential energy and virtual work, may be allocated to other topics.\end{pasaje}
\vspace{3pt}\noindent{\small\textbf{p.\ 1160} \quad \campo{01\_introduccion (Planteamiento): una encuesta profesional recogida por el articulo muestra que solo el 31 \% de los ingenieros se considera conocedor del analisis por elementos finitos. \textbackslash{}autocite[p.\textasciitilde{}1160]\{perezsantiago2023fem\}}}\par
\begin{pasaje}\footnotesize However, the same study reveals that only 31\% of the respondents consider themselves ``knowledgeable'' about FEA, suggesting the need to strengthen FEM education in undergraduate engineering studies. In Mechanical Engineering (ME) programs, FEM education must prepare the students for their proficiency and adequate use in structural and machine design activities. Several factors may impact the course outline, like the a\end{pasaje}
\vspace{3pt}\noindent{\small\textbf{p.\ 1162-1163} \quad \campo{01\_introduccion, 02\_marco\_teorico §1.2 (cuarto principio) y 02b §2.1.4: el consenso enumera 15 conceptos teoricos (Tabla 1) y 34 destrezas practicas (Tabla 2), de las que la tesis toma el universo de contenido. \textbackslash{}autocite[pp.\textasciitilde{}1162-1163]\{perezsantiago2023fem\}}}\par
\begin{pasaje}\footnotesize DFEA4: Meshing FEA13---Selection of elements (1D, 2D, etc.) FEA14---Definition of element size FEA15---Understanding the impact of element aspect ratio and distortion in the solution accuracy FEA16---Understanding the impact of element order on the solution accuracy FEA17---Understanding mesh refinement and mesh density convergence FEA18---Refining the mesh at stress concentration or high?stress gradient regions DFEA5: BC FEA19---\end{pasaje}
\vspace{3pt}\noindent{\small\textbf{p.\ 1168} \quad \campo{02b §2.1.4 y 04 §3.7: las dimensiones de planificacion, mallado y V\&V son las que el consenso califica mas alto. \textbackslash{}autocite[p.\textasciitilde{}1168]\{perezsantiago2023fem\}}}\par
\begin{pasaje}\footnotesize FEA1---Planning 449 5 90 FEA4---Meshing 521 6 87 FEA8---V\&V 251 3 84 FEA2---Modeling and defeaturing 328 4 82 FEA3---Materials 246 3 82 FEA7---Concepts 246 3 82 FEA5---BC 552 7 79 FEA6---Postprocessing 236 3 79 Abbreviations: BC, Boundary Conditions; FEA, Finite Element Analysis; V\&V, Verification and Validation. FIGURE 5 Frequency percentage for all items of the type category. 1168 \textbar{} PEREZ?SANTIAGO and CAMPOS 10990542, 2023, 5, Dow\end{pasaje}
\vspace{3pt}\noindent{\small\textbf{p.\ 1171} \quad \campo{01\_introduccion (Planteamiento): el egresado debe aprender a planificar, verificar y validar sus estudios (conclusion 3). \textbackslash{}autocite[p.\textasciitilde{}1171]\{perezsantiago2023fem\}}}\par
\begin{pasaje}\footnotesize Specialists have confirmed that practical FEA com- petencies should be included in the FEM course content. Besides the generally accepted stages of modeling to postprocessing, engineers must learn how to plan, verify, and validate their FEA studies. (4) The historical focus of FEM?related courses on linear static problems is justified by its predominance in structural design. Nevertheless, procedures to ana- lyze mod\end{pasaje}
\vspace{3pt}\noindent{\small\textbf{p.\ 1171} \quad \campo{02b §2.1.5: el propio estudio declara que sus resultados no son generalizables a otros contextos geograficos. \textbackslash{}autocite[p.\textasciitilde{}1171]\{perezsantiago2023fem\}}}\par
\begin{pasaje}\footnotesize the results of this work may not be generalized to other geographical contexts because the industry's needs may vary from country to country. Despite these limitations, the findings from this study should benefit the design of FEM?related courses and offer guidance to course instructors about the four surveyed aspects. At the authors' university, the traditional 16?week FEM course evolved into the gradual exposition\end{pasaje}
\vspace{6pt}
\subsection{\texttt{reddy2006introduction} --- An Introduction to the Finite Element Method}
{\footnotesize\campo{Ejemplar:} J. N. Reddy. McGraw-Hill (The McGraw-Hill Companies, Inc.); International Edition con derechos exclusivos de McGraw-Hill Education (Asia); impreso en Singapur. 2006 (copyright (c) 2006, 1993, 1984; International Edition 2006). \campo{Archivo:} \texttt{Reddy 2006 - An Introduction to the Finite Element Method.pdf}\par}\vspace{4pt}
\vspace{3pt}\noindent{\small\textbf{p.\ 560 (ejemplo ilustrativo; PDF 569)} \quad \campo{Encuadre: \textquotedbl{}los textos clásicos desarrollan la teoría con un grado de abstracción considerable, donde las funciones de forma, la matriz Jacobiana, la matriz de deformación-desplazamiento y la integración de Gauss aparecen como pasos de una cadena algebraica cuya conexión con la geometría concreta de un elemento no siempre...}}\par
\begin{pasaje}\footnotesize Example 9.3.4 --- Consider the quadratic triangular element shown in Fig. 9.3.7. We wish to calculate $\partial$$\psi$1/$\partial$x, $\partial$$\psi$1/$\partial$y, $\partial$$\psi$4/$\partial$x, and $\partial$$\psi$4/$\partial$y at the point (x, y) = (2, 4) and evaluate the integral of the product ($\partial$$\psi$1/$\partial$x)($\partial$$\psi$4/$\partial$x). [...] [J] = [[-2, -6],[5, -4]],  [J]\textasciicircum{}-1 = (1/38)[[-4, 6],[-5, -2]]\end{pasaje}
\vspace{3pt}\noindent{\small\textbf{p.\ 612} \quad \campo{\textquotedbl{}El MEF parte por ello de la forma débil o variacional, obtenida multiplicando la ecuación de equilibrio fuerte por un campo de desplazamientos virtuales $\delta$u admisible (nulo sobre $\Gamma$u) e integrando por partes mediante el teorema de la divergencia. Esto conduce al principio de los trabajos virtuales\textquotedbl{}, seguido de la ecuación...}}\par
\begin{pasaje}\footnotesize 11.3.2 Principle of Virtual Displacements in Vector Form --- Here, we use the (dynamic version of) the principle of virtual displacements [see Reddy (2002)] applied to a plane elastic finite element $\Omega$\_e with volume V\_e = h\_e $\Omega$\_e (see Fig. 11.3.1)

0 = $\int$\_\{V\_e\} ($\sigma$\_ij $\delta$$\varepsilon$\_ij + $\rho$ ü\_i $\delta$u\_i) dV $\-$ $\int$\_\{V\_e\} f\_i $\delta$u\_i dV $\-$ ?\_\{s\_e\} t?\_i $\delta$u\_i ds     (11.3.1)

where s\_e is the surface of the volume element V\_e, h\_e is the thickness of the element, $\delta$ denotes the variational operator, $\sigma$\_ij and $\varepsilon$\_ij are the components of stress and strain tensors, respectively, and f\_i and t?\_i are the components of the body force and boundary stress vectors, respectively. [...] The first term in Eq. (11.3.1) corresponds to the virtual strain energy stored in the body, the second term corresponds to the kinetic energy stored in the body, the third term represents the virtual work done by the body forces, and the fourth term...\end{pasaje}
\vspace{3pt}\noindent{\small\textbf{p.\ 613} \quad \campo{Misma afirmación de tesis/capitulos/02\_marco\_teorico.tex:49, en su versión matricial-vectorial ($\delta$$\varepsilon$ = D $\delta$u, $\sigma$ = C Du) y en el procedimiento de obtención de la forma débil por multiplicación por función de peso e integración por partes}}\par
\begin{pasaje}\footnotesize Equation (11.3.3) can be rewritten using the notation introduced in Eqs. (11.2.13)--(11.2.15) (note that $\delta$$\varepsilon$ = D$\delta$u and (AB)\textasciicircum{}T = B\textasciicircum{}T A\textasciicircum{}T)

0 = $\int$\_\{$\Omega$\_e\} h\_e [(D$\delta$u)\textasciicircum{}T C (Du) + $\rho$ u\textasciicircum{}T ü] dx $\-$ $\int$\_\{$\Omega$\_e\} ($\delta$u)\textasciicircum{}T f dx $\-$ ?\_\{$\Gamma$\_e\} ($\delta$u)\textasciicircum{}T t ds     (11.3.4)

11.3.3 Weak Form of the Governing Differential Equations --- [...] We use the three-step procedure for each of the two differential equations: multiply the first equation with a weight function w1 and integrate by parts to trade the differentiation equally between the weight function and the dependent variables (u\_x, u\_y).\end{pasaje}
\vspace{3pt}\noindent{\small\textbf{p.\ 526} \quad \campo{\textquotedbl{}La biblioteca de elementos puede ampliarse con elementos triangulares (lineales y cuadráticos)...\textquotedbl{} --- tesis/capitulos/05\_conclusiones.tex:47 (\textbackslash{}autocite\{onate2009structural,reddy2006introduction\}, sin localizador)}}\par
\begin{pasaje}\footnotesize Figure 9.2.1 Topmost six rows of Pascal's triangle for the generation of the Lagrange family of triangular elements. [...] For example, a triangular element of order 2 (i.e., the degree of the polynomial is 2) contains six nodes, as can be seen from the third row of Pascal's triangle. The position of the six nodes in the triangle is at the three vertices and at the midpoints of the three sides. The polynomial involves six constants, which can be expressed in terms of the nodal values of the variable being interpolated: u = $\Sigma$\_\{i=1\}\textasciicircum{}\{6\} u\_i $\psi$\_i(x, y)   (9.2.1) --- where $\psi$\_i are the quadratic interpolation functions obtained following the same procedure as that used for the linear element in Section 8.2. In general, a pth-order triangular element has a number of n nodes: n = ½(p + 1)(p + 2)   (9.2.2)\end{pasaje}
\vspace{3pt}\noindent{\small\textbf{p.\ 566} \quad \campo{\textquotedbl{}...y con cuadriláteros de transición\textquotedbl{} --- tesis/capitulos/05\_conclusiones.tex:47}}\par
\begin{pasaje}\footnotesize Combining elements of different order, say linear to quadratic elements, may be necessary to accomplish local mesh refinements. There are two ways to do this. One way is to use a transition element, which has different number of nodes on different sides [see Fig. 9.4.6(a)]. The other way is to impose a condition that constrains the midside node [...]

Figure 9.4.6 Combining different order elements: (a) use of a transition element that has three sides linear and one side quadratic; and (b) use of a linear constraint equation to connect a linear side to a quadratic side.\end{pasaje}
\vspace{3pt}\noindent{\small\textbf{p.\ 563} \quad \campo{\textquotedbl{}...ofreciendo así una variedad que enriquezca el estudio comparativo de la convergencia y del comportamiento bajo distorsión de malla\textquotedbl{} --- tesis/capitulos/05\_conclusiones.tex:47}}\par
\begin{pasaje}\footnotesize Figure 9.4.1 Finite elements with unacceptable vertex angles: (a) linear quadrilateral elements; (b) linear triangular elements; and (c) quadratic triangular elements. The angles marked are either too small compared with 0° or too large compared with 180°.

For higher-order Lagrange elements (also called the C\textasciicircum{}0 elements), the locations of the interior nodes contribute to the element distortion, and therefore they are constrained to lie within certain distance from the vertex nodes. For example, in the case of a quadratic element the midside node should be at a distance not less than one-fourth of the length of the side from the vertex nodes (see Fig. 9.4.2).\end{pasaje}
\vspace{6pt}
\subsection{\texttt{rutten2012learning} --- The learning effects of computer simulations in science education}
{\footnotesize\campo{Ejemplar:} Nico Rutten, Wouter R. van Joolingen, Jan T. van der Veen (University of Twente). Elsevier (Computers \& Education 58(1):136-153). 2012. \campo{Archivo:} \texttt{Rutten et al 2012 - The Learning Effects of Computer Simulations in Science Education.pdf}\par}\vspace{4pt}
\vspace{3pt}\noindent{\small\textbf{p.\ 136} \quad \campo{01\_introduccion (Justificacion) y 02\_marco\_teorico §1.2, primer principio: la revision de una decada de estudios cuasiexperimentales encuentra evidencia solida de que las simulaciones mejoran la instruccion tradicional, sobre todo en laboratorio, y que el efecto depende de como se visualiza la informacion y del apoyo...}}\par
\begin{pasaje}\footnotesize In particular, we consider the effects of variations in how information is visualized, how instructional support is provided, and how computer simulations are embedded within the lesson scenario. The reviewed literature provides robust evidence that computer simulations can enhance traditional instruction, especially as far as laboratory activities are concerned. However, in most of this research the use of computer\end{pasaje}
\vspace{3pt}\noindent{\small\textbf{p.\ 151} \quad \campo{02\_marco\_teorico §1.2, primer principio: entre los resultados destacados esta el gran efecto de aprendizaje de visualizar fenomenos invisibles. \textbackslash{}autocite[p.\textasciitilde{}151]\{rutten2012learning\}}}\par
\begin{pasaje}\footnotesize There is one outstanding result: the visualization of invisible phenomena that was investigated by Trey and Khan (2008). A large learning effect on these unobservable phenomena was found, which is in line with one of the main advantages of computer simulation hypothesized by de Jong and colleagues (van Berkum \& de Jong, 1991; van Joolingen \& de Jong, 1991). With respect to the level of immersion in 3D simulations, no\end{pasaje}
\vspace{6pt}
\subsection{\texttt{salari2000mms} --- Code Verification by the Method of Manufactured Solutions}
{\footnotesize\campo{Ejemplar:} Kambiz Salari and Patrick Knupp. Sandia National Laboratories (SANDIA REPORT SAND2000-1444, Unlimited Release; Sandia is a multiprogram laboratory operated by Sandia Corporation, a Lockheed Martin Company, para el U.S. Department of Energy, contrato DE-AC04-94AL85000). 2000 (portada: \textquotedbl{}Printed June 2000\textquotedbl{}). \campo{Archivo:} \texttt{Salari y Knupp 2000 - Code Verification by the Method of Manufactured Solutions.pdf}\par}\vspace{4pt}
\vspace{3pt}\noindent{\small\textbf{p.\ 21} \quad \campo{Descripcion del procedimiento MMS: \textquotedbl{}se elige a priori un campo de desplazamientos exacto u\_M, se sustituye en las ecuaciones de equilibrio para deducir la fuerza de cuerpo b = -div(sigma(u\_M)) que lo produce\textquotedbl{} --- tesis/capitulos/02\_marco\_teorico.tex:396}}\par
\begin{pasaje}\footnotesize A manufactured solution is an exact solution to some PDE or set of PDE's that has been constructed by solving the problem backwards. Suppose one is solving a differential equation of the form  Du = g  (4)  where D is the differential operator, u is the solution, and g is a source term. In the method of exact solution (MES), one chooses the function g and then, using methods from classical applied mathematics, inverts the operator to solve for u. In MMS, one first manufactures a solution u and then applies D to u to find g (a much simpler procedure). In the MMS code Verification procedure, one manufactures a solution to the fully general set of interior equations solved by the code that is to be Verified.\end{pasaje}
\vspace{3pt}\noindent{\small\textbf{p.\ 26} \quad \campo{\textquotedbl{}...y se inyectan esa carga y las restricciones de Dirichlet no homogeneas correspondientes\textquotedbl{} --- tesis/capitulos/02\_marco\_teorico.tex:396}}\par
\begin{pasaje}\footnotesize Assume for the moment that both the type and location of the boundary condition has been given. Then, in the MMS procedure the value attribute can be computed from the manufactured solution. [...] Dirichlet Boundary Conditions. If a Dirichlet condition is imposed on a given dependent variable for some portion of the domain boundary, then the value attribute is easily computed by evaluating the manufactured solution to the interior equations at the relevant points in space dictated by the location attribute and the grid on the boundary. The calculated values form part of the code input, along with the boundary condition type and location.\end{pasaje}
\vspace{3pt}\noindent{\small\textbf{p.\ 21 (mecanica) y 20 (rango de \textquotedbl{}tecnica...} \quad \campo{\textquotedbl{}El metodo de soluciones manufacturadas (MMS) constituye la tecnica de referencia para verificar la correccion de un codigo de elementos finitos, ya que permite construir una solucion exacta arbitraria y derivar de ella el termino fuente que la hace satisfacer la ecuacion de gobierno\textquotedbl{} --- tesis/capitulos/04\_resultados.tex:16}}\par
\begin{pasaje}\footnotesize [p.\textasciitilde{}21] In MMS, one first manufactures a solution u and then applies D to u to find g (a much simpler procedure). --- [p.\textasciitilde{}20] In the following description of MMS we advocate that MMS be associated with the order-of-accuracy criterion, resulting in the most comprehensive and rigorous of all code Verification methods.\end{pasaje}
\vspace{3pt}\noindent{\small\textbf{p.\ 28 y 29} \quad \campo{El error se cuantifica con normas de error globales y las normas deben decrecer al refinar h con la tasa teorica, lo que confirma la correcta implementacion --- tesis/capitulos/02\_marco\_teorico.tex:396-406 y 04\_resultados.tex:16}}\par
\begin{pasaje}\footnotesize [p.\textasciitilde{}28] To compute the discretization error, several measures are possible. Let x be a point in R\textasciicircum{}n and dx the local volume. To compare functions u and v on R\textasciicircum{}n the L2 norm of u-v is [Eq. 14] [...] Because the manufactured solution is defined on the continuum, one can evaluate the exact (manufactured) solution at the same locations in time and space as the discrete solution. --- [p.\textasciitilde{}29] Given that we can obtain some measure of the global error on a grid, we then run the code on a series of different mesh sizes h (at least 2) and compute the error for each mesh. From this information, the observed order-of-accuracy can be estimated and compared to the theoretical order-of-accuracy. [...] For consistent methods, the discretization error E in the solution is proportional to h\textasciicircum{}p where p\textgreater{}0 is the theoretical order of the method  E = Ch\textasciicircum{}p + H.O.T  (21)\end{pasaje}
\vspace{3pt}\noindent{\small\textbf{p.\ 29} \quad \campo{Se mide ademas el error del campo de tensiones recuperado porque las normas de desplazamiento no ejercitan la cadena de recuperacion de tensiones --- tesis/capitulos/02\_marco\_teorico.tex:406 (segunda mitad de la linea, actualmente SIN cita)}}\par
\begin{pasaje}\footnotesize Often engineering software computes not only the solution but certain derived variables such as flux or aerodynamic coefficients (lift, drag, moment). For thorough Verification of the code, one should also compute the error in these derived variables if they are provided as code output. For example, Verification of the calculated flux (which involves the gradient of u), can use the following global error measure [Eq. 19] [...] It should be noted that the theoretical order-of-accuracy of calculated fluxes may be less than the theoretical order-of-accuracy of the interior equations.\end{pasaje}
\vspace{3pt}\noindent{\small\textbf{p.\ 11 y 31} \quad \campo{Encuadre metodologico: \textquotedbl{}La evaluacion de la correccion se realizo mediante un esquema de verificacion y validacion: verificacion por el metodo de soluciones manufacturadas y validacion contra la viga de Timoshenko y la membrana de Cook\textquotedbl{} --- tesis/capitulos/02b\_diseno\_metodologico.tex:16; y \textquotedbl{}El marco metodologico de verificacion y...}}\par
\begin{pasaje}\footnotesize [p.\textasciitilde{}11] We argue in this report that a formal procedure known as the Method of Manufactured Solutions (MMS) can be applied to rigorously verify computer codes. [...] By this definition, code Verification is not concerned with asking the question, are we solving the right equation? The is part of code Validation. --- [p.\textasciitilde{}31] In this Section we summarize the MMS code Verification procedure (idealized in Figure 2 on page 36). Step 1. Determine Governing Equations.\end{pasaje}
\vspace{6pt}
\subsection{\texttt{sampieri2018metodologia} --- Metodología de la investigación: las rutas cuantitativa,...}
{\footnotesize\campo{Ejemplar:} Roberto Hernández-Sampieri; Christian Paulina Mendoza Torres. McGraw-Hill Interamericana Editores, S.A. de C.V.. 2018. \campo{Archivo:} \texttt{Hernandez-Sampieri y Mendoza 2018 - Metodologia de la investigacion.pdf}\par}\vspace{4pt}
\vspace{3pt}\noindent{\small\textbf{p.\ 6} \quad \campo{El trabajo tiene «enfoque cuantitativo en su fase de verificación y validación», entendido como establecer la calidad del artefacto «midiendo errores numéricos, tensiones, desplazamientos y tasas de convergencia frente a referencias de exactitud conocida» --- tesis/capitulos/02b\_diseno\_metodologico.tex:4 (cita conjunta...}}\par
\begin{pasaje}\footnotesize La ruta cuantitativa es apropiada cuando queremos estimar las magnitudes u ocurrencia de los fenómenos y probar hipótesis. [...] 4. Los datos se encuentran en forma de números (cantidades) y, por tanto, su recolección se fundamenta en la medición (en los casos se miden las variables contenidas en las hipótesis). Esta recolección se lleva a cabo utilizando procedimientos estandarizados y aceptados por una comunidad científica.\end{pasaje}
\vspace{3pt}\noindent{\small\textbf{p.\ 137} \quad \campo{«La tabla operacionaliza las variables del estudio, es decir, especifica para cada una las operaciones concretas con que se la mide» --- definición de operacionalización de variables. tesis/capitulos/02b\_diseno\_metodologico.tex:21 (sin localizador)}}\par
\begin{pasaje}\footnotesize Definición operacional. Una definición operacional consiste en el conjunto de procedimientos, técnicas y métodos para medir una variable en los casos de la investigación. Constituye las actividades que un observador debe realizar para recibir las impresiones sensoriales que indican la existencia de un concepto teórico en mayor o menor grado (MacGregor, 2006; Reynolds, 1986). La definición operacional nos señala: para recoger datos respecto de una variable es necesario hacer esto y aquello [...]\end{pasaje}
\vspace{3pt}\noindent{\small\textbf{p.\ 215} \quad \campo{El conjunto de casos de estudio «no es una muestra estadística de ese universo, sino una selección intencional por valor probatorio ---una muestra dirigida, no probabilística---» --- equivalencia entre muestra no probabilística y muestra dirigida, y su selección por características del estudio y no por criterio estadístico de...}}\par
\begin{pasaje}\footnotesize Muestras no probabilísticas. Las muestras no probabilísticas, también denominadas muestras dirigidas, como ya te mencionamos, suponen un procedimiento de selección orientado por las características y contexto de la investigación, más que por un criterio estadístico de generalización. [...] seleccionan individuos o casos típicos sin intentar que sean estadísticamente representativos de una población determinada. [...] La ventaja de una muestra no probabilística ---desde la visión cuantitativa--- es su utilidad para determinados diseños de estudio que requieren no tanto una representatividad de elementos de una población, sino una cuidadosa y controlada elección de casos con ciertas características especificadas previamente en el planteamiento del problema. Su valor reside en que las unidades de análisis son estudiadas a profundidad, lo que permite conocer el comportamiento de las variables...\end{pasaje}
\vspace{6pt}
\subsection{\texttt{stimpson2007verdict} --- The Verdict Geometric Quality Library}
{\footnotesize\campo{Ejemplar:} C. J. Stimpson; C. D. Ernst; P. Knupp; P. P. Pébay; D. Thompson (en la lista de distribución, p. 108, figuran los nombres completos: Clinton J. Stimpson, Corey D. Ernst, Patrick Knupp, Philippe P. Pébay, David C. Thompson). Sandia National Laboratories (SANDIA REPORT SAND2007-1751, Unlimited Release; operado por Sandia Corporation, a Lockheed Martin Company, para el U.S. DOE / NNSA, contrato DE-AC04-94-AL85000). 2007 (portada y portadilla: \textquotedbl{}Printed March 2007\textquotedbl{}; metadatos del PDF: creación 2007-03-20). \campo{Archivo:} \texttt{Stimpson et al 2007 - The Verdict Geometric Quality Library.pdf}\par}\vspace{4pt}
\vspace{3pt}\noindent{\small\textbf{p.\ 3} \quad \campo{«La precisión del resultado depende de la calidad geométrica de los elementos. EduFEM evalúa dos métricas ortogonales tomadas de la biblioteca de calidad geométrica Verdict» --- tesis/capitulos/02\_marco\_teorico.tex:362}}\par
\begin{pasaje}\footnotesize Verdict is a collection of subroutines for evaluating the geometric qualities of triangles, quadrilaterals, tetrahedra, and hexahedra using a variety of metrics. A metric is a real number assigned to one of these shapes depending on its particular vertex coordinates. These metrics are used to evaluate the input to finite element, finite volume, boundary element, and other types of solvers that approximate the solution to partial differential equations defined over regions of space. The geometric qualities of these regions is usually strongly tied to the accuracy these solvers are able to obtain in their approximations.\end{pasaje}
\vspace{3pt}\noindent{\small\textbf{p.\ 51} \quad \campo{El Jacobiano escalado «equivale al seno del ángulo local entre ambas tangentes y vive en el rango \$[-1,1]\$»; y «el módulo M0 emplea la forma por vértices de Verdict ---el seno del ángulo interior en cada esquina, tomando el peor---» --- tesis/capitulos/02\_marco\_teorico.tex:362-374}}\par
\begin{pasaje}\footnotesize 5.15 Scaled Jacobian --- The scaled Jacobian is the minimum of the Jacobian at each corner divided by the lengths of the 2 edge vectors (which is the minimum sine of the included angles): q = min\{ a0/(L0 L3), a1/(L1 L0), a2/(L2 L1), a3/(L3 L2) \}. Note that if any edge has \textbar{}L\textbar{} \textless{} DBL MIN, we take q = 0. // quadrilateral scaled Jacobian --- Dimension: 1 $\cdot$ Acceptable Range: [0.3, 1] $\cdot$ Normal Range: [-1, 1] $\cdot$ Full Range: [-1, 1] $\cdot$ q for unit square: 1 $\cdot$ Reference: [5] $\cdot$ Verdict function: v\_quad\_scaled\_jacobian\end{pasaje}
\vspace{3pt}\noindent{\small\textbf{p.\ 57} \quad \campo{«La compacidad o stretch mide la proporción del elemento relacionando su arista más corta \$L\_\{\textbackslash{}min\}\$ con su diagonal más larga \$D\_\{\textbackslash{}max\}\$: \$Q=\textbackslash{}sqrt\{2\}\textbackslash{},L\_\{\textbackslash{}min\}/D\_\{\textbackslash{}max\}\textbackslash{}in[0,1]\$, que vale \$1\$ para un cuadrado perfecto» --- tesis/capitulos/02\_marco\_teorico.tex:362-374 (ec. eq:stretch)}}\par
\begin{pasaje}\footnotesize 5.21 Stretch --- The stretch is q = sqrt(2) $\cdot$ min\_\{i in \{0,1,2,3\}\} \{L\_i\} / D\_max. Note that if D\_max \textless{} DBL MIN, we take q = DBL MAX. // quadrilateral stretch --- Dimension: 1 $\cdot$ Acceptable Range: [0.25, 1] $\cdot$ Normal Range: [0, 1] $\cdot$ Full Range: [0, DBL MAX] $\cdot$ q for unit square: 1 $\cdot$ Reference: [3] $\cdot$ Verdict function: v\_quad\_stretch\end{pasaje}
\vspace{3pt}\noindent{\small\textbf{p.\ 57 (stretch) y 51 (Jacobiano escalado)} \quad \campo{«De la documentación de Verdict provienen los mínimos admisibles, \$q\_\{SJ\}\textbackslash{}ge 0\{,\}50\$ y \$Q\textbackslash{}ge 0\{,\}25\$» --- tesis/capitulos/02\_marco\_teorico.tex:376}}\par
\begin{pasaje}\footnotesize [p. 57] quadrilateral stretch --- Acceptable Range: [0.25, 1]  \textbar{}\textbar{}  [p. 51] quadrilateral scaled Jacobian --- Acceptable Range: [0.3, 1]  \textbar{}\textbar{}  [p. 12, §1.4 Metric Ranges] \textquotedbl{}Acceptable Range : Well-behaved elements will have metrics in this range.\textquotedbl{}\end{pasaje}
\vspace{3pt}\noindent{\small\textbf{p.\ 35 y 36} \quad \campo{«Las dos restantes se construyen sobre una descomposición bilineal del cuadrilátero, equivalente a la que Verdict emplea para sus métricas de forma», con \$\textbackslash{}bm\{X\}\_1=\textbackslash{}tfrac14(-\textbackslash{}bm r\_1+\textbackslash{}bm r\_2+\textbackslash{}bm r\_3-\textbackslash{}bm r\_4)\$, \$\textbackslash{}bm X\_2=\textbackslash{}tfrac14(-\textbackslash{}bm r\_1-\textbackslash{}bm r\_2+\textbackslash{}bm r\_3+\textbackslash{}bm r\_4)\$, \$\textbackslash{}bm X\_3=\textbackslash{}tfrac14(\textbackslash{}bm r\_1-\textbackslash{}bm r\_2+\textbackslash{}bm r\_3-\textbackslash{}bm r\_4)\$ ---...}}\par
\begin{pasaje}\footnotesize [p. 35] \textquotedbl{}The principal axes are  X1 = (P1 - P0) + (P2 - P3),  X2 = (P2 - P1) + (P3 - P0)   (7)\textquotedbl{}  \textbar{}\textbar{}  [p. 36] \textquotedbl{}and the cross derivatives of the map from parametric to world space are oriented along  X12 = (P0 - P1) + (P2 - P3) = X21 = (P0 - P3) + (P2 - P1).   (8)\textquotedbl{}\end{pasaje}
\vspace{3pt}\noindent{\small\textbf{p.\ 58} \quad \campo{«el estrechamiento (taper) \$T\_R\$, que es nulo en todo paralelogramo y crece con la trapecialidad» --- tesis/capitulos/02\_marco\_teorico.tex:376 (parrafo posterior a eq:descomposicion-bilineal)}}\par
\begin{pasaje}\footnotesize 5.22 Taper --- Taper is the maximum ratio of cross derivative magnitude to principal axis magnitude: q = \textbar{}\textbar{}X12\textbar{}\textbar{} / min(\textbar{}\textbar{}X1\textbar{}\textbar{}, \textbar{}\textbar{}X2\textbar{}\textbar{}). Note that if \textbar{}\textbar{}X1\textbar{}\textbar{} or \textbar{}\textbar{}X2\textbar{}\textbar{} \textless{} DBL MIN, we set q = DBL MAX. // quadrilateral taper --- Dimension: 1 $\cdot$ Acceptable Range: [0, 0.7] $\cdot$ Normal Range: [0, DBL MAX] $\cdot$ Full Range: [0, DBL MAX] $\cdot$ q for unit square: 0 $\cdot$ Reference: Adapted from [10] $\cdot$ Verdict function: v\_quad\_taper\end{pasaje}
\vspace{3pt}\noindent{\small\textbf{p.\ 57 y 51} \quad \campo{«El módulo M0 atiende la calidad de malla en el pre-proceso, evaluando el Jacobiano escalado y la compacidad (stretch) con los umbrales bibliográficos de Verdict» --- tesis/capitulos/04\_resultados.tex:209}}\par
\begin{pasaje}\footnotesize [p. 57] quadrilateral stretch --- Acceptable Range: [0.25, 1]  \textbar{}\textbar{}  [p. 51] quadrilateral scaled Jacobian --- Acceptable Range: [0.3, 1]\end{pasaje}
\vspace{6pt}
\subsection{\texttt{strang2008analysis} --- An Analysis of the Finite Element Method}
{\footnotesize\campo{Ejemplar:} Gilbert Strang (Massachusetts Institute of Technology) y George J. Fix (University of Maryland). Prentice-Hall, Inc. (serie: Prentice-Hall Series in Automatic Computation, editor George Forsythe). 1973. \campo{Archivo:} \texttt{Strang y Fix 1973 - An Analysis of the Finite Element Method.pdf}\par}\vspace{4pt}
\vspace{3pt}\noindent{\small\textbf{p.\ 107 (edición Prentice-Hall 1973; PDF...} \quad \campo{Que la teoría de aproximación del MEF establece las tasas asintóticas O(h\textasciicircum{}\{p+1\}) en la norma L2 y O(h\textasciicircum{}p) en H1, con p el orden polinómico del elemento --- tesis/capitulos/02\_marco\_teorico.tex:406 (\textbackslash{}autocite\{strang2008analysis,hughes2000fem\}, sin localizador de página)}}\par
\begin{pasaje}\footnotesize Convergence in strain energy is essentially convergence of the mth derivatives of u\textasciicircum{}h to those of u. This derivative is therefore special, because the Ritz method is minimizing in energy. For the sth derivative, where s may be smaller or larger than m, convergence cannot be faster than O(h\textasciicircum{}\{k-s\}); this is the order of the best approximation to u from a space S\textasciicircum{}h of degree k - 1. This rate of convergence will normally be attained by the Ritz solution u\textasciicircum{}h. Using the \textquotedbl{}Nitsche trick,\textquotedbl{} which gave the O(h\textasciicircum{}2) rate in displacement for one-dimensional linear elements in Section 1.6, we shall show that
(4)     \textbar{}\textbar{}u - u\textasciicircum{}h\textbar{}\textbar{}\_s = O(h\textasciicircum{}\{k-s\} + h\textasciicircum{}\{2(k-m)\}).
In almost all cases the first exponent is smaller, and approximation theory governs the convergence rate.\end{pasaje}
\vspace{3pt}\noindent{\small\textbf{p.\ 106 (edición Prentice-Hall 1973; PDF...} \quad \campo{Idem, replicada en tesis/capitulos/04\_resultados.tex:16 --- \textquotedbl{}para un elemento de orden polinómico p, se espera O(h\textasciicircum{}\{p+1\}) en la norma L2 y O(h\textasciicircum{}p) en la seminorma H1\textquotedbl{} (\textbackslash{}autocite\{strang2008analysis,hughes2000fem,zienkiewicz2013fem\}, sin localizador)}}\par
\begin{pasaje}\footnotesize The basic hypotheses will be, first, that S\textasciicircum{}h is of degree k - 1 -- it contains in each element the complete polynomial of this degree, restricted only at the boundary by essential conditions -- and second, that its basis is uniform as h -\textgreater{} 0. The latter is effectively a geometrical condition on the element regions: If the diameter of e\_i is h\_i, then e\_i should contain a sphere of radius at least tau*h\_i, where tau is bounded away from zero. This forbids arbitrarily small angles in triangular elements. Angles near pi must also be forbidden in quadrilaterals. Under these conditions, the distance between u and S\textasciicircum{}h is
     a(u - u\textasciicircum{}h, u - u\textasciicircum{}h) \textless{}= C\textasciicircum{}2 h\textasciicircum{}\{2(k-m)\} \textbar{}u\textbar{}\_k\textasciicircum{}2.
Therefore, h\textasciicircum{}\{2(k-m)\} is the rate of convergence in strain energy.\end{pasaje}
\vspace{3pt}\noindent{\small\textbf{p.\ 144 (edición Prentice-Hall 1973; PDF...} \quad \campo{Respaldo secundario del mismo enunciado: la estimación de aproximación en normas de media cuadrática que sustenta las tasas h\textasciicircum{}\{k-s\} (Teorema 3.3)}}\par
\begin{pasaje}\footnotesize THEOREM 3.3
Suppose that S\textasciicircum{}h is of degree k - 1, and its basis is uniform to order q. Suppose also that the derivatives D\_j associated with the nodal parameters are all of order \textbar{}D\_j\textbar{} \textless{} k - n/2. Then for any function u(x\_1, ..., x\_n) which has k derivatives in the mean-square sense, and any derivative D\textasciicircum{}alpha of order s \textless{}= q,
(7)     integral\_\{e\_i\} \textbar{}D\textasciicircum{}alpha u(x) - D\textasciicircum{}alpha u\_I(x)\textbar{}\textasciicircum{}2 dx \textless{}= C\_s\textasciicircum{}2 h\_i\textasciicircum{}\{2(k-s)\} \textbar{}u\textbar{}\textasciicircum{}2\_\{k,e\_i\}.
Since the integral over Omega is the sum of the integrals over e\_i,
(8)     \textbar{}u - u\_I\textbar{}\_\{s,Omega\} \textless{}= C\_s h\textasciicircum{}\{k-s\} \textbar{}u\textbar{}\_\{k,Omega\}.\end{pasaje}
\vspace{6pt}
\subsection{\texttt{suarez1998edelas2d} --- ED-Elas2D: An Educational Program for Computer-aided Training in...}
{\footnotesize\campo{Ejemplar:} Benjamín Suárez, Elena Blanco, Lluís Gil, Pere Zapata y Eugenio Oñate. Taylor \& Francis (impreso solo en la hoja de cubierta del distribuidor como Publisher; el artículo en sí no lleva pie editorial). Publicado en European Journal of Engineering Education, Vol. 23, No. 2.. 1998. \campo{Archivo:} \texttt{Suarez et al 1998 - ED-Elas2D.pdf}\par}\vspace{4pt}
\vspace{3pt}\noindent{\small\textbf{p.\ 243} \quad \campo{Bishay se inscribe «en la línea de programas previos como ED-Beams y ED-Frames» (capitulos/02\_marco\_teorico.tex:11)}}\par
\begin{pasaje}\footnotesize The program follows the methodology proposed by the authors in the development of the educational programs ED-Beams and ED-Frames [2, 3], designed for the computer-aided training in linear elastic frames using matrix methods.\end{pasaje}
\vspace{3pt}\noindent{\small\textbf{p.\ 243} \quad \campo{ED-Elas2D es «un programa educativo desarrollado en la Universitat Politècnica de Catalunya para el adiestramiento asistido por computadora en el análisis de sólidos bidimensionales por el MEF» (02\_marco\_teorico.tex:13)}}\par
\begin{pasaje}\footnotesize SUMMARY In this paper, a methodology for computer-aided training of structural problems by the finite element method is presented. The resulting educational program, named ED-Elas2D, has been developed by a group of researchers at the School of Civil Engineering at the Technical University of Catalonia (UPC) in Barcelona, UPC.\end{pasaje}
\vspace{3pt}\noindent{\small\textbf{p.\ 244} \quad \campo{ED-Elas2D comparte con EduFEM el dominio: «elasticidad plana en tensión y deformación plana» (02\_marco\_teorico.tex:13)}}\par
\begin{pasaje}\footnotesize The topic chosen to initiate the study of FEM was the linear elastic analysis of structures under plain stress or plain strain conditions, i.e. structures satisfying the assumptions of 2D elasticity.\end{pasaje}
\vspace{3pt}\noindent{\small\textbf{p.\ 253} \quad \campo{ED-Elas2D comparte «la organización en pre-proceso, proceso y post-proceso» (02\_marco\_teorico.tex:13) y la tabla comparativa le asigna «Media (pre/proc/post)» (02\_marco\_teorico.tex:28)}}\par
\begin{pasaje}\footnotesize The organization of the program in three modular blocks, pre-processing, processing and post-processing, has also proved to be adequate as it follows the three logic steps in the practical solution of any structural problem by FEM.\end{pasaje}
\vspace{3pt}\noindent{\small\textbf{p.\ 244} \quad \campo{«una biblioteca de elementos triangulares y cuadriláteros que incluye los elementos Q4, Q8 y Q9» (02\_marco\_teorico.tex:13) y la celda de tabla «Triángulos y cuadriláteros (Q4/Q8/Q9)» (02\_marco\_teorico.tex:28)}}\par
\begin{pasaje}\footnotesize FIG. 1. Definition of a new problem. Element types accepted by the program. [Lista desplegable visible en la captura: «3-Node triangle / 6-Node triangle / 4-node quadrilateral / 8-node quadrilateral / 9-node quadrilateral»]\end{pasaje}
\vspace{3pt}\noindent{\small\textbf{p.\ 247} \quad \campo{«aquel programa de finales de la década de 1990 [...] ligado a la plataforma y a las posibilidades de su época» (02\_marco\_teorico.tex:13)}}\par
\begin{pasaje}\footnotesize All these considerations have been taken into account in the development of the educational environment ED-Elas2D. This educational program has been developed in the Windows environment and includes three distinct blocks which are a logical consequence of the basic steps to follow for the analysis of a structure with FEM.\end{pasaje}
\vspace{3pt}\noindent{\small\textbf{p.\ 248} \quad \campo{Columna «Canal completo: Sí» para ED-Elas2D, entendido como recorrer mapeo isoparamétrico, Jacobiano, B, constitutiva, rigidez, ensamblaje y recuperación de tensiones (02\_marco\_teorico.tex:28)}}\par
\begin{pasaje}\footnotesize FIG. 9. Help menu in the beginner level showing the steps to solve a problem. [Contenido de la captura: «The program's calculation diagram is the following: 1- Calculation and assembly of the stiffness matrix for each element. 2- Calculation and assembly of the equivalent nodal force vector for each element. 3- Boundary conditions on the prescribed nodes. 4- Solution of the linear system. 5- Calculation of the strain vector for each element. 6- Calculation of the stress vector for each element. 7- Calculation of the reactions on the prescribed nodes.»]\end{pasaje}
\vspace{3pt}\noindent{\small\textbf{p.\ 246} \quad \campo{ED-Elas2D está «disponible en español» y comparte con EduFEM «el idioma»; celda de tabla «Idioma: Español» (02\_marco\_teorico.tex:13 y :28)}}\par
\begin{pasaje}\footnotesize FIG. 6. List of theoretical lessons. [En la captura, la barra de menús de la ventana «ED-Elas2D Processing - Help» está en español: «Archivo  Edición  Marcador  Opciones  Ayuda» y «Contenido \textbar{} Búsqueda \textbar{} Atrás \textbar{} Imprimir». El contenido del índice, en cambio, está en inglés: «Index of theory», «Introduction», «The continuous problem», «Kinematic relations», «Constitutive relation»...]\end{pasaje}
\vspace{3pt}\noindent{\small\textbf{p.\ 254} \quad \campo{«software académico cerrado»; celda de tabla «Licencia: Académico» (02\_marco\_teorico.tex:13, :28 y :30)}}\par
\begin{pasaje}\footnotesize Acknowledgements --- The development of ED-Elas2D has been partially funded by the EC programmes COMETT and Leonardo da Vinci. The authors are also grateful to the International Centre for Numerical Methods in Engineering of Barcelona for the support in the development of this project. [Y en REFERENCES: «[1] ED-ElAs2D (1996) Educational Program for the Analysis of 2D Solids Using the Finite Element Method (Barcelona, CIMNE).»]\end{pasaje}
\vspace{3pt}\noindent{\small\textbf{p.\ 246} \quad \campo{01\_introduccion (Planteamiento): con los programas comerciales el estudiante queda limitado a introducir los datos e interpretar los resultados, sin acceso a los detalles del analisis. \textbackslash{}autocite[p.\textasciitilde{}246]\{suarez1998edelas2d\}}}\par
\begin{pasaje}\footnotesize On the other hand, there are many FEM commercial programs which can be used as a base for computer tutorials. The main difficulty of these programs is the lack of transparency to be used as a fully teacher environment. In many cases, these commercial programs operate like a blackbox and therefore the student is practically limited to the introduction of the problem data and the interpretation of the results and has no access to the details of the analysis.\end{pasaje}
\vspace{3pt}\noindent{\small\textbf{p.\ 253} \quad \campo{02\_marco\_teorico §1.1 y tab:comparativa: los propios autores reconocen que el post-proceso de ED-Elas2D no explica conceptos clave como el suavizado nodal y es una herramienta meramente descriptiva. \textbackslash{}autocite[p.\textasciitilde{}253]\{suarez1998edelas2d\}}}\par
\begin{pasaje}\footnotesize The post-processing module, however, does not explain key concepts such as the nodal smoothing of stresses and strains. It is therefore a mere descriptive tool which simply provides visual information on the main results of the structural analysis.\end{pasaje}
\vspace{6pt}
\subsection{\texttt{timoshenko1970elasticity} --- Theory of Elasticity}
{\footnotesize\campo{Ejemplar:} Timoshenko, S. P. y Goodier, J. N. (portada: \textquotedbl{}By S. TIMOSHENKO And J. N. GOODIER\textquotedbl{}, Professors of Engineering Mechanics, Stanford University). McGraw-Hill Book Company, Inc.. 1951. \campo{Archivo:} \texttt{Timoshenko y Goodier 1951 - Theory of Elasticity.pdf}\par}\vspace{4pt}
\vspace{3pt}\noindent{\small\textbf{p.\ 11} \quad \campo{Delimitacion del alcance: EduFEM resuelve las dos idealizaciones clasicas, tension plana (cuerpos delgados cargados en su plano) y deformacion plana (cuerpos prismaticos largos con seccion y carga invariantes a lo largo del eje longitudinal). tesis/capitulos/01\_introduccion.tex:62}}\par
\begin{pasaje}\footnotesize 7. Plane Stress. If a thin plate is loaded by forces applied at the boundary, parallel to the plane of the plate and distributed uniformly over the thickness (Fig. 8), the stress components sigma\_z, tau\_xz, tau\_yz are zero on both faces of the plate, and it may be assumed, tentatively, that they are zero also within the plate. The state of stress is then specified by sigma\_x, sigma\_y, tau\_xy only, and is called plane stress. [...] 8. Plane Strain. A similar simplification is possible at the other extreme when the dimension of the body in the z-direction is very large. If a long cylindrical or prismatical body is loaded by forces which are perpendicular to the longitudinal elements and do not vary along the length, it may be assumed that all cross sections are in the same condition.\end{pasaje}
\vspace{3pt}\noindent{\small\textbf{p.\ 11-12} \quad \campo{La tension plana modela cuerpos delgados cargados en su plano, en los que la tension normal al plano es despreciable (sigma\_z = 0); la deformacion plana modela cuerpos prismaticos largos en los que la deformacion longitudinal se anula (epsilon\_z = 0); corresponde, por ejemplo, a una presa o un tubo largo donde una rebanada...}}\par
\begin{pasaje}\footnotesize There are many important problems of this kind-a retaining wall with lateral pressure (Fig. 9), a culvert or tunnel (Fig. 10), a cylindrical tube with internal pressure, a cylindrical roller compressed by forces in a diametral plane as in a roller bearing (Fig. 11). In each case of course the loading must not vary along the length. Since conditions are the same at all cross sections, it is sufficient to consider only a slice between two sections unit distance apart. [...] Since the longitudinal displacement w is zero, Eqs. (2) give gamma\_xz = 0, gamma\_yz = 0, e\_z = dw/dz = 0\end{pasaje}
\vspace{3pt}\noindent{\small\textbf{p.\ 12} \quad \campo{Ecuacion eq:principales-vm, tercera linea: sigma\_z = 0 en tension plana y sigma\_z = nu (sigma\_x + sigma\_y) en deformacion plana. tesis/capitulos/02\_marco\_teorico.tex:298-312}}\par
\begin{pasaje}\footnotesize The longitudinal normal stress sigma\_z can be found in terms of sigma\_x and sigma\_y by means of Hooke's law, Eqs. (3). Since e\_z = 0 we find sigma\_z - nu (sigma\_x + sigma\_y) = 0\end{pasaje}
\vspace{3pt}\noindent{\small\textbf{p.\ 14 (tensiones principales) y 149...} \quad \campo{A partir de las componentes cartesianas se obtienen las tensiones principales y la equivalente de von Mises, la medida escalar con que habitualmente se resume un estado tensional biaxial. tesis/capitulos/02\_marco\_teorico.tex:298}}\par
\begin{pasaje}\footnotesize p. 14: \textquotedbl{}It may be seen that the angle alpha can be chosen in such a manner that the shearing stress tau becomes equal to zero. [...] From this equation two perpendicular directions can be found for which the shearing stress is zero. These directions are called principal directions and the corresponding normal stresses principal stresses.\textquotedbl{} \textbar{} p. 149: \textquotedbl{}In order to bring this theory into agreement with the fact that isotropic materials can sustain very large hydrostatic pressures without yielding, it has been proposed to split the strain energy into two parts, one due to the change in volume and the other due to the distortion, and consider only the second part in determining the strength. [...] we can present the part of the total energy due to distortion in the form [Eq. (88)]\textquotedbl{} (nota 2 al pie: \textquotedbl{}M. T. Huber, Czasopismo techniczne, Lwow, 1904. See also R. v. Mises, Gottingen Nachrichten,...\end{pasaje}
\vspace{3pt}\noindent{\small\textbf{p.\ 39-43 (art. 21); sigma\_x ec. (33) en p....} \quad \campo{La viga de Timoshenko simplemente apoyada (bajo carga uniforme) admite solucion analitica de la teoria de la elasticidad, y EduFEM se contrasta contra ella en sigma\_x, en la flecha y en la tension cortante. tesis/capitulos/02\_marco\_teorico.tex:408, tesis/capitulos/04\_resultados.tex:99 y tesis/capitulos/05\_conclusiones.tex:19}}\par
\begin{pasaje}\footnotesize p. 39-40: \textquotedbl{}21. Bending of a Beam by Uniform Load. Let a beam of narrow rectangular cross section of unit width, supported at the ends, be bent by a uniformly distributed load of intensity q, as shown in Fig. 28.\textquotedbl{} \textbar{} p. 41, tras la ec. (33): \textquotedbl{}The first term in this expression represents the stresses given by the usual elementary theory of bending, and the second term gives the necessary correction. This correction does not depend on x and is small in comparison with the maximum bending stress, provided the span of the beam is large in comparison with its depth.\textquotedbl{} \textbar{} p. 41, ec. (d), tercera linea: \textquotedbl{}tau\_xy = -(3q/4c\textasciicircum{}3)(c\textasciicircum{}2 - y\textasciicircum{}2)x = -(q/2I)(c\textasciicircum{}2 - y\textasciicircum{}2)x\textquotedbl{} \textbar{} p. 43: \textquotedbl{}Assuming that the deflection is zero at the ends (x = +-l) of the center line, we find [ec. (34)] The factor before the brackets is the deflection which is derived by the elementary analysis, assuming that cross sections of the beam...\end{pasaje}
\vspace{6pt}
\subsection{\texttt{virtanen2020scipy} --- SciPy 1.0: fundamental algorithms for scientific computing in...}
{\footnotesize\campo{Ejemplar:} Pauli Virtanen, Ralf Gommers, Travis E. Oliphant, Matt Haberland, Tyler Reddy, David Cournapeau, Evgeni Burovski, Pearu Peterson, Warren Weckesser, Jonathan Bright, Stefan J. van der Walt, Matthew Brett, Joshua Wilson, K. Jarrod Millman, Nikolay Mayorov, Andrew R. J. Nelson, Eric Jones, Robert Kern, Eric Larson, C J Carey, Ilhan Polat, Yu Feng, Eric W. Moore, Jake VanderPlas, Denis Laxalde, Josef Perktold, Robert Cimrman, Ian Henriksen, E. A. Quintero, Charles R. Harris, Anne M. Archibald, Antonio H. Ribeiro, Fabian Pedregosa, Paul van Mulbregt y \textquotedbl{}SciPy 1.0 Contributors\textquotedbl{} (autor colectivo n.o 39 en la firma impresa). Springer Nature (revista Nature Methods; la unica mencion impresa de editorial es la \textquotedbl{}Publisher's note Springer Nature...\textquotedbl{} en p. 271). No figura pie de imprenta clasico.. 2020. \campo{Archivo:} \texttt{Virtanen et al 2020 - SciPy 1.0.pdf}\par}\vspace{4pt}
\vspace{3pt}\noindent{\small\textbf{p.\ 261} \quad \campo{El motor se construye con \textquotedbl{}bibliotecas numericas de codigo abierto\textquotedbl{} que permiten exhibir matrices y operaciones intermedias sin las restricciones de un producto cerrado (encuadre de la decision metodologica). tesis/capitulos/01\_introduccion.tex:25}}\par
\begin{pasaje}\footnotesize SciPy is an open-source scientific computing library for the Python programming language. Since its initial release in 2001, SciPy has become a de facto standard for leveraging scientific algorithms in Python, with over 600 unique code contributors, thousands of dependent packages, over 100,000 dependent repositories and millions of downloads per year.\end{pasaje}
\vspace{3pt}\noindent{\small\textbf{p.\ 265} \quad \campo{EduFEM construye la matriz de rigidez dispersa en formato de coordenadas (COO) y la convierte a filas comprimidas (CSR) para el solucionador, apoyandose en las bibliotecas cientificas de Python. tesis/capitulos/02\_marco\_teorico.tex:270}}\par
\begin{pasaje}\footnotesize Data structures. Sparse matrices. scipy.sparse offers seven sparse matrix data structures, also known as sparse formats. The most important ones are the row- and column-compressed formats (CSR and CSC, respectively). These offer fast major-axis indexing and fast matrix-vector multiplication, and are used heavily throughout SciPy and dependent packages. [...] Iterating over and slicing of CSC and CSR matrices is now up to 35\% faster, and the speed of coordinate (COO)/diagonal (DIA) to CSR/CSC matrix format conversions has increased.\end{pasaje}
\vspace{3pt}\noindent{\small\textbf{p.\ 263} \quad \campo{El sistema reducido se resuelve por factorizacion LU directa sobre la matriz dispersa. tesis/capitulos/02\_marco\_teorico.tex:296}}\par
\begin{pasaje}\footnotesize sparse --- This subpackage includes implementations of several representations of sparse matrices. scipy.sparse.linalg provides a collection of linear algebra routines that work with sparse matrices, including linear equation solvers, eigenvalue decomposition, singular value decomposition and LU factorization.\end{pasaje}
\vspace{3pt}\noindent{\small\textbf{p.\ AUSENTE} \quad \campo{El sistema se resuelve \textquotedbl{}con un ordenamiento de columnas de minimo grado que aprovecha la simetria de la matriz para reducir el llenado (fill-in)\textquotedbl{}. tesis/capitulos/02\_marco\_teorico.tex:296}}\par
\begin{pasaje}\footnotesize (no hay pasaje: la busqueda de 'minimum degree', 'COLAMD', 'fill-in', 'permutation' y 'reorder' sobre las 15 paginas del PDF no devuelve ninguna ocurrencia)\end{pasaje}
\vspace{3pt}\noindent{\small\textbf{p.\ 263} \quad \campo{NumPy y SciPy \textquotedbl{}proveen el algebra densa y dispersa, la cuadratura y el solucionador lineal del metodo\textquotedbl{}. tesis/capitulos/03\_diseno\_implementacion.tex:25}}\par
\begin{pasaje}\footnotesize linalg --- Linear algebra functions, including elementary functions of a matrix, such as the trace, determinant, norm and condition number; basic solver for Ax = b; specialized solvers for Toeplitz matrices, circulant matrices, triangular matrices and other structured matrices; least-squares solver and pseudo-inverse calculations; eigenvalue and eigenvector calculations (basic and generalized); matrix decompositions, including Cholesky, Schur, Hessenberg, LU, LDLT, QR, QZ, singular value and polar; and functions to create specialized matrices [...]\end{pasaje}
\vspace{3pt}\noindent{\small\textbf{p.\ 261} \quad \campo{\textquotedbl{}La eleccion de un entorno integro en Python sobre las bibliotecas cientificas de uso extendido favorece la inspeccion y extension del codigo por parte de docentes y estudiantes.\textquotedbl{} tesis/capitulos/04\_resultados.tex:223}}\par
\begin{pasaje}\footnotesize SciPy is an open-source scientific computing library for the Python programming language. Since its initial release in 2001, SciPy has become a de facto standard for leveraging scientific algorithms in Python, with over 600 unique code contributors, thousands of dependent packages, over 100,000 dependent repositories and millions of downloads per year.\end{pasaje}
\vspace{6pt}
\subsection{\texttt{zienkiewicz2013fem} --- The Finite Element Method: Its Basis and Fundamentals}
{\footnotesize\campo{Ejemplar:} O. C. Zienkiewicz; R. L. Taylor; J. Z. Zhu. Elsevier Butterworth-Heinemann (publicado con la cooperacion de CIMNE, Barcelona). 2005. \campo{Archivo:} \texttt{Zienkiewicz Taylor y Zhu 2005 - The Finite Element Method.pdf}\par}\vspace{4pt}
\vspace{3pt}\noindent{\small\textbf{p.\ portada y creditos (PDF 4 y 5)} \quad \campo{ADVERTENCIA GLOBAL SOBRE LAS PAGINAS DE ESTE INFORME (afecta a todas las filas). El .bib declara edicion 7 (2013), ISBN 978-1-85617-633-0; el PDF verificado es la 6.a edicion (portada: 'Sixth edition'; creditos p. PDF 5: 'Sixth edition 2005'; ISBN 0 7506 6320 0; Elsevier Butterworth-Heinemann, con la cooperacion de CIMNE)....}}\par
\begin{pasaje}\footnotesize The Finite Element Method: Its Basis and Fundamentals / Sixth edition / O.C. Zienkiewicz, CBE, FRS [...] R.L. Taylor [...] J.Z. Zhu\textquotedbl{} (PDF 4) y \textquotedbl{}First published in 1967 by McGraw-Hill / Fifth edition published by Butterworth-Heinemann 2000 / Reprinted 2002 / Sixth edition 2005 [...] ISBN 0 7506 6320 0 / Published with the cooperation of CIMNE, the International Centre for Numerical Methods in Engineering, Barcelona, Spain\end{pasaje}
\vspace{3pt}\noindent{\small\textbf{p.\ 69-71 (Sec. 3.4 'Virtual work as the...} \quad \campo{\textquotedbl{}Esto conduce al principio de los trabajos virtuales\textquotedbl{} como forma debil obtenida al multiplicar la ecuacion de equilibrio por un campo de desplazamientos virtuales admisible e integrar por partes mediante el teorema de la divergencia (02\_marco\_teorico.tex:49; cita conjunta con hughes2000fem y reddy2006introduction).}}\par
\begin{pasaje}\footnotesize We shall therefore show in this section that the virtual work statement is simply a 'weak form' of equilibrium equations. [...] To obtain a weak form we shall proceed as before, introducing an arbitrary weighting function vector defined as v = du = [du, dv, dw]\textasciicircum{}T [...] Integrating each term by parts and rearranging we can write this as [...] where t = \{tx, ty, tz\} [...] are tractions acting per unit area of external boundary surface [...] [in (3.43) Green's formulae of Appendix G are again used]. [...] we can write Eq. (3.43) simply as INT\_Omega deps\textasciicircum{}T sigma dOmega - INT\_Omega du\textasciicircum{}T b dOmega - INT\_Gamma du\textasciicircum{}T t dGamma = 0 (3.46) [...] We see from the above that the virtual work statement is precisely the weak form of equilibrium equations and is valid for non-linear as well as linear stress-strain (or stress-strain rate) relations.\end{pasaje}
\vspace{3pt}\noindent{\small\textbf{p.\ 146-147 (Secs. 5.5.1 y 5.5.2)} \quad \campo{\textquotedbl{}El integrando no es polinomico: como B depende de J\textasciicircum{}\{-1\}, que es racional en (xi,eta), la integral carece de primitiva elemental\textquotedbl{} y por eso se aproxima con cuadratura de Gauss-Legendre (02\_marco\_teorico.tex:235; cita conjunta con hughes2000fem).}}\par
\begin{pasaje}\footnotesize To find now the global derivatives we invert J and write \{dNa/dx, dNa/dy, dNa/dz\} = J\textasciicircum{}\{-1\} \{dNa/dxi, dNa/deta, dNa/dzeta\} (5.10) [...] More explicitly we can write this as INT\_Omega G(x,y,z) dOmega = INT\_\{-1\}\textasciicircum{}\{1\} INT\_\{-1\}\textasciicircum{}\{1\} INT\_\{-1\}\textasciicircum{}\{1\} Gbar(xi,eta,zeta) J(xi,eta,zeta) dxi deta dzeta (5.13) [...] While the limits of integration are simple in the above case, unfortunately the explicit form of Gbar is not. Apart from the simplest elements, algebraic integration usually defies our mathematical skill, and numerical integration has to be used.\end{pasaje}
\vspace{3pt}\noindent{\small\textbf{p.\ 161 (Sec. 5.9.2 'Gauss quadrature')} \quad \campo{\textquotedbl{}La regla de m puntos integra exactamente polinomios de grado 2m-1\textquotedbl{} (02\_marco\_teorico.tex:242; cita conjunta con hughes2000fem).}}\par
\begin{pasaje}\footnotesize it is easy to see that for n sampling points we have 2n unknowns (wi and xi\_i) and hence a polynomial of degree 2n-1 could be constructed and exactly integrated [Fig. 5.16(b)]. The error is thus of order O(h\textasciicircum{}\{2n\}).\end{pasaje}
\vspace{3pt}\noindent{\small\textbf{p.\ 162 (Tabla 5.2, 'Gaussian quadrature...} \quad \campo{\textquotedbl{}cuadratura 2x2 para Q4 (cuatro puntos en xi = +-1/sqrt(3)) y 3x3 para Q9 (puntos en xi = +-sqrt(3/5), 0 con pesos 5/9, 8/9, 5/9)\textquotedbl{} (02\_marco\_teorico.tex:242).}}\par
\begin{pasaje}\footnotesize Table 5.2 Gaussian quadrature abscissae and weights for INT\_\{-1\}\textasciicircum{}\{1\} f(x)dx = SUM\_j f(xi\_j) w\_j.   +-xi\_j / w\_j :  n = 1 : 0 , 2.000000000000000 ; n = 2 : 1/sqrt(3) , 1.000000000000000 ; n = 3 : sqrt(0.6) , 5/9 ; 0.000000000000000 , 8/9\end{pasaje}
\vspace{3pt}\noindent{\small\textbf{p.\ 461, 464 y 465 (Sec. 13.2...} \quad \campo{\textquotedbl{}En el elemento Q4 los puntos de la regla 2x2 coinciden con los puntos de Barlow, ubicaciones donde las tensiones convergen con un orden superior al del resto del elemento (superconvergencia)\textquotedbl{} (02\_marco\_teorico.tex:315; repetido en 04\_resultados.tex:58 como explicacion del orden O(h\textasciicircum{}\{1,54\}) del Q4). Cita conjunta con...}}\par
\begin{pasaje}\footnotesize p. 461: \textquotedbl{}The stresses, or gradients, in this case will be optimal at points which correspond to the two Gauss quadrature points for each element as indicated in Fig. 13.4(b). This fact was observed experimentally by Barlow.1 We shall now state in an axiomatic manner that: (a) the displacements are best sampled at the nodes of the element, whatever the order of element used, and (b) the best accuracy for gradients or stresses is obtained at the Gauss points corresponding to the order of polynomial used in the solution.\textquotedbl{}  \textbar{}\textbar{}  p. 462: \textquotedbl{}At such points the order of the convergence of the function or its gradients is at least one order higher than that which would be anticipated from the appropriate polynomial and thus such points are known as superconvergent.\textquotedbl{}  \textbar{}\textbar{}  p. 464: \textquotedbl{}The single point at the centre of an element integrates precisely all linear functions passing through that point and,...\end{pasaje}
\vspace{3pt}\noindent{\small\textbf{p.\ 464-465 (texto y Fig. 13.6) y 467 (Fig....} \quad \campo{\textquotedbl{}en el Q9 los puntos optimos son tambien los de la regla 2x2, y no los nueve de la regla 3x3 con que se integra su rigidez\textquotedbl{} (02\_marco\_teorico.tex:315).}}\par
\begin{pasaje}\footnotesize p. 464: \textquotedbl{}For any higher order polynomial of order p, the Gauss-Legendre points numbering p will also provide points of superconvergent sampling. [...] In Fig. 13.6 we contrast these points with the minimum number of quadrature points necessary for obtaining an accurate (though not always stable) stiffness representation and find these to be almost coincident at all times.\textquotedbl{}  \textbar{}\textbar{}  p. 467 (Fig. 13.8, leida como imagen): \textquotedbl{}Interior superconvergent patches for quadrilateral elements (linear, quadratic, and cubic) and triangles (linear and quadratic)\textquotedbl{}, donde el recuadro rotulado '9-node elements' dibuja cuatro puntos de muestreo superconvergentes (2x2) por elemento.\end{pasaje}
\vspace{3pt}\noindent{\small\textbf{p.\ 466 (Fig. 13.7 y texto de la Sec. 13.3)} \quad \campo{\textquotedbl{}EduFEM evalua las tensiones en los puntos de la cuadratura de rigidez de cada elemento y las lleva a los nodos mediante una matriz precalculada\textquotedbl{}, es decir, extrapolacion nodal desde los puntos de Gauss (02\_marco\_teorico.tex:315; y 05\_conclusiones.tex:13, 'recuperacion de tensiones por extrapolacion desde los puntos de Gauss').}}\par
\begin{pasaje}\footnotesize Fig. 13.7 Cantilever beam with four quadratic (Q8) elements. Stress sampling at cubic order (2 x 2) Gauss points with extrapolation to nodes.  \textbar{}\textbar{}  \textquotedbl{}Here other procedures were necessary, for instance Hinton and Campbell9 suggested a method in which stresses at all nodes were calculated by extrapolating the Gauss point values.\textquotedbl{}\end{pasaje}
\vspace{3pt}\noindent{\small\textbf{p.\ 207 (Sec. 6.4) y 476-477 (Sec. 13.6...} \quad \campo{El promediado nodal \textquotedbl{}suaviza el campo y produce un contorno continuo\textquotedbl{} y \textquotedbl{}el salto entre los valores sin promediar es un indicador util del error de discretizacion -es el principio en que se apoyan los estimadores de error por recuperacion-\textquotedbl{} (02\_marco\_teorico.tex:357).}}\par
\begin{pasaje}\footnotesize p. 207: \textquotedbl{}The answers which are obtained by the finite element calculation result in discontinuities of both strains and stresses between adjacent elements. [...] In the very early days of finite element calculation with simple C0 continuity elements an averaging of element strains and stresses [...] was made at each node. This of course gave improved results at most points - except at those which were on the boundary.\textquotedbl{}  \textbar{}\textbar{}  pp. 476-477: \textquotedbl{}One of the most important applications of the recovery methods is its use in the computation of a posteriori error estimators. With the recovered solutions available, we can now evaluate errors simply by replacing the exact values of quantities such as u, sigma, etc., which are in general unknown, [...] by the recovered values which are more accurate than the direct finite element solution. [...] \textbar{}\textbar{}e\_sigma\textbar{}\textbar{}\_L2 \textasciitilde{} \textbar{}\textbar{}sigma* - sigma\_hat\textbar{}\textbar{}\_L2 (13.39) [...]...\end{pasaje}
\vspace{3pt}\noindent{\small\textbf{p.\ 38 (Sec. 2.6 'Discretization error and...} \quad \campo{\textquotedbl{}para un elemento de orden polinomico p, se espera O(h\textasciicircum{}\{p+1\}) en la norma L2 y O(h\textasciicircum{}p) en la seminorma H1\textquotedbl{} (04\_resultados.tex:16; cita conjunta con strang2008analysis y hughes2000fem; reiterado en 05\_conclusiones.tex:13).}}\par
\begin{pasaje}\footnotesize p. 38: \textquotedbl{}If within an element of 'size' h a polynomial expansion complete to degree p is employed, this can fit locally the Taylor expansion up to that degree and, as x - xa and y - ya are of the order of magnitude h, the error in u will be of the order O(h\textasciicircum{}\{p+1\}). [...] By a similar argument the strains (or stresses) which are given by the mth derivatives of displacement should converge with an error of O(h\textasciicircum{}\{p+1-m\}), i.e., as O(h) in the plane stress example, where m = 1.\textquotedbl{}  \textbar{}\textbar{}  p. 458: \textquotedbl{}We noted there that with trial functions in the displacement formulation of degree p, the errors in the stresses were of the order O(h\textasciicircum{}p). This order of error should therefore apply to the energy norm error \textbar{}\textbar{}e\textbar{}\textbar{}.\textquotedbl{}\end{pasaje}
\vspace{3pt}\noindent{\small\textbf{p.\ 145-147 (Sec. 5.5) y 163 (Eq. 5.47)} \quad \campo{\textquotedbl{}la formulacion isoparametrica, la cuadratura de Gauss y el ensamblaje son directamente generalizables a tres dimensiones\textquotedbl{} (01\_introduccion.tex:64; cita conjunta con bathe2014fem).}}\par
\begin{pasaje}\footnotesize p. 147: \textquotedbl{}For two-dimensional problems we drop all the terms containing z and/or zeta in Eqs (5.8) to (5.11).\textquotedbl{}  \textbar{}\textbar{}  p. 147: \textquotedbl{}This type of transformation is valid irrespective of the number of coordinates used. [...] One- and two-dimensional problems will similarly result in integrals with respect to one or two coordinates within simple limits.\textquotedbl{}  \textbar{}\textbar{}  p. 163: \textquotedbl{}For a brick we have similarly I = INT INT INT f(xi,eta,zeta) dxi deta dzeta = SUM\_k SUM\_j SUM\_i f(xi\_i, eta\_j, xi\_k) wi wj wk (5.47)\textquotedbl{}\end{pasaje}
\vspace{6pt}
\end{document}